\chapter{Valuations of linear growth and toric models}
\index{valuation!linear growth}\index{toric geometry!BJ model}
\label{chap:bj-linear-growth}

The finite-energy theory identifies the analytically relevant part of the
valuation space.  A valuation belongs to it exactly when its vanishing on
sections grows at most linearly with the degree.  Dirac masses then embed
this locus into the strong space of measures.  In the toric case the entire
picture becomes convex analysis: psh functions are convex functions,
Monge--Amp\`ere is Alexandrov's measure, and energy is an integral of a
Legendre transform.

Throughout, $X$ is an integral projective unibranch variety over an algebraically
closed trivially valued field, $L$ is an ample line bundle,
$\omega=c_1(L)$, and $v_{\triv}$ is the trivial valuation of $k(X)$.
A point with proper support is retained in $X^{\An}$, but it is not a
valuation of the function field and hence does not enter $X^{\lin}$.

\section{Orientation}
\label{sec:ch49-orientation}

\subsection{What this chapter does}
\label{sec:ch49-what}

\Cref{chap:bj-energy,chap:bj-finite-energy,chap:bj-strong-topologies} built
an analysis on $X^{\An}$: an energy, a finite-energy class
$\cE^1(\omega)$, a dual energy on measures, and strong topologies for
which the Monge--Amp\`ere operator gives a homeomorphism from
$\cE^1(\omega)/\mathbf R$ to $\cM^1(\omega)$.  All of it
was carried out with the points of $X^{\An}$ as an undifferentiated
compact Hausdorff space, with distinguished loci: $X^{\divis}$ supplies
the weak topology on psh functions, while every psh function is finite on
$X^{\lin}$.  Differences of finite-energy potentials are also defined
almost everywhere for each finite-energy measure by integrability;
this does not require that measure to be carried by $X^{\lin}$.
The linear-growth locus was
introduced in \Cref{def:bj-linear-growth}, its elementary properties were
proved in \Cref{chap:trivial-analytification}, and every later chapter has
used it as a black box.  \emph{This chapter opens the box.}

Two subjects organize the chapter.  The first is an exact identification: the set of
points at which the analysis of the preceding three chapters is
nondegenerate --- equivalently, the set of points whose Dirac mass has
finite dual energy, equivalently the set of nonpluripolar points --- is
precisely the linear-growth locus of \Cref{def:bj-linear-growth}, a purely
algebro-geometric condition on the vanishing orders of a valuation on
sections of $mL$.  Around that identification there is a
metric geometry: a complete metric $d_{\infty,\omega}$ on $X^{\lin}$,
under which $v\mapsto\delta_v$ is a bi-Lipschitz embedding into the
strong space $\cM^1(\omega)$ of \Cref{chap:bj-strong-topologies}, and in
which divisorial points are dense in characteristic zero.

The second is a complete dictionary in the toric case.  If $X$ is toric
and $\phi$ is invariant under the torus, then $\phi$ \emph{is} a convex
function on $N_{\mathbf R}$, its Monge--Amp\`ere measure \emph{is} the
Alexandrov measure of that convex function, and its energy \emph{is} the
integral of the Legendre transform over the polytope.  The three formulae
are stated in \Cref{thm:bj-toric-psh-convex,thm:bj-toric-ma-energy}, and
they are the formulae that earlier chapters have been using to check
themselves: \Cref{ex:ch46-toric-p2} computes an energy on $\mathbf P^2$
three ways and quotes \eqref{eq:bj-toric-energy} for one of them;
\Cref{ex:ch47-toric-cy} solves two Calabi--Yau equations by reading them
off a polytope; \Cref{ex:ch47-dirac-green} points at
\eqref{eq:bj-toric-green} for the general closed formula for Green functions at toric valuations.  Because those computations are cited elsewhere as
consistency checks, the formulae below display every sign and normalization.
The section, bounded measure, and energy calculations are supplied here;
general invariant extension and the full strong/$L^1$ correspondence
retain their named Appendix~B inputs.

Three features distinguish this chapter from its predecessors.

First, \emph{the geometry is on the base, not on the potentials}.  The
objects metrized here are valuations.  All the analytic work has been done;
what remains is to transport it along the map $v\mapsto\delta_v$, and the
transport is exact --- $d_{\infty,\omega}$ and the measure gauge
$I^\vee$ are bi-Lipschitz equivalent on the image, by
\Cref{thm:bj-linear-topology-complete}\,\cref{itm:ch49-top:embedding}.  A
reader who wants a slogan may take this one: \emph{the linear-growth locus
is the set of nonpluripolar points, and its intrinsic metric is uniformly
comparable to the finite-energy measure gauge on Dirac masses}.

Second, \emph{the published Radon--Riesz question has a negative answer}.
\Cref{rem:ch49-open-question} records the exact implication asked in
\cite[Question~11.16]{BJ22}, and
\cref{thm:ch49-global-radon-riesz-counterexample} disproves it on
$\mathbf P^n$ for every $n\ge2$.  The equality of exhaustion values at
the weak limit is essential: uniform comparability in
\eqref{eq:bj-dirac-energy-T} alone would not supply it.

Third, \emph{the toric half is an equality of two theories, not an
analogy}.  \Cref{chap:model-calculations} developed toric Monge--Amp\`ere
over a nontrivially valued field, with tropical coordinate
$u=-\log|\chi|$, reflected support function $\check h_P$, and the
factor $n!$ of \eqref{eq:appJ-toric-degree}.  The trivially valued theory
of this chapter has a \emph{different} polytope convention, inherited from
\cite[Appendix~B]{BJ22}, and the difference is a reflection.
\Cref{warn:ch49-toric-signs} converts one into the other once and for all,
and every toric formula below is checked against
\Cref{chap:model-calculations} and against
\Cref{thm:ch26-toric-energy} in \Cref{sec:ch49-unit-tests}.

\subsection{Standing hypotheses and notation}
\label{sec:ch49-hypotheses}

The hypotheses are those of
\Cref{chap:bj-energy,chap:bj-finite-energy,chap:bj-strong-topologies}.
Thus $k$ is algebraically closed and trivially valued, $X$ is an
integral projective $k$-variety of dimension $n\ge1$ which is
\emph{unibranch}, $L$ is an ample line bundle with
$\omega:=c_1(L)\in\Amp(X)$, and
\[
 V:=(\omega^n)>0 .
\]
We write $\MA_\omega(\phi)=V^{-1}(\omega+\ddc\phi)^n$ for the
\emph{probability} Monge--Amp\`ere measure and $\MA^{\mathrm u}_\omega$
for the unnormalized one, exactly as in \Cref{warn:ch46-normalization}.
The classes $\cE^1(\omega)$, $\cE^1_0(\omega)$,
$\PSH^{\sup}(\omega)$, $\cM^1(\omega)$ are those of
\eqref{eq:ch47-normalized-classes} and
\Cref{def:bj-e1,def:bj-finite-energy-measure}; the functionals
$E_\omega$, $E^\vee_\omega$, $I$, $J_\psi$, $J$ are those of
\eqref{eq:bj-energy-definition}, \eqref{eq:bj-dual-energy},
\eqref{eq:bj-I-J-potentials} and \eqref{eq:ch47-J-functional}.

Three loci of $X^{\An}$ occur constantly, and were fixed in
\eqref{eq:ch34-loci}:
\[
 X^{\divis}\ \subset\ X^{\lin}\ \subset\ X^{\val}\ \subset\ X^{\An},
\]
the divisorial points, the linear-growth locus of
\Cref{def:bj-linear-growth}, the full-support points
\eqref{eq:triv-full-support-locus}, and the whole analytification.  The first inclusion is strict because the divisorial convention uses
rational scaling; the middle inclusion is strict in examples of dimension
at least two (\Cref{ex:ch49-no-linear-growth}).  Proper-support points
give the last strict inclusion.  For $n=1$ one has
$X^{\lin}=X^{\val}=X^{\An}\setminus\{v_{p,\triv}:p\in X(k)\}$, a statement proved
in \Cref{ex:ch49-curve-case}.

The partial order on semivaluations is \eqref{eq:semivaluation-order}:
$v\le w$ means $v(\mathfrak a)\le w(\mathfrak a)$ for every
coherent ideal $\mathfrak a$.  Equivalently, the centers and local regular
functions satisfy \Cref{thm:ch32-order-centers}.  The
scaling action is $v\mapsto tv$, and its action on functions is
\eqref{eq:function-scaling-action}.  Every $\theta$-psh function is
decreasing for $\le$ and satisfies the maximum principle
$\sup_{X^{\An}}\phi=\phi(v_{\triv})$
(\Cref{prop:ch34-pl-operations}\,\cref{itm:ch34-plop:decreasing}, extended
to $\PSH$ by decreasing regularization
\Cref{thm:ch34-fs-regularization}); this is used so often below that we
shall not always cite it.

For the toric part, $T\simeq\mathbf G_m^n$ is a split torus with
character lattice $M$ and cocharacter lattice $N=M^\vee$, the pairing
is $\langle\,\cdot\,,\,\cdot\,\rangle:N_{\mathbf R}\times M_{\mathbf R}\to
\mathbf R$, and $\rd m$ is the lattice Lebesgue measure on
$M_{\mathbf R}$ normalized so that a fundamental parallelepiped of $M$
has volume one, as in \eqref{eq:appJ-alexandrov}.  Convex-analytic notation
is that of \Cref{app:convexity}: $f^\vee$ is the Legendre--Fenchel
conjugate \eqref{eq:appD-conjugate}, $\partial f$ the subdifferential
\eqref{eq:appD-subgradient}, and $\MA_{\mathbf R}$ the Alexandrov measure
\eqref{eq:appD-alexandrov-measure}.

\subsection{The normalization, inherited}
\label{sec:ch49-normalization}

\begin{warning}{Probability normalization, and the two toric conventions}
\label{warn:ch49-normalization}
\Cref{warn:ch46-normalization} declared a departure from the front-matter
convention for the whole of
\Crefrange{chap:bj-energy}{chap:bj-stability}, and it is in force here
without amendment: $\MA_\omega$ is a probability measure,
$E_\omega(\phi+c)=E_\omega(\phi)+c$, and $E_\omega$ is normalized by
$E_\omega(0)=0$.  Two consequences are used constantly below.  The dual
energy $E^\vee_\omega$ is finite only on probability measures
(\Cref{lem:ch47-mass-forced}), so the Dirac masses considered here are
automatically of the right mass.  And the toric energy formula
\eqref{eq:bj-toric-energy} carries the factor $V_P^{-1}$, which the
printed source suppresses by integrating against normalized Lebesgue
probability measure on the polytope; the source note
\Cref{sec:ch49-source-note} records this.

Separately, and independently of the probability normalization, there are
\emph{two} toric polytope conventions in this book.
\Cref{chap:model-calculations} and \Cref{chap:bfj-examples} use $P$ and
the reflected support function $\check h_P$ of
\eqref{eq:appJ-toric-recession}; this chapter uses the Boucksom--Jonsson
polytope $P_{\mathrm{BJ}}=-P$ and its support function $f_P$.  The
conversion is \eqref{eq:bj-toric-sign-conversion} and is proved in
\Cref{warn:ch49-toric-signs}.  Mixing them produces an energy of the wrong
sign, which the trivial metric detects immediately.
\end{warning}

\subsection{Extended-real ledger}
\label{sec:ch49-ledger}

The extended-real conventions of
\Cref{sec:ch46-ledger,sec:ch47-ledger,sec:ch48-ledger} remain in force.
The following domains are especially relevant here.

\begin{enumerate}
\item\label{itm:ch49-ledger:T}
      The functions $T_L$ of \eqref{eq:bj-maximal-vanishing} and
      $T_\omega$ of \eqref{eq:bj-T-valuation} take values in
      $[0,+\infty]$.  The value $+\infty$ is genuinely attained, on
      every proper-support point (\Cref{prop:bj-linear-growth-potential})
      and on some full-support points as soon as $n\ge2$
      (\Cref{ex:ch49-no-linear-growth}).  No difference of two such values
      is ever formed: $d_{\infty,\omega}$ is defined only on
      $X^{\lin}$, where both are finite.
\item\label{itm:ch49-ledger:legendre}
      In \eqref{eq:bj-toric-legendre} the conjugate $f^\vee$ takes the
      value $+\infty$, and it does so on a set of positive measure
      when the closure of its effective domain is a proper subset of
      $P_{\mathrm{BJ}}$, equivalently when the recession function differs
      from $f_P$.  Even if the conjugate is finite almost everywhere,
      its integral can diverge.  The value $-\infty$ never occurs, because $f$ is finite on
      $N_{\mathbf R}\subset X^{\lin}$.  The integral in
      \eqref{eq:bj-toric-energy} is therefore an integral of a function
      with values in $(-\infty,+\infty]$, well defined in
      $(-\infty,+\infty]$, and $E_\omega(\phi)=-\infty$ exactly when it
      diverges.  This is the content of \eqref{eq:bj-toric-e1}.
\item\label{itm:ch49-ledger:sup}
      The supremum $\sup_{X^{\An}}\phi$ is a real number for every
      $\phi\in\PSH(\omega)$
      (\Cref{lem:ch47-three-normalizations}), so every drop
      $\sup\phi-\phi(v)$ occurring in \eqref{eq:bj-T-valuation} and
      \eqref{eq:bj-dinfty} is a well-defined element of
      $[0,+\infty]$, never of the form $(+\infty)-(+\infty)$.
\end{enumerate}

\subsection{What is imported, and what is proved}
\label{sec:ch49-imported}

\begin{important}{Published interfaces and local proofs}
\label{imp:ch49-imported}
The chapter uses the earlier analytic theory and the following precise
interfaces of \cite{BJ22}.  The proofs below identify which parts are
derived here and which are quoted.
\begin{enumerate}
\item\label{itm:ch49-import:sections}
      \emph{Section formula}, \cite[Lemma~11.7]{BJ22}:
      \eqref{eq:bj-dinfty-sections}.  We prove it from finite maxima
      and the fixed-class regularization
      \Cref{thm:ch34-fs-regularization}; no varying-class limit is needed.
\item\label{itm:ch49-import:embedding}
      \emph{Dirac-mass comparison}, \cite[Theorem~11.13]{BJ22}:
      there are dimensional constants $a_n,b_n>0$ such that
      \[
      a_n d_{\infty,\omega}(v,w)
      \le I^\vee(\delta_v,\delta_w)
      \le b_n d_{\infty,\omega}(v,w)
      \qquad(v,w\in X^{\lin}).
      \]
      The upper bound with constant $2$ is proved in
      \Cref{prop:ch49-embedding-easy}; the general lower bound is quoted.
      The target is a quasi-metric space, so this does not identify
      $I^\vee$ with a metric.
\item\label{itm:ch49-import:approximation}
      \emph{Divisorial approximation from below},
      \cite[Corollary~A.3]{BJ22}: in characteristic zero, any
      $v\in X^{\An}$ is the weak limit of divisorial $v_i\le v$.
      The net need not be nondecreasing.  The corollary includes singular
      $X$ by passing to a resolution.  The SNC-test-configuration route
      on $X$ itself assumes $X$ smooth.
\item\label{itm:ch49-import:toriccore}
      \emph{Toric dictionary}, \cite[Appendix~B, pp.~830--831]{BJ22}:
      invariant psh functions correspond to finite convex
      $f\le f_P+O(1)$, with the retraction identity
      \eqref{eq:bj-toric-retraction-identity}.  The elementary
      finite-level calculation is proved in
      \Cref{lem:ch49-toric-finite-level}; the general existence and
      invariant-extension statement are quoted.  Bounded MA,
      extended energy, and the full strong/L1 correspondence are
      separate interfaces in the same appendix.
\item\label{itm:ch49-import:extreme}
      \emph{Integrability characterization},
      \cite[Theorem~9.20]{BJ22}: a Radon probability measure has finite
      dual energy if and only if all $\cE^1(\omega)$ functions are
      integrable.  It gives stability under domination and the converse
      in the extreme-point description
      \cite[Proposition~11.3]{BJ22}.
\item\label{itm:ch49-import:equicontinuity}
      \emph{Bounded-set continuity and sublinear growth},
      \cite[Theorems~11.11--11.12]{BJ22}:
      the assertions of \Cref{prop:bj-psh-continuous-bounded-valuations}.
      The first concerns each fixed psh function, not equicontinuity
      of the whole psh family on a $T$-ball.  Its proof uses decreasing
      approximation and the uniform pairing estimate of Theorem~10.3.
\item\label{itm:ch49-import:green}
      \emph{Continuous Green envelopes and isometry},
      \cite[Theorem~12.12 and Proposition~12.14]{BJ22}.
      \Cref{cor:ch49-green-invariants} proves the envelope, exact height,
      and sup-norm identities from the existing domination principle,
      with continuity supplied by Theorem~12.12.
\item\label{itm:ch49-import:snc}
      \emph{SNC equicontinuity},
      \cite[Theorem~A.4]{BJ22}: for a smooth projective $X$ in
      characteristic zero, each SNC dual complex is strongly compact,
      and the restrictions of $\PSH(\omega)$ are equi-Lipschitz on
      its faces, including their boundaries.
\end{enumerate}
Convex and toric intersection inputs are recorded in
\Cref{sec:ch49-source-note}.  Fixed-measure continuity on bounded
J-sublevels is \cite[Proposition~9.19]{BJ22}, already used in
\Cref{chap:bj-finite-energy,chap:bj-strong-topologies}.
\end{important}

\begin{important}{Standing envelope input}
\label{imp:ch49-envelope}
As in \Cref{imp:ch46-envelope-imported,imp:ch47-envelope,imp:ch48-envelope},
the variety $X$ is unibranch and the envelope property for $\omega$ is
an explicit input. The general route through \Cref{thm:bj-final-envelope}
retains $\mathrm{(SD)}_k$; independent instances retain their scope.
With this input the firewall of \Cref{imp:ch45-firewall-discharged} is
discharged and $P_\omega$ is an operator
(\Cref{imp:ch45-obstacle-operator}).  Three consequences are used here.
\begin{enumerate}
\item\label{itm:ch49-env:cy}
      The Calabi--Yau theorem \Cref{cor:bj-cy-unibranch} solves
      $\MA_\omega(\phi)=\delta_v$ for every $v\in X^{\lin}$; this is
      what makes the Green function $\phi_v$ of
      \Cref{ex:ch47-dirac-green} available, and it is used in the sharp
      upper bound of \Cref{thm:bj-dirac-finite-energy}.
\item\label{itm:ch49-env:homeo}
      $\MA_\omega$ is a homeomorphism
      $\cE^1(\omega)/\mathbf R\simeq\cM^1(\omega)$ for the strong
      topologies (\Cref{thm:ch48-ma-homeomorphism}); this identifies the analytic solution spaces.  The direct completeness
      proof below uses bounded-set continuity and compact $T$-sublevels.
\item\label{itm:ch49-env:polar}
      The equivalence of pluripolarity, negligibility and vanishing outer
      capacity (\Cref{thm:bj-small-set-equivalence}) is used only to phrase
      \Cref{cor:ch49-nonpluripolar}, and nothing later depends on that
      corollary.
\end{enumerate}
\end{important}

\subsection{Relation to \texorpdfstring{\Cref{part:bfj}}{Part III} and to
\texorpdfstring{\Cref{chap:model-calculations}}{Chapter 8}}
\label{sec:ch49-part3}

The trivially valued base supplies a scaling action
$v\mapsto tv$ on $X^{\An}$ and a distinguished vertex $v_{\triv}$.
The identity $T_\omega(tv)=tT_\omega(v)$ makes $T_\omega$ the distance
from that vertex for $d_{\infty,\omega}$.  It is finite exactly on
$X^{\lin}$ and infinite at proper-support points; some full-support
valuations also have infinite growth in dimension at least two.
This cone structure is specific to the trivial absolute value.
For a fixed nontrivial absolute value, scaling a valuation changes its
restriction to the ground field and therefore leaves that fixed-base
analytification.  This difference does not assert that all nontrivially
valued analytic points have finite-energy Dirac masses.

The toric numerical formulae, by contrast, correspond after the stated
normalization and reflection.
\Cref{chap:model-calculations} proved
$\trop_*c_1(L,\|\cdot\|)^n=n!\MA_{\mathbf R}(F)$ for a semipositive toric
model metric over a nontrivially valued field
(\Cref{thm:toric-ma-comparison}), and
\Cref{chap:bfj-examples} extended it to the BFJ setting
(\Cref{thm:toric-bfj-alexandrov}); \Cref{thm:ch26-toric-energy} proved the
corresponding energy formula in terms of the roof function
\Cref{def:ch26-roof}.  The formulae of \Cref{sec:ch49-toric-ma} are the
corresponding formulae in the trivially valued theory, separately defined
and proved and rewritten in the reflected convention.  This numerical
correspondence does not identify the ambient model-intersection pairings
over the two bases.  \Cref{sec:ch49-unit-tests} carries out the
conversion on the nose, on the two examples that
\Cref{part:bfj} computed to numbers.

The following table records the correspondence and is the toric half of the
master dictionary.

\begin{center}
\begin{tabular}{p{.24\textwidth}p{.68\textwidth}}
\toprule
quantity & conventions\\
\midrule
base  & \emph{Nontrivial base:} $K$ nontrivially valued\par\smallskip \emph{Trivial base:} $k$ trivially valued\\

polytope  & \emph{Nontrivial base:} $P\subset M_{\mathbf R}$\par\smallskip \emph{Trivial base:} $P_{\mathrm{BJ}}=-P$\\

support function  & \emph{Nontrivial base:} $\check h_P(u)=h_P(-u)$\par\smallskip \emph{Trivial base:} $f_P=f_{P_{\mathrm{BJ}}}$,
 and $f_P=\check h_P$\\

potential  & \emph{Nontrivial base:} convex $F$ with $F-\check h_P$ bounded\par\smallskip \emph{Trivial base:} convex $f_\phi=(\phi+f_P)|_{N_{\mathbf R}}$, $f_\phi\le f_P+O(1)$\\

measure  & \emph{Nontrivial base:} $\trop_*c_1(L,\|\cdot\|)^n=n!\MA_{\mathbf R}(F)$\par\smallskip \emph{Trivial base:} $\MA_\omega(\phi)=V_P^{-1}\iota_*\MA_{\mathbf R}(f_\phi)$\\

energy  & \emph{Nontrivial base:} $E(F)-E(F_0)=n!\int_P(g_F-g_{F_0})\rd\lambda$\par\smallskip \emph{Trivial base:} $E_\omega(\phi)=-V_P^{-1}\int_{P_{\mathrm{BJ}}}f_\phi^\vee\,\rd m$\\

roof vs.\ conjugate  & \emph{Nontrivial base:} $g_F(y)=-F^\vee(-y)$\par\smallskip \emph{Trivial base:} $f^\vee_\phi$ on
 $P_{\mathrm{BJ}}$\\

mass test  & \emph{Nontrivial base:} $n!\vol_M(P)=(L^n)$\par\smallskip \emph{Trivial base:} $V=n!\,V_P$\\
\bottomrule
\end{tabular}
\end{center}

\subsection{Roadmap}
\label{sec:ch49-roadmap}

\Cref{sec:ch49-growth} constructs $T_L$ and $T_\omega$, proves they
agree, and establishes the elementary calculus of $T_\omega$:
homogeneity, monotonicity, lower semicontinuity, and independence of the
polarization.  \Cref{sec:ch49-dirac} proves the finite-energy dichotomy
\Cref{thm:bj-dirac-finite-energy} with both constants, and identifies the
extreme points of $\cM^1(\omega)$.  \Cref{sec:ch49-topologies}
constructs $d_{\infty,\omega}$, proves the comparison lemma that drives
every later approximation argument, and assembles the topological theorem.
\Cref{sec:ch49-density} proves divisorial density and the dual-complex
statements.  \Cref{sec:ch49-toric} proves the toric dictionary and
\Cref{sec:ch49-toric-ma} the explicit formulae; \Cref{sec:ch49-unit-tests}
checks them.  \Cref{sec:ch49-examples} carries the examples to explicit numbers.  \Cref{sec:ch49-scope} audits the hypotheses,
\Cref{sec:ch49-dictionary} is the chapter dictionary, and
\Cref{sec:ch49-source} the provenance.

\section{Maximal vanishing and linear growth}
\label{sec:ch49-growth}

\subsection{The vanishing sequence and its Fekete limit}
\label{sec:ch49-vanishing}

Choose $d\geq1$ such that $mL$ has nonzero sections for every
$m\in d\mathbf Z_{>0}$.  For $v\in X^{\val}$ and such an $m$, put

\[
 a_m(v):=\max\{v(s)\mid 0\ne s\in H^0(X,mL)\},
\tag{49.1}\label{eq:bj-max-vanishing}
\]

where $v(s)$ is computed in a local frame at the center of $v$.
Changing that frame multiplies $s$ by a unit and does not change the
value.  If $mL$ has a base point at the center, sections vanishing there
may have positive value; the maximum is nevertheless finite for a
function-field valuation.

That last sentence deserves an argument, because it is the only place where
full support is used.

\begin{lemma}
\label{lem:ch49-am-finite}
Let $v\in X^{\val}$, let $W$ be a nonzero finite-dimensional subspace of
$H^0(X,mL)$, and set $F^tW:=\{s\in W\mid v(s)\ge t\}$ for
$t\in\mathbf R$.  Then each $F^tW$ is a linear subspace, the family is
decreasing with at most $\dim W$ jumps, $F^tW=0$ for $t$ large, and
consequently
\[
 \max\{v(s)\mid 0\ne s\in W\}
\]
is a finite real number, attained.  If instead $v\in X^{\An}$ has proper
support, then $v(s)=+\infty$ for every section $s$ vanishing on the
support $s(v)$, and such a section exists in $H^0(X,mL)$ for all large
$m$.
\end{lemma}

\begin{proof}
Fix a local frame $\tau$ of $mL$ at $c(v)$ and write $s=f_s\tau$
with $f_s\in\cO_{X,c(v)}$, so that $v(s)=v(f_s)$
(\Cref{lem:ch34-frame-independence}).  The map $s\mapsto f_s$ is
$k$-linear, and $v$ satisfies $v(f+g)\ge\min\{v(f),v(g)\}$ and
$v(\lambda f)=v(f)$ for $\lambda\in k^\times$, the valuation being
trivial on $k$.  Hence $F^tW$ is a subspace, and it decreases in $t$.
Since $v\in X^{\val}$, the support of $v$ is the generic point of
$X$, so $v(s)<+\infty$ for every $s\ne0$; therefore
$\bigcap_{t}F^tW=0$.  A decreasing family of subspaces of a
$\dim W$-dimensional space that meets in $0$ has at most $\dim W+1$ distinct members and is eventually $0$.
Let $U$ be its smallest nonzero member.  All nonzero elements of $U$
have the same valuation: unequal values would produce a smaller nonzero
filtration member.  Every section outside $U$ has smaller value, again
by the filtration.  Thus any nonzero element of $U$ attains the finite
maximum.

For the last assertion, let $Z:=s(v)\subsetneq X$ be the support, a
proper integral closed subscheme (\Cref{def:ch32-proper-support}).  Since
$L$ is ample, $\cI_Z\otimes mL$ is globally generated for $m\gg0$ and
has nonzero global sections, $Z\ne X$.  Any such section $s$ satisfies
$f_s\in\mathfrak p_v$, so $v(s)=+\infty$.
\end{proof}

Multiplication of sections gives

\[
                         a_{m+m'}(v)\geq a_m(v)+a_{m'}(v).
\tag{49.2}\label{eq:bj-vanishing-superadditive}
\]

Indeed, if $s$ and $s'$ attain the two maxima, then $ss'$ is a
nonzero section of $(m+m')L$ and, in a product frame at $c(v)$,
$v(ss')=v(s)+v(s')$ by multiplicativity of $v$.

Fekete's lemma on this Veronese semigroup therefore gives

\[
 T_L(v):=\lim_{\substack{m\to\infty\\d\mid m}}\frac{a_m(v)}m
        =\sup_{\substack{m\geq1\\d\mid m}}\frac{a_m(v)}m
        \in[0,+\infty].
\tag{49.3}\label{eq:bj-maximal-vanishing}
\]

The value is independent of the admissible Veronese index $d$: passing
to a common multiple leaves the Fekete limit unchanged.  Equivalently, one
may replace $L$ by a sufficiently large positive power and then take all
positive degrees.

\begin{lemma}
\label{lem:ch49-vanishing-elementary}
Let $v\in X^{\val}$.
\begin{enumerate}
\item\label{itm:ch49-vs:allm}
      $a_m(v)/m\le T_L(v)$ for \emph{every} $m\ge1$ with
      $H^0(X,mL)\ne0$, not merely for $m\in d\mathbf Z$.
\item\label{itm:ch49-vs:d}
      $T_L(v)$ does not depend on the choice of admissible $d$.
\item\label{itm:ch49-vs:homog}
      $a_m(tv)=t\,a_m(v)$ and $T_L(tv)=t\,T_L(v)$ for $t>0$.
\item\label{itm:ch49-vs:power}
      $T_{cL}(v)=c\,T_L(v)$ for every $c\in\mathbf Q_{>0}$ such that
      $cL$ is an honest line bundle.
\item\label{itm:ch49-vs:order}
      $v\le w$ implies $T_L(v)\le T_L(w)$.
\end{enumerate}
\end{lemma}

\begin{proof}
\cref{itm:ch49-vs:allm}  Choose $d'\ge1$ with $d\mid md'$.  Iterating
\eqref{eq:bj-vanishing-superadditive} gives
$a_{md'}(v)\ge d'a_m(v)$, whence
$a_m(v)/m\le a_{md'}(v)/(md')\le T_L(v)$.

\cref{itm:ch49-vs:d}  If $d\mid d'$, then $d'\mathbf Z_{>0}\subset
d\mathbf Z_{>0}$ is cofinal, so the limit in
\eqref{eq:bj-maximal-vanishing} along $d'\mathbf Z$ exists and equals the
one along $d\mathbf Z$.  Two admissible indices have a common multiple,
which is admissible.

\cref{itm:ch49-vs:homog}  Immediate from $(tv)(s)=t\,v(s)$.

\cref{itm:ch49-vs:power}  Choose sufficiently divisible degrees $m$
so that $mc$ is an integer, the rational identification
$m(cL)=(mc)L$ is an identification of line bundles, and both degree
sequences are admissible.  Then $a_m(v;cL)=a_{mc}(v;L)$.
Dividing by $m$ and using the independence of admissible Veronese
indices gives $T_{cL}(v)=cT_L(v)$.

\cref{itm:ch49-vs:order}  For each $s$, $v(s)\le w(s)$; take maxima and
limits.
\end{proof}

\begin{definition}[Linear growth]
\label{def:bj-linear-growth}
A valuation $v\in X^{\val}$ has \emph{linear growth} if
$T_L(v)<+\infty$.  The resulting locus is denoted
$X^{\lin}\subset X^{\val}\subset X^{\An}$.
\end{definition}

The definition is independent of $L$: if $L'$ is ample, powers of $L$
and $L'$ dominate one another up to globally generated twists, giving
two-sided linear estimates for $a_m$.  Every divisorial valuation has
linear growth \cite[Proposition~1.28]{BJ22}.

We prove both assertions.

\begin{proposition}
\label{prop:ch49-change-polarization}
Let $L$ and $L'$ be ample line bundles on $X$ and choose
$p,q\in\mathbf Z_{>0}$ with $pL-L'$ and $qL'-L$ ample.  Then
\[
 q^{-1}T_L(v)\ \le\ T_{L'}(v)\ \le\ p\,T_L(v)
 \qquad\text{for every }v\in X^{\val}.
\]
In particular $T_L$ and $T_{L'}$ are finite on exactly the same set, so
$X^{\lin}$ does not depend on the polarization, and the two functions are
uniformly comparable.
\end{proposition}

\begin{proof}
Put $A:=pL-L'$, an ample line bundle, so that $L'+A=pL$.  Choose
$d_A\ge1$ with $H^0(X,mA)\ne0$ for every $m\in d_A\mathbf Z_{>0}$,
and enlarge $d$ so that $d_A\mid d$.  For $m\in d\mathbf Z_{>0}$
pick $0\ne\tau\in H^0(X,mA)$ and let $s'\in H^0(X,mL')$ attain
$a_m(v;L')$.  Then $s'\tau$ is a nonzero section of
$m(L'+A)=mpL$, and
\[
 a_{mp}(v;L)\ \ge\ v(s'\tau)=v(s')+v(\tau)\ \ge\ a_m(v;L') .
\]
Dividing by $mp$ and letting $m\to\infty$ along $d\mathbf Z$, and
using \Cref{lem:ch49-vanishing-elementary}\,\cref{itm:ch49-vs:allm}, gives
$T_L(v)\ge p^{-1}T_{L'}(v)$, which is the right-hand inequality.  The
left-hand one is the same argument with the roles of $L$ and $L'$
exchanged and $A$ replaced by $qL'-L$.
\end{proof}

\subsection{Divisorial points and the pseudoeffective threshold}
\label{sec:ch49-psef}

For a divisorial point the invariant $T_L$ has a purely birational
description, which is both the shortest proof that divisorial points have
linear growth and the fastest way to compute $T_L$ in examples.

\begin{proposition}
\label{prop:ch49-psef-threshold}
\index{pseudoeffective threshold}
Let $\mu:Y\to X$ be a normal projective birational model, $E\subset Y$
a prime divisor, and $t\in\mathbf R_{>0}$.  Every such valuation has linear growth.  If $E$ is additionally
$\mathbf Q$-Cartier, then
\[
 T_L(t\ord_E)=t\cdot\tau_L(E),
 \qquad
 \tau_L(E):=\sup\{\,c\ge0\ :\ \mu^*L-cE\ \text{is pseudoeffective}\,\},
\tag{49.4}\label{eq:ch49-psef-threshold}
\]
and $\tau_L(E)<\infty$.  In particular $X^{\divis}\subset X^{\lin}$.
\end{proposition}

\begin{proof}
First let $E$ be any prime Weil divisor on the normal model $Y$.
For $0\ne s\in H^0(X,mL)$, put $a=\ord_E(s)$.  The cycle
$\operatorname{div}(\mu^*s)-aE$ is effective.  Intersecting it with
$H^{n-1}$ for an ample line bundle $H$ on $Y$ gives
\[
 a(E\cdot H^{n-1})\le m(\mu^*L\cdot H^{n-1}),
 \qquad E\cdot H^{n-1}>0.
\]
This bounds $a_m/m$ uniformly and proves divisorial linear growth
without a Cartier hypothesis, as in \cite[Proposition~1.28]{BJ22}.

For the threshold formula, assume $E$ is $\mathbf Q$-Cartier and
reduce to $t=1$ by homogeneity.  The numerical pseudoeffective cone
is closed and convex.  Its intersection with
$\{\mu^*L-cE:c\ge0\}$ is an interval $[0,\tau]$, bounded by the
same ample intersection estimate.  Every section with order $a$
makes $\mu^*L-(a/m)E$ pseudoeffective, so $T_L(\ord_E)\le\tau$.

Conversely, for rational $0\le c<\tau$, choose rational
$c<c'<\tau$ and a small rational $\varepsilon>0$ such that
$c/(1-\varepsilon)\le c'$.  Then
\[
 \mu^*L-cE=\varepsilon\mu^*L+
 (1-\varepsilon)\left(\mu^*L-\frac{c}{1-\varepsilon}E\right)
\]
is big.  Sufficiently divisible multiples therefore have nonzero
sections.  We justify their descent when $X$ is not normal.
Factor $\mu$ through the finite normalization $\nu:X^\nu\to X$.
Normal-target birational pushforward gives
$\mu_*\cO_Y=\nu_*\cO_{X^\nu}$.  Choose a nonzero section
$\tau_0\in H^0(X,bL\otimes\cI)$ of a conductor ideal
$\cI$ satisfying $\cI\,\nu_*\cO_{X^\nu}\subset\cO_X$.
Such a section exists for $b\gg0$ by ampleness.  A nonzero section
of $m\mu^*L-mcE$, viewed as a section of $m\mu^*L$, multiplied by
$\mu^*\tau_0$ descends to a nonzero section of $(m+b)L$ on $X$.
Its order along $E$ is at least $mc$, since $\ord_E(\tau_0)\ge0$.
Thus
\[
 T_L(\ord_E)\ge\lim_{m\to\infty}\frac{mc}{m+b}=c.
\]
Letting rational $c\uparrow\tau$ proves equality.  The conductor
and normal-target pushforward inputs are those of
\Cref{lem:ch45-normalization-data,app:birational-geometry}.
\end{proof}

\begin{example}[Three divisorial points on $\mathbf P^2$]
\label{ex:ch49-p2-divisorial}
Let $X=\mathbf P^2$, $L=\cO(1)$.
\begin{enumerate}
\item\label{itm:ch49-p2d:line}
      For a line $\ell$, $\tau_L(\ell)=1$: the class
      $\cO(1)-c\,\ell=\cO(1-c)$ is pseudoeffective exactly for $c\le1$.
      So $T_L(\ord_\ell)=1$.
\item\label{itm:ch49-p2d:point}
      For the exceptional divisor $E_1$ of the blowup at a point $p$,
      $\ord_{E_1}=\ord_p$ is the multiplicity valuation and
      $H-cE_1$ is pseudoeffective exactly for $c\le1$, since
      $H-E_1$ is nef with $(H-E_1)^2=0$ and
      $(H-cE_1)\cdot(H-E_1)=1-c$.  So $T_L(\ord_p)=1$.
\item\label{itm:ch49-p2d:conic}
      For a smooth conic $C$, $\cO(1)-cC$ has degree $1-2c$, so
      $\tau_L(C)=\tfrac12$ and $T_L(\ord_C)=\tfrac12$.  Directly:
      $\ord_C(s)=j$ for $s$ of degree $m$ forces $q^j\mid s$ with
      $q$ the conic equation, hence $2j\le m$.
\end{enumerate}
All three are recomputed in \Cref{sec:ch49-examples} from the toric formula
\eqref{eq:ch49-toric-T} or from an intersection number, and the numbers
agree.
\end{example}

\subsection{Linear growth as a potential-theoretic condition}
\label{sec:ch49-T-omega}

\begin{proposition}\label{prop:bj-linear-growth-potential}
For $v\in X^{\An}$, set

\[
 T_\omega(v):=
 \sup_{\phi\in\PSH(\omega)}
       \bigl(\sup_{X^{\An}}\phi-\phi(v)\bigr)\in[0,+\infty].
\tag{49.5}\label{eq:bj-T-valuation}
\]

Then $T_\omega(v)<\infty$ if and only if $v\in X^{\lin}$.  In
particular, a proper-support point is pluripolar and has
$T_\omega(v)=+\infty$.
\end{proposition}

\begin{proof}
The exact identity \Cref{prop:ch49-T-equals-T}, proved next directly
from fixed-class regularization, gives the finiteness criterion.
A regular section vanishing on the support supplies a psh pole at a
proper-support point.  The published identity is
\cite[Lemma~4.46]{BJ22}; no normalization constant occurs.
\end{proof}

\begin{proposition}
\label{prop:ch49-T-equals-T}
For every $v\in X^{\val}$,
\[
 T_\omega(v)=T_L(v),
\tag{49.6}\label{eq:ch49-T-equals-T}
\]
both sides being elements of $[0,+\infty]$; and $T_\omega(v)=+\infty$
for every $v\in X^{\An}\setminus X^{\val}$.
\end{proposition}

\begin{proof}
\emph{Step 1: $T_\omega\ge T_L$.}  Let $m\ge1$ with $H^0(X,mL)\ne0$
and let $0\ne s\in H^0(X,mL)$.  Since $X$ is integral, $s$ is a
regular section and $\phi_s:=m^{-1}\log|s|$ belongs to
$\cH^{\mathrm{gf}}_{\mathbf R}(L)\subset\PSH(\omega)$
(\Cref{def:gf-fs-functions}, \Cref{lem:ch34-regular-section}).  By
\eqref{eq:section-log-value}, $\phi_s\le0$ with
$\phi_s(v_{\triv})=-m^{-1}v_{\triv}(s)=0$, so
$\sup_{X^{\An}}\phi_s=0$ and
\[
 \sup\phi_s-\phi_s(v)=m^{-1}v(s).
\]
Taking the maximum over $s$ and then the supremum over $m$ gives
$T_\omega(v)\ge\sup_m a_m(v)/m=T_L(v)$, by
\eqref{eq:bj-maximal-vanishing}.

\emph{Step 2: $T_\omega\le T_L$ for Fubini--Study functions.}  Let
$\phi=\frac1m\max_{1\le i\le N}\{\log|s_i|+\lambda_i\}$ lie in
$\cH^{\mathrm{gf}}_{\mathbf R}(L)$, with $s_i\in H^0(X,mL)$ not all
zero and $\lambda_i\in\mathbf R$.  Discard the indices with $s_i=0$.
By \Cref{cor:ch34-generic-constancy},
$\sup_{X^{\An}}\phi=m^{-1}\max_i\lambda_i$; fix $i_0$ attaining that
maximum.  Then
\[
 \phi(v)\ \ge\ \frac1m\bigl(\log|s_{i_0}|(v)+\lambda_{i_0}\bigr)
        \ =\ \sup_{X^{\An}}\phi-\frac{v(s_{i_0})}{m},
\]
so that $\sup\phi-\phi(v)\le a_m(v)/m\le T_L(v)$, the last step by
\Cref{lem:ch49-vanishing-elementary}\,\cref{itm:ch49-vs:allm}.

\emph{Step 3: passage to all of $\PSH(\omega)$.}  Let
$\phi\in\PSH(\omega)$.  Since $\omega=c_1(L)$ is the class of an ample
line bundle, \Cref{thm:ch34-fs-regularization} supplies a decreasing net
$(\phi_j)$ in $\cH^{\mathrm{gf}}_{\mathbf R}(L)$, each $\phi_j$ being
$\omega$-psh --- \emph{the class does not vary} --- with
$\phi_j\downarrow\phi$ pointwise on $X^{\An}$.  By Step~2,
$\sup\phi_j-\phi_j(v)\le T_L(v)$ for every $j$.  Both
$\sup\phi_j\downarrow\sup\phi$, by
\Cref{lem:ch47-three-normalizations} and the maximum principle
$\sup\psi=\psi(v_{\triv})$, and $\phi_j(v)\downarrow\phi(v)$; hence the
difference converges in $\Rbar$ to $\sup\phi-\phi(v)$, which is
therefore $\le T_L(v)$.  Taking the supremum over $\phi$ gives
$T_\omega(v)\le T_L(v)$.

\emph{Step 4: proper support.}  By \Cref{lem:ch49-am-finite} there are
$m$ and $0\ne s\in H^0(X,mL)$ with $v(s)=+\infty$; then
$\sup\phi_s-\phi_s(v)=+\infty$.
\end{proof}

\begin{remark}[Why the equality is exact, and not up to a constant]
\label{rem:ch49-no-constant}
Two normalizations conspire to remove every constant from
\eqref{eq:ch49-T-equals-T}.  First, the trivial metric on $L$ is the
reference: $\log|s|\le0$ with equality at $v_{\triv}$, by
\eqref{eq:section-log-value}.  Second, $\sup_{X^{\An}}\phi=\phi(v_{\triv})$
for every $\omega$-psh $\phi$, so the drop
$\sup\phi-\phi(v)$ in \eqref{eq:bj-T-valuation} is the drop \emph{from
the vertex of the cone}, which is exactly what $v(s)$ measures.  The fixed-base nontrivially valued theory has no such scaling action
with a trivial-valuation vertex; see \Cref{sec:ch49-part3}.
\end{remark}

\subsection{The calculus of \texorpdfstring{$T_\omega$}{T}}
\label{sec:ch49-T-calculus}

\begin{proposition}
\label{prop:ch49-T-elementary}
Let $v,w\in X^{\An}$ and $t>0$.
\begin{enumerate}
\item\label{itm:ch49-T:homog}
      $T_\omega(tv)=t\,T_\omega(v)$.
\item\label{itm:ch49-T:order}
      If $v\le w$, then $T_\omega(v)\le T_\omega(w)$.
\item\label{itm:ch49-T:lsc}
      $T_\omega$ is lower semicontinuous on $X^{\An}$; consequently
      for every finite $C\ge0$, the sublevel set $\{T_\omega\le C\}$
      is a closed, hence compact, subset of $X^{\An}$ contained in
      $X^{\lin}$.
\item\label{itm:ch49-T:zero}
      $T_\omega(v)=0$ if and only if $v=v_{\triv}$.
\item\label{itm:ch49-T:cone}
      $X^{\lin}$ is stable under the scaling action and is a cone with
      vertex $v_{\triv}$; it is a Borel, indeed
      $\sigma$-compact, subset of $X^{\An}$, and it is dense in
      $X^{\An}$.
\end{enumerate}
\end{proposition}

\begin{proof}
\cref{itm:ch49-T:homog}  For $v\in X^{\val}$ this is
\Cref{lem:ch49-vanishing-elementary}\,\cref{itm:ch49-vs:homog} together with
\eqref{eq:ch49-T-equals-T}; for a proper-support point both sides are
$+\infty$, the scaling action preserving the support
(\Cref{prop:ch32-scaling-action}).  A proof from
\eqref{eq:bj-T-valuation} alone is also available and is worth recording:
the map $\phi\mapsto t^{-1}\cdot\phi$ of
\eqref{eq:function-scaling-action} is a bijection of $\PSH(\omega)$ onto
itself by \Cref{prop:ch34-pl-operations}\,\cref{itm:ch34-plop:const} and
\Cref{thm:bj-basic-closure}, it multiplies suprema by $t^{-1}$ by
\eqref{eq:scaling-fs-weights}, and
$(t^{-1}\cdot\phi)(v)=t^{-1}\phi(tv)$; hence
\[
 \sup\phi-\phi(tv)
 =t\Bigl(\sup(t^{-1}\cdot\phi)-(t^{-1}\cdot\phi)(v)\Bigr),
\]
and taking suprema gives the claim.

\cref{itm:ch49-T:order}  Every $\omega$-psh function is decreasing for
\eqref{eq:semivaluation-order}, so $v\le w$ gives
$\phi(v)\ge\phi(w)$ and hence
$\sup\phi-\phi(v)\le\sup\phi-\phi(w)$.

\cref{itm:ch49-T:lsc}  Each $\phi\in\PSH(\omega)$ is upper
semicontinuous, so $v\mapsto\sup\phi-\phi(v)$ is lower semicontinuous
with values in $[0,+\infty]$; a supremum of lower semicontinuous
functions is lower semicontinuous.  Therefore $\{T_\omega\le C\}$ is
closed in the compact space $X^{\An}$, hence compact, and it is contained
in $X^{\lin}$ by \Cref{prop:ch49-T-equals-T}.

\cref{itm:ch49-T:zero}  $T_\omega(v_{\triv})=0$ because
$\sup\phi=\phi(v_{\triv})$.  Conversely suppose $v\ne v_{\triv}$.  If
$v$ has proper support then $T_\omega(v)=+\infty$.  Otherwise
$v\in X^{\val}$ is a nontrivial valuation of $k(X)$, so $v(f)>0$ for
some $f\in k(X)^\times$.  Writing $f=s_1/s_2$ with
$s_1,s_2\in H^0(X,mL)$ for a suitable $m$, which is possible because
$L$ is ample, we get $v(s_1)=v(f)+v(s_2)>0$, whence
$a_m(v)>0$ and $T_\omega(v)=T_L(v)>0$.

\cref{itm:ch49-T:cone}  Stability under scaling and the cone property are
\cref{itm:ch49-T:homog}; $X^{\lin}=\bigcup_{N\in\mathbf N}\{T_\omega\le N\}$
is $\sigma$-compact by \cref{itm:ch49-T:lsc}, hence Borel.  Density
follows from $X^{\divis}\subset X^{\lin}$
(\Cref{prop:ch49-psef-threshold}) and
\Cref{thm:triv-divisorial-density}.
\end{proof}

\begin{warning}{$T_\omega$ need not be weakly continuous}
\label{warn:ch49-T-not-continuous}
The function is lower semicontinuous for the weak topology and
linear on each scaling ray.  It is not weakly continuous in general.
The following limits distinguish these properties.

On $\mathbf P^1$ with $\omega=c_1(\cO(1))$ the points $t\ord_p$
converge, as $t\to+\infty$, to the proper-support point $v_{p,\triv}$
(\Cref{prop:ch32-scaling-limits}); along that net
$T_\omega(t\ord_p)=t\to+\infty=T_\omega(v_{p,\triv})$, and lower
semicontinuity holds with equality.

By contrast, \Cref{ex:ch49-weak-not-strong} produces a net of divisorial
points $v_S$ on $\mathbf P^1$, all with $T_\omega(v_S)=d$, converging
weakly to $v_{\triv}$, for which $T_\omega(v_{\triv})=0$: the function
$T_\omega$ drops in the limit, so it is not upper semicontinuous and not
continuous.  Lower semicontinuity does pass a uniform upper bound to a weak limit:
if $T_\omega(v_i)\le C$, then $T_\omega(v)\le C$.  It does not prevent
the strict drop in this example.  Strong convergence makes
$T_\omega$ converge because it is the distance to the vertex.
\end{warning}

\begin{example}[The curve case]
\label{ex:ch49-curve-case}
Let $n=1$.  If $X$ is smooth, its analytification is the tree with
one ray for each closed point and common root $v_{\triv}$
(\cite[\S1.1.6]{BJ22}); this includes the projective-line picture
of \Cref{thm:ch32-p1-classification}.  For a singular unibranch
curve, pass to its smooth normalization.  The normalization is a
universal homeomorphism and reflects psh functions, as in
\Cref{lem:ch45-normalization-data,sec:ch45-curve-conventions}.
We write $\ord_p$ for the order at the unique point over $p$.
Thus every nontrivial function-field valuation is $t\ord_p$,
$t>0$, and proper-support points are $v_{p,\triv}$.
Since
$\tau_L(p)=\deg L$ --- the class $L-c\,p$ has degree
$\deg L-c$ --- \Cref{prop:ch49-psef-threshold} gives
\[
 T_\omega(t\ord_p)=t\deg L=tV,
 \qquad
 T_\omega(v_{p,\triv})=+\infty .
\]
Hence $X^{\lin}=X^{\val}=X^{\An}\setminus\{v_{p,\triv}:p\in X(k)\}$: on a
curve, linear growth fails only at the proper-support points.  The curve case has substantial metric and energy content:
$d_{\infty,\omega}$ is the scaled tree metric
\cite[Example~11.8]{BJ22}.  It is also why \eqref{eq:bj-dirac-energy-T} has coinciding constants
there: $\tfrac1{n+1}=\tfrac n{n+1}=\tfrac12$.
\end{example}

\section{Dirac masses and the finite-energy locus}
\label{sec:ch49-dirac}

\subsection{A truncation estimate}
\label{sec:ch49-truncation}

The whole dichotomy rests on one estimate, which says that a normalized
psh function of oscillation at most $a$ cannot lose more than the
fraction $n/(n+1)$ of $a$ in energy.  The centering is essential:
without it, the constant function $\psi\equiv-a$ would contradict the
claimed lower bound.  We therefore state the estimate in a
translation-covariant form.  It is the exact analogue, in this
normalization, of the elementary oscillation estimate in the complex
theory.

\begin{lemma}
\label{lem:ch49-truncation-energy}
Let $\psi\in\PSH(\omega)$, put
$c:=\sup_{X^{\An}}\psi=\psi(v_{\triv})$, and assume that
$c\in\mathbf R$ and, for some $a\ge0$,
\[
                         c-a\ \le\ \psi\ \le\ c.
\]
Then $\psi\in\cE^1(\omega)$ and
\[
 c-\frac{n}{n+1}\,a\ \le\ E_\omega(\psi)\ \le\ c .
\]
More precisely, writing
$\mu_j:=V^{-1}(\omega+\ddc\psi)^j\wedge\omega^{\,n-j}$ for the
probability mixed measures of \Cref{thm:bj-mixed-ma}, one has
$\mu_0=\delta_{v_{\triv}}$ and
\[
 (n+1)\bigl(E_\omega(\psi)-c\bigr)
 =\sum_{j=1}^{n}\int_{X^{\An}}(\psi-c)\,\rd\mu_j .
\]
In particular, one may take
$a=\operatorname{osc}_{X^{\An}}\psi=c-\inf_{X^{\An}}\psi$.  When
$\sup\psi=0$, this gives
$-\frac{n}{n+1}a\le E_\omega(\psi)\le0$.
\end{lemma}

\begin{proof}
We first make explicit the precise results from
\Cref{chap:bj-energy,chap:bj-finite-energy} that enter the argument.  The
normalized energy is monotone and satisfies
\[
 E_\omega(t)=t,
 \qquad
 E_\omega(\eta+t)=E_\omega(\eta)+t,
 \qquad
 \eta\le\chi\Longrightarrow E_\omega(\eta)\le E_\omega(\chi).
\tag{49.6a}\label{eq:ch49-truncation-input-energy}
\]
The extended $(n+1)$-fold pairing is continuous along coordinatewise
decreasing nets.  If $\eta\in\cE^1(\omega)$, its normalized mixed
measures satisfy, for $0\le j\le n$,
\[
 \begin{aligned}
 \nu_j(\eta)&:=V^{-1}(\omega+\ddc\eta)^j\wedge\omega^{\,n-j}\ge0,
 &\nu_j(\eta)(X^{\An})&=1,\\
 \int f\,\rd\nu_j(\eta)
 &=V^{-1}(0,f)(\omega,\eta)^j(\omega,0)^{n-j}
 &&(f\in\cE^1(\omega)).
 \end{aligned}
\tag{49.6b}\label{eq:ch49-truncation-input-mixed}
\]
These are, respectively,
\Cref{prop:ch46-energy-usc,thm:bj-extended-energy-pairing,thm:bj-mixed-ma};
the formula in \eqref{eq:ch49-truncation-input-mixed} also fixes every
normalizing factor needed below.

The bounds on $\psi$, together with
\eqref{eq:ch49-truncation-input-energy}, give
\[
 c-a=E_\omega(c-a)\le E_\omega(\psi)\le E_\omega(c)=c.
\]
Thus $E_\omega(\psi)$ is finite and
$\psi\in\cE^1(\omega)$, without any appeal to an unrecorded
integrability criterion.  We also recall the short proof of the maximum
principle in the form used here.  Choose a decreasing Fubini--Study
regularization $\psi_\alpha\downarrow\psi$ in the fixed ample class
$\omega$.  Each $\psi_\alpha$ attains its maximum at
$v_{\triv}$: in a weighted-section presentation every nonzero section
satisfies $\log|s|\le0$ on $X^{\An}$ and
$\log|s|(v_{\triv})=0$, so the weighted maximum is everywhere at most
its value at $v_{\triv}$.  Consequently
\[
 \sup\psi\le\inf_\alpha\sup\psi_\alpha
 =\inf_\alpha\psi_\alpha(v_{\triv})
 =\psi(v_{\triv})\le\sup\psi,
\]
and hence $c=\psi(v_{\triv})$.

For completeness, we now pass the PL telescoping identity to $\psi$,
rather than merely citing the resulting mixed-energy formula.  At every
Fubini--Study level, multilinearity gives
\[
 (\omega,\psi_\alpha)^{n+1}-(\omega,0)^{n+1}
 =\sum_{j=0}^n
   (0,\psi_\alpha)(\omega,\psi_\alpha)^j(\omega,0)^{n-j}.
\tag{49.6c}\label{eq:ch49-truncation-telescope-pl}
\]
For each $j$, the summand on the right is the difference
\[
 (\omega,\psi_\alpha)^{j+1}(\omega,0)^{n-j}
 -(\omega,\psi_\alpha)^j(\omega,0)^{n-j+1}.
\]
Thus coordinatewise decreasing continuity of the extended pairing passes
\eqref{eq:ch49-truncation-telescope-pl} to the limit without subtracting
two infinite quantities: all limiting pairings are finite because
$\psi$ is bounded.  Moreover $(\omega,0)^{n+1}=0$, since on the
product compactification it is the intersection of $n+1$ pullbacks from
the $n$-fold $X$.  Dividing the limiting
identity by $V$, using
$(n+1)E_\omega(\psi)=V^{-1}(\omega,\psi)^{n+1}$, and then applying
\eqref{eq:ch49-truncation-input-mixed} with $f=\psi$, gives
\[
 (n+1)E_\omega(\psi)
 =\sum_{j=0}^n\int_{X^{\An}}\psi\,\rd\mu_j .
\tag{49.6d}\label{eq:ch49-truncation-mixed-energy}
\]

It remains to identify the zeroth term and estimate the others.  On the
trivial test configuration, the central fibre is one reduced component,
its associated point is $v_{\triv}$, and its intersection coefficient
is $(\omega^n)=V$.  Hence the unnormalized measure $\omega^n$ is
$V\delta_{v_{\triv}}$, so
\[
                         \mu_0=V^{-1}\omega^n=\delta_{v_{\triv}}.
\]
The mass-one assertion in \eqref{eq:ch49-truncation-input-mixed} now gives
$\int c\,\rd\mu_j=c$ for every $j$, while the preceding maximum
principle gives $\int\psi\,\rd\mu_0=c$.  Subtracting $(n+1)c$ from
\eqref{eq:ch49-truncation-mixed-energy} therefore yields
\[
 (n+1)\bigl(E_\omega(\psi)-c\bigr)
 =\sum_{j=1}^n\int_{X^{\An}}(\psi-c)\,\rd\mu_j.
\]
Every integral on the right lies in $[-a,0]$, since
$-a\le\psi-c\le0$ and $\mu_j$ is a probability measure.  Division
by $n+1$ proves the asserted bounds.  Taking the actual oscillation gives
the first specialization.  Finally, $\ddc(\psi+t)=\ddc\psi$ and
\eqref{eq:ch49-truncation-input-energy} is translation covariant, so
adding $t$ changes both $c$ and $E_\omega(\psi)$ by $t$ and
leaves $a$ and all $\mu_j$ unchanged; this gives the supremum-zero
form as well.
\end{proof}

\begin{remark}[Both constants in the lemma are attained]
\label{rem:ch49-truncation-sharp}
The upper bound is attained by every constant function $\psi\equiv c$.
The coefficient in the lower bound is sharp in every dimension.  Indeed,
take $X=\mathbf P^n$, $\omega=c_1(\cO(1))$, with homogeneous
coordinates $z_0,\ldots,z_n$, and set
\[
 \psi:=\max\!\left\{\log|z_1|,
              \max_{j\ne1}\bigl(\log|z_j|-1\bigr)\right\}.
\tag{49.6e}\label{eq:ch49-truncation-sharp-fs}
\]
This is the weighted Fubini--Study function in which $z_1$ has weight
zero and every other coordinate section has weight $-1$.  Since the
coordinate sections generate $\cO(1)$, one has
$\max_j\log|z_j|=0$ on all of $X^{\An}$.  Comparing each weight with
$-1$ and with $0$ therefore gives, pointwise,
\[
 -1=-1+\max_j\log|z_j|\le\psi
 \le\max_j\log|z_j|=0.
\]
At $v_{\triv}$ every $\log|z_j|$ is zero, so $\psi(v_{\triv})=0$
and the maximum principle gives $\sup\psi=0$.

Here is the complete energy computation.  Put
$P=\operatorname{conv}\{0,e_1^\vee,\ldots,e_n^\vee\}$,
$P_{\mathrm{BJ}}=-P$, and
\[
 f_P(w)=\max\{0,-w_1,\ldots,-w_n\},\qquad u=e_1.
\]
On the toric skeleton, the convex function associated to
\eqref{eq:ch49-truncation-sharp-fs} is
\[
 f_\psi(w)=f_P(w-u)-1
 =\max\{-1,-w_1,-1-w_2,\ldots,-1-w_n\}.
\]
Since the conjugate of $f_P$ is zero on $-P$ and $+\infty$
off $-P$, translation of the argument and subtraction of one give
\[
 f_\psi^\vee(m)=1+\langle u,m\rangle=1+m_1
 \quad(m\in-P),
\]
and $f_\psi^\vee=+\infty$ off $-P$.  The normalized toric energy
identity, in the precise convention of this chapter, is
\[
 E_\omega(\eta)=-\frac1{V_P}\int_{-P}f_\eta^\vee(m)\,\rd m,
 \qquad V_P:=\vol(P).
\]
The barycenter of $-P$ is
$-\frac1{n+1}(1,\ldots,1)$; hence
\[
 E_\omega(\psi)
 =-\frac1{V_P}\int_{-P}(1+m_1)\,\rd m
 =-\left(1-\frac1{n+1}\right)
 =-\frac n{n+1}.
\]
Thus the lower coefficient is attained in every dimension.

For comparison, the two surface assertions quoted later in the chapter
can also be checked locally.  On $\mathbf P^2$, let $p$ be a point
and
$\eta:=\max\{\log|\mathfrak m_p|,-1\}$.  On the compactified
deformation to the normal cone of $p$, the function $\eta$ corresponds
to $H-E$, where $H$ is the pullback of the hyperplane class and $E$
is the exceptional divisor.  Here $E\simeq\mathbf P^2$, the restriction
of $H$ to $E$ is trivial, and
$\cO_E(E)=\cO_{\mathbf P^2}(-1)$.  Therefore
\[
 H^3=H^2E=HE^2=0,\qquad E^3=1,\qquad V=1.
\]
Consequently
\[
 E_\omega(\eta)=\frac{(H-E)^3}{3V}=-\frac13.
\]
Since $-1\le\eta\le0$, this is not the lower endpoint $-2/3$ of the
estimate in dimension two.

Finally, let $C\subset\mathbf P^2$ be a smooth conic and put
$\eta_c:=c\max\{\log|\mathcal I_C|,-1\}$.  On the compactified
deformation to the normal cone of $C$, it corresponds to $H-cE$.  If
$\widetilde P$ is the strict transform of the central $\mathbf P^2$
and $F$ its total transform, then
$F=\widetilde P+E$ and $E|_{\widetilde P}=C$.  Thus
\[
 HE\widetilde P=\deg C=2,
 \qquad E^2\widetilde P=C^2=4.
\]
The class $F$ comes from the base, so its restriction to $E$ is
trivial; and $H|_E$ comes from the curve $C$, so $H^2E=0$.  Hence
\[
 0=HEF=HE\widetilde P+HE^2,
 \qquad
 0=E^2F=E^2\widetilde P+E^3,
\]
which gives $HE^2=-2$ and $E^3=-4$.
The endpoint
$\eta_{1/2}=\tfrac12\max\{\log|q|,-1\}$ is $\omega$-psh,
where $q\in H^0(\mathbf P^2,\cO(2))$ defines $C$.
Convex combinations with $0$ give $\eta_c$ for
$0\le c\le\tfrac12$.  The line and exceptional-fibre tests also
give the necessary bounds $1-2c\ge0$ and $c\ge0$.  Moreover
\[
 H^3=H^2E=0,\qquad HE^2=-2,\qquad E^3=-4,
\]
and hence
\[
 E_\omega(\eta_c)=\frac{(H-cE)^3}{3}
 =-2c^2+\frac43c^3.
\]
At $c=\tfrac12$, one has
$-\tfrac12\le\eta_{1/2}\le0$,
$\sup\eta_{1/2}=0$, and
\[
 E_\omega(\eta_{1/2})=-\frac13
 =-\frac23\cdot\frac12.
\]
Thus $\eta_{1/2}$ is another equality case.  If $C$ is not a
coordinate line, then $\ord_C$ is not a toric divisorial valuation for
the standard torus, so this equality case is non-toric in the fixed toric
structure.  These two computations are the ones recorded, in expanded
form, in \Cref{ex:ch46-toric-p2,ex:ch49-conic}.
\end{remark}

\subsection{The dichotomy}
\label{sec:ch49-dichotomy}

\begin{theorem}\label{thm:bj-dirac-finite-energy}
For $v\in X^{\An}$,

\[
 \delta_v\in\mathcal M^1(\omega)
       \quad\Longleftrightarrow\quad v\in X^{\lin}.
\tag{49.7}\label{eq:bj-dirac-e1-iff}
\]

When these conditions hold,

\[
 \frac1{n+1}T_\omega(v)
 \leq E_\omega^\vee(\delta_v)
 \leq\frac n{n+1}T_\omega(v).
\tag{49.8}\label{eq:bj-dirac-energy-T}
\]
\end{theorem}

\begin{proof}
This is \cite[Proposition~11.1]{BJ22}.  The following proof derives
the lower bound from truncation and the sharp upper bound from the
Green function.  The latter route uses the envelope property.
The published proof obtains the upper bound without that property,
by almost-maximizer estimates \cite[Corollary~9.18]{BJ22}.


\emph{Step 1: $E^\vee_\omega(\delta_v)\le T_\omega(v)$, for every
$v\in X^{\An}$.}  Let $\phi\in\cE^1(\omega)$.  Replacing $\phi$ by
$\phi-\sup\phi$, which changes neither side of
\eqref{eq:bj-dual-energy} by
\Cref{prop:ch47-dual-elementary}\,\cref{itm:ch47-dual:nonneg}, we may assume
$\phi\in\cE^1_0(\omega)$.  Then
$E_\omega(\phi)\le\sup_{X^{\An}}\phi=0$ by
\Cref{prop:ch46-energy-usc}, and
$-\phi(v)\le\sup\phi-\phi(v)\le T_\omega(v)$ by
\eqref{eq:bj-T-valuation}.  Hence
$E_\omega(\phi)-\int\phi\,\rd\delta_v=E_\omega(\phi)-\phi(v)
\le T_\omega(v)$, and taking the supremum gives the claim.  In
particular $v\in X^{\lin}$ implies $\delta_v\in\cM^1(\omega)$.

\emph{Step 2: $E^\vee_\omega(\delta_v)\ge\frac1{n+1}T_\omega(v)$, for
every $v\in X^{\An}$.}  Let $\phi\in\PSH^{\sup}(\omega)$, so
$\phi\le0$ and $\phi(v_{\triv})=0$, and let $a\in\mathbf R$ satisfy
$0\le a\le-\phi(v)$.  Put
\[
 \psi:=\max\{\phi,-a\}\in\PSH(\omega),
\]
which is $\omega$-psh by \Cref{thm:bj-basic-closure}\,(2) applied to
$\phi$ and the constant $-a$.  Then
$-a\le\psi\le0$, $\sup\psi=\psi(v_{\triv})=0$, and, since
$\phi(v)\le-a$, also $\psi(v)=-a$.  By
\Cref{lem:ch49-truncation-energy}, $\psi\in\cE^1(\omega)$ and
$E_\omega(\psi)\ge-\frac n{n+1}a$.  Therefore
\[
 E^\vee_\omega(\delta_v)
 \ \ge\ E_\omega(\psi)-\int\psi\,\rd\delta_v
 \ =\ E_\omega(\psi)+a
 \ \ge\ a-\frac n{n+1}a=\frac{a}{n+1}.
\]
Letting $a\uparrow-\phi(v)$ and then taking the supremum over
$\phi\in\PSH^{\sup}(\omega)$ gives
$E^\vee_\omega(\delta_v)\ge\frac1{n+1}T_\omega(v)$, by
\eqref{eq:bj-T-valuation} and \Cref{thm:bj-basic-closure}\,(1).  In
particular $\delta_v\in\cM^1(\omega)$ implies $T_\omega(v)<\infty$,
that is $v\in X^{\lin}$ by \Cref{prop:ch49-T-equals-T}; together with
Step~1 this proves \eqref{eq:bj-dirac-e1-iff} and the left-hand inequality
of \eqref{eq:bj-dirac-energy-T}.

\emph{Step 3: the sharp upper bound.}  Let $v\in X^{\lin}$.  By
\Cref{ex:ch47-dirac-green}, which uses the Calabi--Yau theorem
\Cref{cor:bj-cy-unibranch} and hence the envelope property
(\Cref{imp:ch49-envelope}\,\cref{itm:ch49-env:cy}), there is a unique
$\phi_v\in\cE^1(\omega)$ with
\[
 \MA_\omega(\phi_v)=\delta_v,\qquad\phi_v(v)=0 .
\]
Put $S:=\sup_{X^{\An}}\phi_v$.  Then $0\le S\le T_\omega(v)$, the first
inequality because $\phi_v(v)=0$ and the second by
\eqref{eq:bj-T-valuation}.  Set $\psi:=\phi_v-S\in\cE^1_0(\omega)$, so
that $\MA_\omega(\psi)=\delta_v$ and $\psi(v)=-S$.  Then
\[
 I(\psi,0)=\int(\psi-0)\bigl(\MA_\omega(0)-\MA_\omega(\psi)\bigr)
 =\psi(v_{\triv})-\psi(v)=0+S=S,
\]
using $\MA_\omega(0)=\delta_{v_{\triv}}$, and
\[
 J_0(\psi)=E_\omega(0)-E_\omega(\psi)+\int\psi\,\MA_\omega(0)
 =-E_\omega(\psi)+\psi(v_{\triv})=-E_\omega(\psi),
\]
by \eqref{eq:bj-I-J-potentials} and $E_\omega(0)=0$.  The comparison
\eqref{eq:bj-I-J-comparison} therefore reads
\[
 \frac{S}{n+1}\ \le\ -E_\omega(\psi)\ \le\ \frac{n\,S}{n+1}.
\]
On the other hand \eqref{eq:bj-dual-energy-ma} gives
\[
 E^\vee_\omega(\delta_v)=E^\vee_\omega\bigl(\MA_\omega(\psi)\bigr)
 =E_\omega(\psi)-\int\psi\,\rd\delta_v=E_\omega(\psi)+S .
\]
Combining the last two displays,
\[
 \frac{S}{n+1}\ \le\ E^\vee_\omega(\delta_v)\ \le\ \frac{n\,S}{n+1}
 \ \le\ \frac{n}{n+1}T_\omega(v),
\tag{49.9}\label{eq:ch49-dirac-refined}
\]
which is the right-hand inequality of \eqref{eq:bj-dirac-energy-T}.
\end{proof}

\begin{corollary}
\label{cor:ch49-green-invariants}
Let $v\in X^{\lin}$ and let $\phi_v\in\cE^1(\omega)$ be its Green
function, normalized by $\phi_v(v)=0$.  Then $\phi_v$ is continuous,
\[
 0\le\phi_v\le T_\omega(v),\qquad
 S:=\sup_{X^{\An}}\phi_v=T_\omega(v),
\]
and
\[
 E^\vee_\omega(\delta_v)=E_\omega(\phi_v),\qquad
 \frac{T_\omega(v)}{n+1}\le E^\vee_\omega(\delta_v)
 \le\frac{nT_\omega(v)}{n+1}.
\]
More precisely,
\[
 \phi_v=\sup\{\psi\in\PSH(\omega):\psi(v)\le0\}
 \quad\text{pointwise on }X^{\An}.
 \tag{49.9a}\label{eq:ch49-green-envelope}
\]
For every $w\in X^{\lin}$,
\[
 d_{\infty,\omega}(v,w)
 =\max\{\phi_v(w),\phi_w(v)\}
 =\|\phi_v-\phi_w\|_{C^0(X^{\An})}.
 \tag{49.9b}\label{eq:ch49-green-isometry}
\]
\end{corollary}

\begin{proof}
Continuity is \cite[Theorem~12.12]{BJ22}, already recorded in
\Cref{ex:ch47-dirac-green}, because the measure is the single Dirac
mass $\delta_v$ at finite $T_\omega(v)$.
If $\psi\in\PSH(\omega)$ satisfies $\psi(v)\le0$, then
$\psi\le\phi_v$ almost everywhere for $\MA_\omega(\phi_v)=\delta_v$.
The domination principle \Cref{thm:ch47-domination} therefore gives
$\psi\le\phi_v$ everywhere.  Conversely, $\phi_v$ itself is a
competitor, proving the envelope identity.
The competitor $0$ shows $\phi_v\ge0$.

Every $\eta\in\PSH(\omega)$ is finite at $v$.  Apply that identity
to $\eta-\eta(v)$ and evaluate at the trivial valuation:
\[
 \sup\eta-\eta(v)\le\phi_v(v_{\triv})=\sup\phi_v.
\]
Taking the supremum gives $T_\omega(v)\le S$; choosing $\eta=\phi_v$
gives the reverse inequality.  Thus $S=T_\omega(v)$, and
$0\le\phi_v\le S$.  The dual-energy identity and its bounds follow
from the preceding theorem and \eqref{eq:ch49-dirac-refined}.

The envelope identity also gives
$\eta(w)-\eta(v)\le\phi_v(w)$ for every psh $\eta$.
Exchanging $v,w$ and taking suprema yields the first metric equality;
the reverse bound follows by using $\eta=\phi_v$ or $\phi_w$.
Finally,
$\phi_v-\phi_v(w)\le\phi_w$ and
$\phi_w-\phi_w(v)\le\phi_v$ by the same identity.
Hence $-\phi_w(v)\le\phi_v-\phi_w\le\phi_v(w)$.
The endpoint values at $v,w$ attain the corresponding absolute
bounds, proving the sup-norm equality.
These are precisely \cite[Proposition~12.14(i)--(iv)]{BJ22}.
\end{proof}

\begin{remark}[The Green height is exactly the exhaustion]
\label{rem:ch49-S-vs-T}
The identity $S=T_\omega(v)$ holds for every linear-growth valuation
under the envelope property, not only for the toric examples.
It is the general published statement
\cite[Proposition~12.14(iv)]{BJ22}.
The weaker factor-$n$ bound from the energy comparison remains valid,
but it does not describe the exact Green height.
The isometric embedding in \eqref{eq:ch49-green-isometry} is into
continuous potentials with their sup norm; it is distinct from the
bi-Lipschitz comparison with the measure quasi-metric $I^\vee$.
\end{remark}

\begin{corollary}\label{cor:bj-extremal-measures-dirac}
The extremal points of the convex set $\mathcal M^1(\omega)$ are exactly
$\delta_v$ with $v\in X^{\lin}$.
\end{corollary}

\begin{proof}
Dirac masses are extreme in the full probability simplex, hence in any
convex subset containing them.  The finite-energy Dirac masses are
exactly those in \Cref{thm:bj-dirac-finite-energy}.

For the converse, a non-Dirac Radon probability measure $\mu$ admits
a Borel set $A$ with $0<\mu(A)<1$.  For example, choose disjoint
neighborhoods of two distinct points in its support, both of positive
mass.  The normalized restrictions
\[
 \mu_0=\frac{\mu|_A}{\mu(A)},\qquad
 \mu_1=\frac{\mu|_{A^c}}{\mu(A^c)}
\]
are distinct Radon probabilities and split $\mu$ nontrivially.
For each $\phi\in\cE^1_0(\omega)$, the nonnegative function
$-\phi$ is $\mu$-integrable, and domination gives finite integrals
against both restrictions.  Translations give integrability of every
$\cE^1$ function.  By \cite[Theorem~9.20]{BJ22},
$\mu_0,\mu_1\in\cM^1(\omega)$, so $\mu$ is not extreme.
This is the proof of \cite[Proposition~11.3]{BJ22}.
\end{proof}

\begin{remark}[The precise domination input]
\label{rem:ch49-extreme-import}
If $\nu\le c\mu$ is a probability and $\mu\in\cM^1(\omega)$, then
integrability of $\cE^1$ functions passes to $\nu$.
The converse integrability characterization, Theorem~9.20, converts this
to finite dual energy.  Integrability of a smaller continuous test class
alone would not suffice.  The argument uses that exact published
characterization rather than an unproved uniform bound on its integrals.
\end{remark}

\subsection{Consequences}
\label{sec:ch49-dirac-consequences}

\begin{corollary}
\label{cor:ch49-nonpluripolar}
For $v\in X^{\An}$ the following are equivalent:
\begin{enumerate}
\item\label{itm:ch49-npp:lin} $v\in X^{\lin}$;
\item\label{itm:ch49-npp:T} $T_\omega(v)<\infty$;
\item\label{itm:ch49-npp:dirac} $\delta_v\in\cM^1(\omega)$;
\item\label{itm:ch49-npp:polar} $\{v\}$ is not pluripolar, that is,
      there is no $\phi\in\PSH(\omega)$ with $\phi(v)=-\infty$;
\item\label{itm:ch49-npp:cap} $\Capa^*_\omega(\{v\})>0$.
\end{enumerate}
\end{corollary}

\begin{proof}
\cref{itm:ch49-npp:lin}$\Leftrightarrow$\cref{itm:ch49-npp:T} is
\Cref{prop:ch49-T-equals-T}, and
\cref{itm:ch49-npp:T}$\Leftrightarrow$\cref{itm:ch49-npp:dirac} is
\Cref{thm:bj-dirac-finite-energy}.

\cref{itm:ch49-npp:T}$\Rightarrow$\cref{itm:ch49-npp:polar}: if
$T_\omega(v)<\infty$, then every $\phi\in\PSH(\omega)$ satisfies
$\phi(v)\ge\sup\phi-T_\omega(v)>-\infty$.

\cref{itm:ch49-npp:polar}$\Rightarrow$\cref{itm:ch49-npp:T}: suppose
$T_\omega(v)=+\infty$ and choose $\phi_j\in\PSH^{\sup}(\omega)$ with
$\phi_j(v)\le-2^{j}$.  Setting $c_j:=2^{-j}$, the partial sums
$\psi_N:=\sum_{j\le N}c_j\phi_j$ form a decreasing net of
$\omega$-psh functions --- decreasing because each $\phi_j\le0$, and
$\omega$-psh because $\sum_{j\le N}c_j\le1$ and
\Cref{thm:bj-basic-closure}\,(3) allows the convex combination of
$\phi_1,\dots,\phi_N$ and the constant $0$ with total weight one.
Their limit $\psi:=\sum_{j\ge1}c_j\phi_j$ satisfies
$\psi(v_{\triv})=0$, so it is not identically $-\infty$ and is
$\omega$-psh by \Cref{thm:bj-basic-closure}\,(4); and
$\psi(v)\le\sum_j 2^{-j}(-2^{j})=-\infty$.

\cref{itm:ch49-npp:polar}$\Leftrightarrow$\cref{itm:ch49-npp:cap} is the
equivalence \eqref{eq:bj-small-set-equivalence} of
\Cref{thm:bj-small-set-equivalence}, applied to the compact set
$\{v\}$; it is the only place where
\Cref{imp:ch49-envelope}\,\cref{itm:ch49-env:polar} is used.  Alternatively,
and without that import, the implication
\cref{itm:ch49-npp:T}$\Rightarrow$\cref{itm:ch49-npp:cap} follows from the
compact Alexander--Taylor lower bound proved in
\Cref{prop:ch48-at-compact}, applied to $A=\{v\}$, for which
$T_\omega(A)=T_\omega(v)$ by
\eqref{eq:bj-alexander-taylor-height}.
\end{proof}

\begin{corollary}
\label{cor:ch49-mass-in-lin}
Let $\mu\in\cM^1(\omega)$.  Every atom of $\mu$ belongs to $X^{\lin}$.
In particular, $\mu(X^{\lin})=1$ if $\mu$ is purely atomic, or if $X$ is
a smooth projective curve over an algebraically closed field.
For an arbitrary Radon probability measure, the additional condition
$\int T_\omega\,\rd\mu<\infty$ also implies $\mu(X^{\lin})=1$.
\end{corollary}

\begin{proof}
If $v\notin X^{\lin}$, then $\{v\}$ is pluripolar by
\Cref{cor:ch49-nonpluripolar}.  By \cite[Corollary~7.48]{BJ22}, there
is a finite-energy psh function $u$ with $u(v)=-\infty$.
\Cref{prop:ch47-no-polar-mass} gives $\mu(\{v\})=0$.
This proves the assertion about atoms and its purely atomic consequence.

For a smooth curve, the complement of $X^{\lin}$ consists of the
proper-support endpoints of the branches.
It is Borel, since
$X^{\lin}=\bigcup_{N\geq1}\{T_\omega\leq N\}$ and $T_\omega$ is lower
semicontinuous.  Every compact subset $K$ of this complement is finite:
the open branches, with their common trivial valuation removed, give
an open cover of $K$, and each branch contains just one endpoint.
Compactness gives a finite subcover.  Each such $K$ has measure zero by
the first assertion, so inner regularity of $\mu$ gives
$\mu(X^{\An}\setminus X^{\lin})=0$.

Finally, if $M:=\int T_\omega\,\rd\mu<\infty$, then
$\mu(\{T_\omega>N\})\leq M/N$ for every $N>0$.
Letting $N\to\infty$ gives $\mu(\{T_\omega=\infty\})=0$, as required.
\end{proof}

\begin{example}
\label{ex:ch49-uniform-tail-capacity}
Uniform tails of the pointwise height need not have small capacity.
Take $X=\mathbf P^1$ over an algebraically closed field and
$\omega=c_1(\cO(1))$.  For $R>0$, set
$B_R=\{v:T_\omega(v)>R\}$.
Choose an integer $m>R$ and distinct points $p_1,\dots,p_m$.
Define $\psi_m$ to vanish off their branches and put
\[
 \psi_m(t\,\ord_{p_i})=-\min\{t/m,1\},\qquad 1\leq i\leq m.
\]
Assign value $-1$ at the selected proper-support endpoints.
This gives a continuous PL function with $-1\leq\psi_m\leq0$.
Each branch profile is convex, and the sum of its initial outgoing
slopes is $-m/m=-1$, so the curve criterion in
\Cref{ex:ch34-psh-check} gives $\psi_m\in\PSH(\omega)$.
Equivalently, if $s$ is a section of $\cO(m)$ with simple zeros
$p_1,\dots,p_m$, then
$\psi_m=\max\{m^{-1}\log|s|,-1\}$.
The graph Laplacian gives
\[
 \MA_\omega(\psi_m)
   =\frac1m\sum_{i=1}^m\delta_{m\,\ord_{p_i}}.
\]
There is no remaining atom at $v_{\triv}$: its mass is
$1-m/m=0$.  Since $T_\omega(m\,\ord_{p_i})=m>R$, all the mass lies in
$B_R$.  The defining capacity test therefore proves
\[
 \Capa_\omega(B_R)=\Capa_\omega^*(B_R)=1
 \qquad\text{for every }R>0.
\]
In particular, the upper bound $1/R$ fails when $R>1$.
The estimate in \Cref{prop:ch48-at-upper} uses the height of a
\emph{set}, not the infimum of its pointwise heights.  These quantities
have different orders of supremum and infimum; see
\eqref{eq:bj-alexander-taylor-height}.
This example concerns a family of measures varying with $R$.  It does
not contradict the full-mass statement for any fixed finite-energy
measure on a curve proved above.
\end{example}

\begin{remark}
\label{rem:ch49-consistency-ch46}
The conclusions needed for integration in \Cref{chap:bj-energy} concern
each finite collection of psh functions.  Their polar sets are null,
so their differences have well-defined integrable representatives,
by \Cref{thm:bj-mixed-ma,prop:ch46-no-mass-outside-lin}.
The atomic conclusion of \Cref{cor:ch49-mass-in-lin} and the full-mass
conclusion under its stated additional hypotheses are compatible with
this construction.  A general assertion that every finite-energy measure
is carried by $X^{\lin}$ does not follow from the individual null-set
statements, or from the capacity argument excluded by
\Cref{ex:ch49-uniform-tail-capacity}; it is not used here.
\end{remark}

\section{Weak and strong topologies on
\texorpdfstring{$X^{\lin}$}{the linear-growth locus}}
\label{sec:ch49-topologies}

\subsection{The metric \texorpdfstring{$d_{\infty,\omega}$}{d infinity}}
\label{sec:ch49-dinfty}

\begin{definition}[The $d_\infty$-metric]
\label{def:bj-dinfty}
For $v,w\in X^{\lin}$, set

\[
 d_{\infty,\omega}(v,w):=
 \sup_{\phi\in\PSH(\omega)}|\phi(v)-\phi(w)|.
\tag{49.10}\label{eq:bj-dinfty}
\]
\end{definition}

The expression is finite because it is at most
$T_\omega(v)+T_\omega(w)$.  Psh functions separate valuations, so it is
a metric.  It is the smallest metric for which every $\omega$-psh
function is $1$-Lipschitz, and

\[
 d_{\infty,\omega}(v,v_{\triv})=T_\omega(v),\qquad
 d_{\infty,\omega}(tv,tw)=t\,d_{\infty,\omega}(v,w).
\tag{49.11}\label{eq:bj-dinfty-homogeneity}
\]

For a $\mathbf Q$-line bundle $L$, there is also the section formula

\[
 d_{\infty,L}(v,w)
 =\sup_{\substack{m\geq1\\0\ne s\in H^0(X,mL)}}
      m^{-1}|v(s)-w(s)|.
\tag{49.12}\label{eq:bj-dinfty-sections}
\]

This is \cite[Lemma~11.7]{BJ22}.  The class inclusions proved next
show that changing the ample class changes $d_\infty$ only up to
bi-Lipschitz equivalence.

We prove all of this, and give the change-of-polarization statement a proof
that does not use \eqref{eq:bj-dinfty-sections}.

\begin{proposition}
\label{prop:ch49-dinfty-basic}
Let $v,w,u\in X^{\lin}$ and $t>0$.
\begin{enumerate}
\item\label{itm:ch49-dinf:finite}
      $d_{\infty,\omega}(v,w)\le T_\omega(v)+T_\omega(w)<\infty$, and the
      supremum in \eqref{eq:bj-dinfty} is unchanged if it is restricted to
      $\PSH^{\sup}(\omega)$, or to $\cE^1_0(\omega)$, or to the
      Fubini--Study functions $\cH^{\mathrm{gf}}_{\mathbf R}(L)$.
\item\label{itm:ch49-dinf:metric}
      $d_{\infty,\omega}$ is a metric on $X^{\lin}$.
\item\label{itm:ch49-dinf:smallest}
      It is the smallest pseudometric on $X^{\lin}$ for which every
      $\omega$-psh function is $1$-Lipschitz.
\item\label{itm:ch49-dinf:vertex}
      $d_{\infty,\omega}(v,v_{\triv})=T_\omega(v)$; consequently
      $|T_\omega(v)-T_\omega(w)|\le d_{\infty,\omega}(v,w)$ and
      $T_\omega$ is $1$-Lipschitz.
\item\label{itm:ch49-dinf:homog}
      $d_{\infty,\omega}(tv,tw)=t\,d_{\infty,\omega}(v,w)$.
\item\label{itm:ch49-dinf:monotone}
      If $v\le w\le w'\le u$ in the order
      \eqref{eq:semivaluation-order}, then
      $d_{\infty,\omega}(w,w')\le d_{\infty,\omega}(v,u)$.
\item\label{itm:ch49-dinf:polarization}
      If $L'$ is another ample line bundle and $p,q\in\mathbf Z_{>0}$
      are such that $pL-L'$ and $qL'-L$ are ample, then
      \[
       q^{-1}d_{\infty,\omega}\ \le\ d_{\infty,\omega'}\ \le\
       p\,d_{\infty,\omega},
       \qquad \omega'=c_1(L') .
      \]
      In particular the topology and the uniform structure of
      $d_{\infty,\omega}$ do not depend on the polarization.
\end{enumerate}
\end{proposition}

\begin{proof}
\cref{itm:ch49-dinf:finite}  For $\phi\in\PSH(\omega)$,
$\phi(v)-\phi(w)=(\sup\phi-\phi(w))-(\sup\phi-\phi(v))$, a difference of
two numbers in $[0,T_\omega(w)]$ and $[0,T_\omega(v)]$; hence
$|\phi(v)-\phi(w)|\le T_\omega(v)+T_\omega(w)$.  Adding a constant to
$\phi$ does not change $|\phi(v)-\phi(w)|$, so the supremum may be
restricted to $\PSH^{\sup}(\omega)$.  It may be restricted further to
$\cE^1_0(\omega)$: given $\phi\in\PSH^{\sup}(\omega)$, the truncations
$\phi_m:=\max\{\phi,-m\}$ are bounded, hence lie in
$\cE^1_0(\omega)$, and $\phi_m(v)\to\phi(v)$,
$\phi_m(w)\to\phi(w)$ as $m\to\infty$, both values being finite.
Finally, a decreasing Fubini--Study regularization
\Cref{thm:ch34-fs-regularization} of $\phi$ converges pointwise at
$v$ and at $w$, so the supremum over
$\cH^{\mathrm{gf}}_{\mathbf R}(L)$ is already the full one.

\cref{itm:ch49-dinf:metric}  Symmetry and the triangle inequality are
immediate from the definition as a supremum of $|\phi(v)-\phi(w)|$.  For
separation, suppose $\phi(v)=\phi(w)$ for every $\phi\in\PSH(\omega)$.
Taking $\phi=m^{-1}\log|s|$ for $0\ne s\in H^0(X,mL)$ and using
\eqref{eq:section-log-value} gives $v(s)=w(s)$ for all such $s$.  Since
$L$ is ample, every element of $k(X)^\times$ is a quotient
$s_1/s_2$ of two nonzero sections of a common $mL$; hence
$v(f)=v(s_1)-v(s_2)=w(s_1)-w(s_2)=w(f)$ for every $f\in k(X)^\times$,
that is $v=w$.

\cref{itm:ch49-dinf:smallest}  If $d'$ is a pseudometric for which every
$\omega$-psh function is $1$-Lipschitz, then
$|\phi(v)-\phi(w)|\le d'(v,w)$ for every $\phi$, so
$d_{\infty,\omega}\le d'$; and $d_{\infty,\omega}$ itself has the
property, by definition.

\cref{itm:ch49-dinf:vertex}  By the maximum principle
$\phi(v_{\triv})=\sup\phi$, so
$|\phi(v)-\phi(v_{\triv})|=\sup\phi-\phi(v)$, and taking suprema gives
$d_{\infty,\omega}(v,v_{\triv})=T_\omega(v)$ by
\eqref{eq:bj-T-valuation}.  The Lipschitz statement is the triangle
inequality.

\cref{itm:ch49-dinf:homog}  As in the proof of
\Cref{prop:ch49-T-elementary}\,\cref{itm:ch49-T:homog}: the map
$\phi\mapsto t^{-1}\cdot\phi$ is a bijection of $\PSH(\omega)$, and
$|\phi(tv)-\phi(tw)|
=t\,|(t^{-1}\cdot\phi)(v)-(t^{-1}\cdot\phi)(w)|$.

\cref{itm:ch49-dinf:monotone}  Every $\omega$-psh function is decreasing
for \eqref{eq:semivaluation-order}, so
$\phi(v)\ge\phi(w)\ge\phi(w')\ge\phi(u)$ and therefore
$|\phi(w)-\phi(w')|\le|\phi(v)-\phi(u)|$.

\cref{itm:ch49-dinf:polarization}  Let $\phi\in\PSH(\omega')$.  Since
$\omega-p^{-1}\omega'=p^{-1}(p\omega-\omega')$ is ample, the zero
function is $(\omega-p^{-1}\omega')$-psh by
\Cref{lem:ch34-psh-nonempty}, and
\Cref{thm:bj-basic-closure}\,(3) applied with the coefficients
$p^{-1}$ and $1$ shows that $p^{-1}\phi$ is $\omega$-psh.
Consequently
\[
 |\phi(v)-\phi(w)|
 =p\,\bigl|(p^{-1}\phi)(v)-(p^{-1}\phi)(w)\bigr|
 \le p\,d_{\infty,\omega}(v,w),
\]
and taking the supremum over $\phi$ gives
$d_{\infty,\omega'}\le p\,d_{\infty,\omega}$.  Exchanging the roles of
$L$ and $L'$ gives $d_{\infty,\omega}\le q\,d_{\infty,\omega'}$.
Note that $X^{\lin}$ is the same set for the two polarizations, by
\Cref{prop:ch49-change-polarization}.
\end{proof}

\begin{remark}[The section formula, and what it adds]
\label{rem:ch49-section-formula}
The inequality $\ge$ in \eqref{eq:bj-dinfty-sections} is immediate from
\Cref{prop:ch49-dinfty-basic}\,\cref{itm:ch49-dinf:finite}: the function
$\phi=m^{-1}\log|s|$ is $\omega$-psh with
$\phi(v)-\phi(w)=m^{-1}(w(s)-v(s))$.  The inequality $\le$ holds for
every Fubini--Study function, by the elementary estimate
$\max_ia_i-\max_ib_i\le\max_i(a_i-b_i)$ applied to
$a_i=\lambda_i-v(s_i)$ and $b_i=\lambda_i-w(s_i)$.  For $\phi\in\PSH(\omega)$, choose a decreasing regularization
$\phi_j\in\cH_{\mathrm{dom}}(L)\subset\cH(L)$ in the fixed class.
Both values $\phi_j(v)$ and $\phi_j(w)$ converge to finite values.
The finite-level bound therefore passes to $\phi$.  Taking the
supremum proves \eqref{eq:bj-dinfty-sections}.  This is the fixed-class
argument of \cite[Lemma~11.7 and Theorem~4.15]{BJ22}.
For a $\mathbf Q$-line bundle, restrict to degrees clearing its denominator.
\end{remark}

\subsection{The multiplicative comparison lemma}
\label{sec:ch49-comparison-lemma}

Every approximation argument in this chapter rests on the following
estimate, which converts a \emph{multiplicative} two-sided bound between
valuations into an \emph{additive} bound for $d_{\infty,\omega}$.  It is
the reason quasimonomial and divisorial approximation works, and it is the
only place where the convexity of psh functions along scaling rays
(\Cref{prop:ch35-scaling-convex}) is used.

\begin{lemma}
\label{lem:ch49-multiplicative-comparison}
Let $v\in X^{\lin}$, $w\in X^{\val}$ and $\lambda\ge1$, and suppose
\[
 v\ \le\ w\ \le\ \lambda v
\]
in the order \eqref{eq:semivaluation-order}.  Then $w\in X^{\lin}$ and
\[
 0\ \le\ \phi(v)-\phi(w)\ \le\ (\lambda-1)\,T_\omega(v)
 \quad\text{for every }\phi\in\PSH^{\sup}(\omega),
\]
so that
\[
 d_{\infty,\omega}(v,w)\ \le\ (\lambda-1)\,T_\omega(v).
\tag{49.13}\label{eq:ch49-multiplicative-comparison}
\]
More generally, if $\lambda^{-1}v\le w\le\lambda v$ with
$\lambda\ge1$, then
\[
 d_{\infty,\omega}(v,w)\le(\lambda-1)\max\{T_\omega(v),T_\omega(w)\}.
\]
\end{lemma}

\begin{proof}
That $w\in X^{\lin}$ follows from $w\le\lambda v$,
\Cref{prop:ch49-T-elementary}\,\cref{itm:ch49-T:order} and
\cref{itm:ch49-T:homog}: $T_\omega(w)\le T_\omega(\lambda v)
=\lambda T_\omega(v)<\infty$.

Let $\phi\in\PSH^{\sup}(\omega)$, so that $\phi\le0$ and
$\phi(v_{\triv})=0$.  By \Cref{prop:ch35-scaling-convex} the function
\[
 g:[0,+\infty)\longrightarrow\mathbf R\cup\{-\infty\},
 \qquad g(t):=\phi(tv),
\]
is convex and nonincreasing, with $g(0)=\phi(v_{\triv})=0$ because
$0\cdot v=v_{\triv}$.  A convex function on $[0,\infty)$ vanishing at
the origin has nondecreasing difference quotient $t\mapsto g(t)/t$;
hence $g(\lambda)\ge\lambda\,g(1)$ for $\lambda\ge1$, that is
\[
 \phi(\lambda v)\ \ge\ \lambda\,\phi(v).
\]
Since $\phi$ is decreasing for \eqref{eq:semivaluation-order} and
$v\le w\le\lambda v$, we get
$\phi(v)\ge\phi(w)\ge\phi(\lambda v)\ge\lambda\phi(v)$, whence
\[
 0\ \le\ \phi(v)-\phi(w)\ \le\ \phi(v)-\lambda\phi(v)
 =(\lambda-1)\bigl(-\phi(v)\bigr)\ \le\ (\lambda-1)T_\omega(v),
\]
the last step by \eqref{eq:bj-T-valuation}.  Taking the supremum over
$\phi$ and using
\Cref{prop:ch49-dinfty-basic}\,\cref{itm:ch49-dinf:finite} gives
\eqref{eq:ch49-multiplicative-comparison}.

For the last assertion, fix $\phi\in\PSH^{\sup}(\omega)$.  If
$\phi(v)\ge\phi(w)$, then from $w\le\lambda v$ and
$\phi(\lambda v)\ge\lambda\phi(v)$ we get
$\phi(v)-\phi(w)\le\phi(v)-\phi(\lambda v)\le(\lambda-1)T_\omega(v)$.  If
$\phi(w)\ge\phi(v)$, the symmetric computation with $v\le\lambda w$
gives $\phi(w)-\phi(v)\le(\lambda-1)T_\omega(w)$.  In both cases
$|\phi(v)-\phi(w)|\le(\lambda-1)\max\{T_\omega(v),T_\omega(w)\}$.
\end{proof}

\begin{remark}[The lemma is sharp, and why the estimate is multiplicative]
\label{rem:ch49-comparison-sharp}
Take $v\in X^{\lin}$ and $w=\lambda v$.  Then
$d_{\infty,\omega}(v,\lambda v)\ge
|T_\omega(\lambda v)-T_\omega(v)|=(\lambda-1)T_\omega(v)$ by
\Cref{prop:ch49-dinfty-basic}\,\cref{itm:ch49-dinf:vertex}, so
\eqref{eq:ch49-multiplicative-comparison} is an equality in that case.  On a common $T_\omega$-ball of radius $C$, define Thompson's
extended part metric by
\[
 \rho(v,w)=\log\inf\{\lambda\ge1:\lambda^{-1}v\le w\le\lambda v\}.
\]
The lemma yields $d_{\infty,\omega}(v,w)\le C(e^{\rho(v,w)}-1)$
when $\rho$ is finite, and always $d_{\infty,\omega}(v,w)\le2C$.
For $\rho\le1$, convexity of the exponential gives
$e^\rho-1\le(e-1)\rho$; for $\rho\ge1$, use $2C\le2C\rho$.
Hence $d_{\infty,\omega}\le2C\rho$ on each such ball.
The estimate concerns this order metric; no ambient vector-space norm
on valuations has been defined.
\end{remark}

\subsection{The topological theorem}
\label{sec:ch49-topology-theorem}

\noindent\begin{minipage}{\linewidth}
\begin{theorem}\label{thm:bj-linear-topology-complete}
The following hold.

\begin{enumerate}
\item\label{itm:ch49-top:pointwise}
      The topology of $\PSH(\omega)$, originally defined by pointwise
      convergence on $X^{\divis}$, is also pointwise convergence on
      $X^{\lin}$.
\item\label{itm:ch49-top:complete}
      The metric space $(X^{\lin},d_\infty)$ is complete.
\item\label{itm:ch49-top:embedding}
      The map $v\mapsto\delta_v$ is a bi-Lipschitz embedding
\[
 (X^{\lin},d_\infty)\hookrightarrow(\mathcal M^1,I^\vee).
\tag{49.14}\label{eq:bj-valuations-measures-embedding}
\]
\item\label{itm:ch49-top:criterion}
      A net $v_i$ converges strongly to $v$ if and only if
      $v_i\to v$ weakly in $X^{\An}$ and
      $E^\vee(\delta_{v_i})\to E^\vee(\delta_v)$.
\end{enumerate}
\par
\end{theorem}
\end{minipage}\par

\begin{proof}[Published interfaces and proof order]
The four assertions are \cite[Theorems~11.4, 11.9, 11.13 and
Corollary~11.15]{BJ22}.  The proof below derives the upper embedding
bound, criterion, completeness, and pointwise assertion.  The lower
embedding estimate is the exact interface of Theorem~11.13.
That theorem uses the section-weighted maximizing sequence of
Lemma~11.17, strong convergence of its MA measures, and the quantitative
pairing estimates.  Its lower bound is not inferred merely from
point separation.
\end{proof}

We first prove the elementary upper comparison.  The subsequent proof
uses fixed-measure continuity for the pointwise assertion in every
characteristic; divisorial strong density is not a dependency.

\begin{proposition}
\label{prop:ch49-embedding-easy}
For $v,w\in X^{\lin}$,
\[
 \bigl|E^\vee_\omega(\delta_v)-E^\vee_\omega(\delta_w)\bigr|
 \ \le\ d_{\infty,\omega}(v,w)
 \qquad\text{and}\qquad
 I^\vee\bigl(\delta_v,\delta_w\bigr)\ \le\ 2\,d_{\infty,\omega}(v,w).
\tag{49.15}\label{eq:ch49-embedding-easy}
\]
\end{proposition}

\begin{proof}
For the first estimate, write
$E^\vee_\omega(\delta_v)=\sup_{\phi\in\cE^1}(E_\omega(\phi)-\phi(v))$ by
\eqref{eq:bj-dual-energy}.  For each $\phi$,
\[
 E_\omega(\phi)-\phi(v)
 \le E_\omega(\phi)-\phi(w)+|\phi(v)-\phi(w)|
 \le E_\omega(\phi)-\phi(w)+d_{\infty,\omega}(v,w),
\]
whence $E^\vee_\omega(\delta_v)\le E^\vee_\omega(\delta_w)
+d_{\infty,\omega}(v,w)$; exchanging $v$ and $w$ gives the claim.
Note that all the quantities are finite by
\Cref{thm:bj-dirac-finite-energy}.

For the second, let $\phi_v\in\cE^1(\omega)$ be the Green function of
$v$, normalized by $\phi_v(v)=0$
(\Cref{ex:ch47-dirac-green}).  By
\Cref{prop:ch48-defect-elementary} the Fenchel defect
\eqref{eq:bj-fenchel-defect} satisfies
$J_{\delta_v}(\phi_v)=0$, because $\MA_\omega(\phi_v)=\delta_v$.
Hence, by the definition \eqref{eq:bj-measure-Idual} of $I^\vee$ as an
infimum,
\[
 I^\vee(\delta_v,\delta_w)
 \ \le\ J_{\delta_v}(\phi_v)+J_{\delta_w}(\phi_v)
 \ =\ J_{\delta_w}(\phi_v)
 \ =\ E^\vee_\omega(\delta_w)-E_\omega(\phi_v)+\phi_v(w).
\]
By \Cref{cor:ch49-green-invariants},
$E_\omega(\phi_v)=E^\vee_\omega(\delta_v)$; and $\phi_v(v)=0$.
Therefore
\[
 I^\vee(\delta_v,\delta_w)
 \le\bigl(E^\vee_\omega(\delta_w)-E^\vee_\omega(\delta_v)\bigr)
   +\bigl(\phi_v(w)-\phi_v(v)\bigr)
 \le d_{\infty,\omega}(v,w)+d_{\infty,\omega}(v,w),
\]
the first bracket by the first estimate and the second because $\phi_v$
is $1$-Lipschitz for $d_{\infty,\omega}$.
\end{proof}

\begin{proof}[Proof of \Cref{thm:bj-linear-topology-complete} with its stated source input]
\emph{\cref{itm:ch49-top:embedding}.}  The upper bound
\[
 I^\vee(\delta_v,\delta_w)\le2\,d_{\infty,\omega}(v,w)
\]
is \Cref{prop:ch49-embedding-easy}, and the lower bound
$a_n\,d_{\infty,\omega}(v,w)\le I^\vee(\delta_v,\delta_w)$
is the imported estimate
\Cref{imp:ch49-imported}\,\cref{itm:ch49-import:embedding}.
Injectivity is
\Cref{prop:ch49-dinfty-basic}\,\cref{itm:ch49-dinf:metric}, and
$\delta_v\in\cM^1(\omega)$ by \Cref{thm:bj-dirac-finite-energy}.

\emph{\cref{itm:ch49-top:criterion}.}  We first record that, for
$v_i,v\in X^{\An}$, the net $\delta_{v_i}$ converges weakly to
$\delta_v$ if and only if $v_i\to v$ in $X^{\An}$.  Indeed weak
convergence means $f(v_i)\to f(v)$ for every $f\in C^0(X^{\An})$, and
on a compact Hausdorff space the continuous functions separate points and
generate the topology.

Suppose $d_{\infty,\omega}(v_i,v)\to0$.  By
\Cref{prop:ch49-embedding-easy}, $I^\vee(\delta_{v_i},\delta_v)\to0$, so
by \Cref{thm:bj-strong-measures}, in the implication
\cref{itm:ch48-sm:gauge}$\Rightarrow$\cref{itm:ch48-sm:weak}, the measures converge weakly and their dual
energies converge; by the previous paragraph this is exactly
$v_i\to v$ weakly together with
$E^\vee_\omega(\delta_{v_i})\to E^\vee_\omega(\delta_v)$.  Conversely, if
those two conditions hold, then \Cref{thm:bj-strong-measures}, in the
implication \cref{itm:ch48-sm:weak}$\Rightarrow$\cref{itm:ch48-sm:gauge}, gives
$I^\vee(\delta_{v_i},\delta_v)\to0$, and the imported lower bound of
\cref{itm:ch49-top:embedding} gives $d_{\infty,\omega}(v_i,v)\to0$.

\emph{\cref{itm:ch49-top:complete}.}  Let $(v_i)_{i\in I}$ be a
$d_{\infty,\omega}$-Cauchy net in $X^{\lin}$.  It is eventually bounded: fixing
$i_0$ with $d_{\infty,\omega}(v_i,v_j)\le1$ for $i,j\ge i_0$ and
using \Cref{prop:ch49-dinfty-basic}\,\cref{itm:ch49-dinf:vertex}, we get
\[
 T_\omega(v_i)\le T_\omega(v_{i_0})+1=:C
 \qquad\text{for all }i\ge i_0 .
\]
The set $K:=\{T_\omega\le C\}$ is a compact subset of $X^{\lin}$ by
\Cref{prop:ch49-T-elementary}\,\cref{itm:ch49-T:lsc}, and the tail of the
net lies in $K$; hence some subnet $(v_{i_\alpha})$ converges in
$X^{\An}$ to a point $v\in K\subset X^{\lin}$.

We claim $d_{\infty,\omega}(v_i,v)\to0$.  Let $\varepsilon>0$ and
choose $i_1\ge i_0$ with $d_{\infty,\omega}(v_i,v_j)\le\varepsilon$ for
$i,j\ge i_1$.  Fix $i\ge i_1$ and $\phi\in\PSH(\omega)$.  For every
$\alpha$ with $i_\alpha\ge i_1$,
$|\phi(v_i)-\phi(v_{i_\alpha})|\le\varepsilon$.  By
\Cref{prop:bj-psh-continuous-bounded-valuations} the function $\phi$ is
weakly continuous on $K$, and $v_{i_\alpha}\to v$ inside $K$;
therefore $\phi(v_{i_\alpha})\to\phi(v)$ and
$|\phi(v_i)-\phi(v)|\le\varepsilon$.  Taking the supremum over $\phi$
gives $d_{\infty,\omega}(v_i,v)\le\varepsilon$ for every $i\ge i_1$.

\emph{\cref{itm:ch49-top:pointwise}.}
Pointwise convergence on $X^{\lin}$ implies convergence on $X^{\divis}$.
Conversely, let $\phi_i\to\phi$ in the psh weak topology and fix
$v\in X^{\lin}$.  The maxima $c_i=\sup\phi_i$ converge to
$c=\sup\phi$ by \Cref{thm:ch34-psh-compactness}.  Since
\[
 c_i-T_\omega(v)\le\phi_i(v)\le c_i,
\]
these values and $\phi(v)$ have a common finite lower bound on a tail.
Choose $t>0$ such that they all exceed $-t$ there.
The truncations $\psi_i=\max\{\phi_i,-t\}$ and
$\psi=\max\{\phi,-t\}$ preserve their values at $v$ and converge
pointwise on $X^{\divis}$.  They are bounded above uniformly on the
tail and below by $-t$, so their normalized J-values are uniformly
bounded by monotonicity of energy.  By
\cite[Proposition~9.19]{BJ22}, applied to the fixed finite-energy
measure $\delta_v$,
\[
 \psi_i(v)=\int\psi_i\,\rd\delta_v\longrightarrow
 \int\psi\,\rd\delta_v=\psi(v).
\]
Thus $\phi_i(v)\to\phi(v)$.  This is the argument of
\cite[Theorem~11.4]{BJ22}, and it has no characteristic restriction.
\end{proof}

\begin{remark}[The two topologies on $X^{\lin}$, and the third]
\label{rem:ch49-three-topologies}
Three descriptions of topology on $X^{\lin}$ have now been named.  The
strong topology and the trace topology coincide.  The weak topology
can be strictly coarser, already on $\mathbf P^1$.
\begin{enumerate}
\item\label{itm:ch49-3top:weak}
      The \emph{weak} topology, induced from $X^{\An}$.  $X^{\lin}$ is not closed in the compact analytification:
      on a curve the net $t\,\ord_p$ converges to its proper-support
      endpoint as $t\to\infty$.  Thus the inherited compact-space
      uniformity is not complete.  Moreover,
      $T_\omega$ is only lower semicontinuous for it
      (\Cref{warn:ch49-T-not-continuous}).
\item\label{itm:ch49-3top:strong}
      The \emph{strong} topology, that of $d_{\infty,\omega}$.  It is
      complete (\Cref{thm:bj-linear-topology-complete}\,\cref{itm:ch49-top:complete}),
      it refines the weak topology, and $T_\omega$ is
      $1$-Lipschitz for it.
\item\label{itm:ch49-3top:trace}
      The trace of the strong topology of $\cM^1(\omega)$ under
      $v\mapsto\delta_v$.  By
      \Cref{thm:bj-linear-topology-complete}\,\cref{itm:ch49-top:embedding}
      this coincides with \cref{itm:ch49-3top:strong}, and that coincidence
      is the whole point of the theorem.
\end{enumerate}
Weak and strong topology need not agree on a compact $T$-ball:
\Cref{ex:ch49-weak-not-strong} gives a counterexample on $\mathbf P^1$
for every positive radius after scaling.  Continuity of each fixed psh
function there does not imply equicontinuity of the whole family.
On an SNC dual complex the stronger published equicontinuity statement
does hold, including at boundaries; see
\Cref{prop:ch49-dual-complex}\,\cref{itm:ch49-dc:compact}.
\end{remark}

\begin{remark}[Question 11.16 answered negatively]
\label{rem:ch49-open-question}
BJ22 Question~11.16 asks whether $E^\vee(\delta_v)$ in part (4) may be
replaced by $T_\omega(v)$.  The implication is false by
\cref{thm:ch49-global-radon-riesz-counterexample}: the exhaustion is
identically one along a weakly convergent sequence and at its limit, while
the strong distance stays at least one.  This book-level counterexample is
not attributed to \cite{BJ22}; that source supplies the definitions, the
strong metric, and the question.
\end{remark}

\begin{remark}[What comparability does and does not give]
\label{rem:ch49-open-question-detail}
Write $A(v):=E^\vee_\omega(\delta_v)$ and $B(v):=T_\omega(v)$.  By
\eqref{eq:bj-dirac-energy-T}, $\tfrac1{n+1}B\le A\le\tfrac n{n+1}B$, so
$A$ and $B$ have the same finiteness locus, the same bounded sets, and
the same order of magnitude.  This already gives one implication in
\Cref{rem:ch49-open-question}: if $v_i\to v$ weakly and
$T_\omega(v_i)\to T_\omega(v)$, \emph{and} if in addition
$\limsup_iA(v_i)\le A(v)$, then
\Cref{thm:bj-linear-topology-complete}\,\cref{itm:ch49-top:criterion}
applies, because $A$ is lower semicontinuous
(\Cref{lem:ch49-dual-lsc}).  What is missing is precisely the
upper semicontinuity of $A$ along nets on which $B$ converges, and
comparability of two functions says nothing about the semicontinuity of
one along the level sets of the other.  For $n=1$ the two constants in
\eqref{eq:bj-dirac-energy-T} coincide, so $A=\tfrac12B$ and the question
has a positive answer; \Cref{ex:ch49-curve-case} makes this explicit.  For
$n\ge2$, \cref{thm:ch49-global-radon-riesz-counterexample} shows that
the missing upper semicontinuity cannot hold in the generality needed for
the proposed criterion.
\end{remark}

\section{Dual complexes and divisorial density}
\label{sec:ch49-density}

\subsection{Dual complexes inside the linear-growth locus}
\label{sec:ch49-dual-complexes}

For the SNC-test-configuration statements, assume additionally that
$X$ is smooth and $\operatorname{char}k=0$.
A regular total space with generic fibre $X$ requires smooth $X$ over
the algebraically closed field.  The approximation theorem below
extends to singular $X$ through a resolution, by Corollary~A.3 of
\cite{BJ22}; it does not require an SNC test configuration of singular $X$.
Let $\cX$ be an SNC test configuration
(\Cref{def:ch34-snc-tc}) with central fibre $\cX_0=\sum_{i\in I}b_iE_i$,
and let
\[
 \Delta_{\cX}=\bigcup_J\{\sigma^{-1}(w_\alpha)\mid\alpha\in\sigma_J\}
 \subset X^{\An},
 \qquad
 \sigma_J=\Bigl\{\alpha\in\mathbf R^J_{\ge0}\ \Bigm|\
 \sum_{j\in J}b_j\alpha_j=1\Bigr\},
\]
be its dual complex \eqref{eq:ch34-dual-complex}, with one face
for each connected component of a nonempty $E_J$.  The component
index is suppressed in $\sigma_J$ and $w_\alpha$ throughout this
discussion; it must be retained if $E_J$ is disconnected.
This is a compact subset of
$X^{\An}$ homeomorphic to the dual complex of $\cX_0$
(\Cref{lem:ch34-skeleton-in-gauss}), with vertices the divisorial points
$v_{E_i}=b_i^{-1}\ord_{E_i}$.  We write $w_\alpha$ for the point of
$\Delta_\cX$ attached to $\alpha$, suppressing $\sigma^{-1}$; on
$k(X)$ it is the restriction of the monomial valuation
\[
 w_\alpha(f)=\min\bigl\{\langle\alpha,\beta\rangle
 \ \bigm|\ \beta\in\mathbf Z^J_{\ge0},\ c_\beta\ne0\bigr\},
 \qquad
 \hat f=\sum_\beta c_\beta z^\beta,
\]
of \eqref{eq:quasimonomial-valuation} for regular functions, extended
to rational functions by quotients.  The expansion is taken in the
completed local ring of $\cX$ at the generic point of the stratum
$E_J$, with $z_j$ local equations of $E_j$.  Zero weights pass to
the corresponding face and its generized center.  Coordinate comparison
holds for regular functions at the center of the larger monomial
valuation, and hence for coherent ideals by
\Cref{thm:ch32-order-centers}; it is not obtained by subtracting
inequalities for arbitrary rational quotients.  Thus:
\begin{equation*}
 \alpha\le\alpha'\ \text{coordinatewise}
 \quad\Longrightarrow\quad
 w_\alpha\le w_{\alpha'}
 \ \text{in the order }\eqref{eq:semivaluation-order},
\end{equation*}
and $w_{t\alpha}=t\,w_\alpha$ for $t>0$, the normalization
$\sum b_j\alpha_j=1$ being irrelevant for either.

\begin{proposition}
\label{prop:ch49-dual-complex}
Let $X$ be smooth in characteristic zero and let $\cX$ be an SNC
test configuration.
\begin{enumerate}
\item\label{itm:ch49-dc:contained}
      $\Delta_\cX\subset X^{\lin}$, and
      $C_\cX:=\sup_{\Delta_\cX}T_\omega<\infty$.
\item\label{itm:ch49-dc:lipschitz}
      For every face $\sigma_J$ and every $c>0$, the map
      $\alpha\mapsto w_\alpha$ is Lipschitz from
      $\{\alpha\in\sigma_J\mid\alpha_j\ge c\ \text{for all }j\in J\}$
      with the sup-norm to $(X^{\lin},d_{\infty,\omega})$, with constant
      $C_\cX/c$.
\item\label{itm:ch49-dc:density}
      $X^{\divis}\cap\Delta_\cX$ is $d_{\infty,\omega}$-dense in
      $\Delta_\cX$.  More precisely, if
      $J_+:=\{j\mid\alpha_j>0\}$ and
      $c_\alpha:=\frac12\min_{j\in J_+}\alpha_j$, then for every
      $\varepsilon>0$ there is a rational point
      $\alpha'\in\sigma_{J_+}$ with
      \[
       \alpha'_j\ge c_\alpha\quad(j\in J_+),\qquad
       \|\alpha'-\alpha\|_\infty
       <\frac{\varepsilon c_\alpha}{\max\{C_\cX,1\}},
      \]
      such that $w_{\alpha'}\in X^{\divis}$ and
      $d_{\infty,\omega}(w_{\alpha'},w_\alpha)<\varepsilon$.
\item\label{itm:ch49-dc:compact}
      The restrictions of $\PSH(\omega)$ are equi-Lipschitz on every
      closed face for its Euclidean metric, with a constant depending on
      $\omega$ and $\cX$.  Hence the inclusion
      $\Delta_\cX\hookrightarrow(X^{\lin},d_{\infty,\omega})$
      is continuous, $\Delta_\cX$ is strongly compact, and its weak and
      strong topologies coincide.
\item\label{itm:ch49-dc:affine}
      Every $\psi\in\PL_{\mathbf R}(X)$ determined on $\cX$ restricts
      to an affine function on each face, with
      $\psi(w_\alpha)=\sum_{j\in J}a_j\alpha_j$ when
      $\psi=\phi_D$ and $D=\sum a_iE_i$.
\end{enumerate}
\end{proposition}

\begin{proof}
\cref{itm:ch49-dc:contained}  Fix a face $\sigma_J$ and set
$\alpha^0_j:=(\#J\cdot b_j)^{-1}$ for $j\in J$; then
$\alpha^0\in\sigma_J$ has positive rational coordinates, so $w_{\alpha^0}$
is divisorial by
\Cref{thm:ch32-quasimonomial-invariants}\,\cref{itm:ch32-qi:divisorial},
hence $T_\omega(w_{\alpha^0})<\infty$ by
\Cref{prop:ch49-psef-threshold}.  For $\alpha\in\sigma_J$ one has
$b_j\alpha_j\le1$, that is $\alpha_j\le b_j^{-1}=\#J\cdot\alpha^0_j$;
therefore $\alpha\le(\#J)\,\alpha^0$ coordinatewise and
$w_\alpha\le(\#J)\,w_{\alpha^0}$.  By
\Cref{prop:ch49-T-elementary}\,\cref{itm:ch49-T:order} and
\cref{itm:ch49-T:homog},
\[
 T_\omega(w_\alpha)\le\#J\cdot T_\omega(w_{\alpha^0})<\infty .
\]
Taking the maximum over the finitely many faces bounds $T_\omega$ on
$\Delta_\cX$.

\cref{itm:ch49-dc:lipschitz}  Let $\alpha,\alpha'$ lie in the indicated
set and put $\eta:=\|\alpha-\alpha'\|_\infty$.  Then
$\alpha\le\alpha'+\eta\mathbf 1\le(1+\eta/c)\alpha'$ and symmetrically,
so $\lambda^{-1}\alpha'\le\alpha\le\lambda\alpha'$ with
$\lambda=1+\eta/c$.  Hence
$\lambda^{-1}w_{\alpha'}\le w_\alpha\le\lambda w_{\alpha'}$, and
\Cref{lem:ch49-multiplicative-comparison} gives
\[
 d_{\infty,\omega}(w_\alpha,w_{\alpha'})
 \le(\lambda-1)\max\{T_\omega(w_\alpha),T_\omega(w_{\alpha'})\}
 \le\frac{C_\cX}{c}\,\|\alpha-\alpha'\|_\infty .
\]

\cref{itm:ch49-dc:density}  The rational points of the relative simplex
$\sigma_{J_+}$, whose defining equation has rational coefficients, are
dense.  Choose one satisfying the two displayed bounds in the statement;
for a sufficiently close choice all its $J_+$-coordinates are at least
$c_\alpha$.  Applying \cref{itm:ch49-dc:lipschitz} on this relative face
gives
\[
 d_{\infty,\omega}(w_{\alpha'},w_\alpha)
 \le\frac{C_\cX}{c_\alpha}\|\alpha'-\alpha\|_\infty
 <\varepsilon.
\]
Since $\alpha'$ is already a normalized rational point of the same face,
$w_{\alpha'}\in X^{\divis}\cap\Delta_\cX$ by
\Cref{thm:ch32-quasimonomial-invariants}\,\cref{itm:ch32-qi:divisorial}.

\cref{itm:ch49-dc:compact}  This is the exact content of
\cite[Theorem~A.4]{BJ22}.  Its proof passes to $k((\varpi))$ through
the Gauss section, uses the uniform facewise Lipschitz estimate
\cite[Theorem~6.1]{BFJ16} for psh model functions, and passes to
decreasing fixed-class limits.  Taking the supremum of
$|\phi(w_\alpha)-\phi(w_{\alpha'})|$ preserves that Lipschitz bound
for $d_{\infty,\omega}$ on each face.  The finitely many closed faces
give a continuous inclusion of the whole dual complex.
Compactness and the Hausdorff property make it a homeomorphism onto
its image.  This global boundary conclusion is quoted, rather than
deduced from the interior estimate with $C_\cX/c$.

\cref{itm:ch49-dc:affine}  This is
\Cref{prop:ch34-affine-on-faces}, restated.
\end{proof}

\begin{remark}[Interior estimates and the boundary theorem]
\label{rem:ch49-dual-complex-lipschitz}
The elementary estimate with $C_\cX/c$ degenerates as $c\downarrow0$.
It suffices to approximate a point inside its minimal face by rational
weights.  The separate published theorem
\cref{itm:ch49-dc:compact} controls the closed faces as well.
Thus the degeneration of that elementary bound does not represent a
failure of continuity or compactness of an SNC dual complex.
For toric valuations the closed formula
\eqref{eq:ch49-toric-dinfty} also gives an explicit boundary estimate.
\end{remark}

\subsection{Retractions decrease valuations}
\label{sec:ch49-retractions}

\begin{lemma}
\label{lem:ch49-retraction-decreases}
Let $X$ be smooth in characteristic zero, let $\cX$ be an SNC test
configuration, and let $r_\cX:X^{\An}\to\Delta_\cX$
its retraction \eqref{eq:ch34-retraction}.  Then for every
$v\in X^{\An}$:
\begin{enumerate}
\item\label{itm:ch49-retr:order}
      $r_\cX(v)\le v$ in the order \eqref{eq:semivaluation-order};
\item\label{itm:ch49-retr:T}
      $T_\omega(r_\cX(v))\le T_\omega(v)$; in particular
      $r_\cX(X^{\An})\subset X^{\lin}$, consistently with
      \Cref{prop:ch49-dual-complex}\,\cref{itm:ch49-dc:contained};
\item\label{itm:ch49-retr:limit}
      the net $(r_\cX(v))_\cX$, indexed by the directed set of SNC test
      configurations, is nondecreasing and converges to $v$ in
      $X^{\An}$.
\end{enumerate}
\end{lemma}

\begin{proof}
\cref{itm:ch49-retr:order}  Let $\mathfrak a$ be a coherent ideal on
$X$.  By \Cref{def:ch34-snc-tc}, $\cX$ dominates the constant model;
let $\mathfrak b$ be the pullback of $\mathfrak a$ to $\cX$.
For every integer $m\ge1$, the ideal
$\mathfrak b+(\varpi^m)$ has zero locus in $\cX_0$.
Applying \Cref{prop:ch34-retraction}\,\cref{itm:ch34-retr:ideal}
and the constant-model ideal identity of
\Cref{thm:ch32-gauss-transfer}\,\cref{itm:ch32-gt:ideals} gives
\[
 \min\{r_\cX(v)(\mathfrak a),m\}
 \le \min\{v(\mathfrak a),m\}.
\]
Both Gauss valuations have $\varpi$-value one.  Letting $m$ increase
gives $r_\cX(v)(\mathfrak a)\le v(\mathfrak a)$, including infinite
ideal values.  Since $\mathfrak a$ was arbitrary, this is
\eqref{eq:semivaluation-order}.

\cref{itm:ch49-retr:T}  Immediate from \cref{itm:ch49-retr:order} and
\Cref{prop:ch49-T-elementary}\,\cref{itm:ch49-T:order}.

\cref{itm:ch49-retr:limit}  For $\cX'\ge\cX$,
$r_\cX\circ r_{\cX'}=r_\cX$; applying
\cref{itm:ch49-retr:order} to $r_{\cX'}(v)$ proves monotonicity.
Convergence is \cite[Theorem~A.1]{BJ22}, at its smooth,
characteristic-zero signature.  Its proof uses PL functions factoring
through determining SNC test configurations and separating points.
It does not assume equivariant models cofinal among all
nonequivariant SNC models.
\end{proof}

\begin{lemma}
\label{lem:ch49-dual-lsc}
The function
\[
 X^{\An}\longrightarrow[0,+\infty],
 \qquad
 v\longmapsto E^\vee_\omega(\delta_v)
 =\sup_{\phi\in\cE^1(\omega)}\bigl(E_\omega(\phi)-\phi(v)\bigr),
\]
is lower semicontinuous, and it is nondecreasing for the order
\eqref{eq:semivaluation-order}.
\end{lemma}

\begin{proof}
Each $\phi\in\cE^1(\omega)$ is upper semicontinuous with values in
$[-\infty,+\infty)$, so $v\mapsto E_\omega(\phi)-\phi(v)$ is lower
semicontinuous with values in $(-\infty,+\infty]$; a supremum of lower
semicontinuous functions is lower semicontinuous.  For monotonicity, if
$v\le w$ then $\phi(v)\ge\phi(w)$ for every $\omega$-psh $\phi$,
hence $E_\omega(\phi)-\phi(v)\le E_\omega(\phi)-\phi(w)$, and the
suprema compare in the same way.
\end{proof}

\subsection{The density theorem}
\label{sec:ch49-density-theorem}

\begin{theorem}\label{thm:bj-strong-divisorial-density}
If $\operatorname{char}k=0$, then $X^{\divis}$ is dense in
$X^{\lin}$ for $d_\infty$.
\end{theorem}

\begin{proof}
The theorem is \cite[Theorem~11.18]{BJ22}.  The detailed argument
below uses approximation from below in Corollary~A.3, lower
semicontinuity and monotonicity of the dual energy, and the already
proved strong convergence criterion.  The resolution route in that
corollary also covers singular $X$.
\end{proof}

We spell the argument out, because the three ingredients are exactly the
three statements just proved and because the same pattern is used again in
\Cref{sec:ch49-toric}.

\begin{proof}[Proof of \Cref{thm:bj-strong-divisorial-density}, in detail]
Let $v\in X^{\lin}$ and $\varepsilon>0$.

\emph{Step 1: an approximating net from below.}  By
\Cref{imp:ch49-imported}\,\cref{itm:ch49-import:approximation} there is a
net $(v_i)$ in $X^{\divis}$ with $v_i\le v$ for all $i$ and
$v_i\to v$ in $X^{\An}$.  If $X$ is smooth, an alternative with quasimonomial rather than
divisorial members is $v_i:=r_{\cX_i}(v)$ for the net of
SNC test configurations, by
\Cref{lem:ch49-retraction-decreases}; the retraction route is valid at that smooth signature.
For singular $X$ we use Corollary~A.3, which first lifts to a resolution
and then projects the divisorial approximants back to $X$.
To keep the rational-scaling convention for $X^{\divis}$, take the
source's intermediate factors $t_0>1$ rational and tending to one.
Its divisorial approximants satisfy $v_0\le t_0v$; rescaling to
$v_0/t_0\le v$ preserves rational divisorial scaling and convergence.
No monotonicity of the divisorial net is needed.

\emph{Step 2: the energies converge.}  By
\Cref{lem:ch49-dual-lsc} the function $w\mapsto E^\vee_\omega(\delta_w)$
is nondecreasing for the order, so
$E^\vee_\omega(\delta_{v_i})\le E^\vee_\omega(\delta_v)$ for all $i$;
and it is lower semicontinuous, so
$\liminf_iE^\vee_\omega(\delta_{v_i})\ge E^\vee_\omega(\delta_v)$.  Hence
$E^\vee_\omega(\delta_{v_i})\to E^\vee_\omega(\delta_v)$, a finite limit
by \Cref{thm:bj-dirac-finite-energy}.

\emph{Step 3: strong convergence.}  Steps~1 and~2 verify the two
conditions of
\Cref{thm:bj-linear-topology-complete}\,\cref{itm:ch49-top:criterion},
whence $d_{\infty,\omega}(v_i,v)\to0$.  Choose $i$ with
$d_{\infty,\omega}(v_i,v)\le\varepsilon/2$.

\emph{Step 4: from quasimonomial to divisorial, if necessary.}  If the net
of Step~1 was taken to be $v_i=r_{\cX_i}(v)\in\Delta_{\cX_i}$, apply
\Cref{prop:ch49-dual-complex}\,\cref{itm:ch49-dc:density} to obtain
$w\in X^{\divis}\cap\Delta_{\cX_i}$ with
$d_{\infty,\omega}(w,v_i)\le\varepsilon/2$.  Then
$d_{\infty,\omega}(w,v)\le\varepsilon$.
\end{proof}

The characteristic-zero hypothesis belongs to the cited approximation
theorem, not to divisorial weak density in the whole analytification.  It
must not be erased by appealing to a dual complex that need not exist in
positive characteristic.

\begin{warning}{Where characteristic zero enters, and where it does not}
\label{warn:ch49-characteristic}
Weak density of $X^{\divis}$ in $X^{\An}$ is
\Cref{thm:triv-divisorial-density} and holds over any algebraically closed
trivially valued field: it uses only that rational PL functions form a
point-separating lattice.  What needs resolution of singularities is
\emph{approximation from below}, that is, the production of divisorial
$v_i\le v$ converging to $v$; this is
\Cref{imp:ch49-imported}\,\cref{itm:ch49-import:approximation}, and it requires characteristic zero.  The SNC statements above
also have their stated smooth, characteristic-zero signature.  In
positive characteristic the following remain valid without change:
\Cref{prop:ch49-T-equals-T},
\Cref{thm:bj-dirac-finite-energy},
\Cref{prop:ch49-dinfty-basic},
\Cref{lem:ch49-multiplicative-comparison},
\Cref{thm:bj-linear-topology-complete}\,\cref{itm:ch49-top:embedding} and
\cref{itm:ch49-top:criterion}, and the whole of
\Cref{sec:ch49-toric,sec:ch49-toric-ma}, since a toric variety carries its
own combinatorial resolutions.  The pointwise psh topology and completeness statements also hold in
every characteristic, by the proofs given here.  Strong divisorial density
is quoted in characteristic zero.  The explicit toric formulas do not
need a general resolution theorem.
\end{warning}

\begin{corollary}
\label{cor:ch49-separable}
Assume $\operatorname{char}k=0$.  Then
$(X^{\lin},d_{\infty,\omega})$ is a complete metric space in which the
divisorial points are dense, and on which the scaling action
$t\cdot$ is a one-parameter group of homotheties of ratio $t$ fixing
$v_{\triv}$.  If moreover $k$ is countable, then
$(X^{\lin},d_{\infty,\omega})$ is separable, hence a Polish space.
\end{corollary}

\begin{proof}
Completeness is
\Cref{thm:bj-linear-topology-complete}\,\cref{itm:ch49-top:complete} and
density is \Cref{thm:bj-strong-divisorial-density}; the homothety property
is \Cref{prop:ch49-dinfty-basic}\,\cref{itm:ch49-dinf:homog} together with
$T_\omega(v_{\triv})=0$.  For separability, note that
$X^{\divis}$ is contained in the countable union of rays
$\mathbf R_{>0}\ord_E$ over
the prime divisors $E$ of the countably many normal projective
birational models of $X$ --- countably many because $k$ is countable
and $X$ is of finite type --- and each ray is separable: for $s,t\ge0$,
$d_{\infty,\omega}(s\ord_E,t\ord_E)=|s-t|T_\omega(\ord_E)$ by the
equality case of \Cref{lem:ch49-multiplicative-comparison}.
Rational parameters on the countably many rays give a countable
dense subset.  This is compatible with the book's rational-scaling
convention for $X^{\divis}$.
\end{proof}

\subsection{Bounded-set continuity and growth}
\label{sec:ch49-bounded}

On an SNC dual complex at the signature above, the entire psh family
is equi-Lipschitz by \Cref{prop:ch49-dual-complex}\,\cref{itm:ch49-dc:compact}.
A general $T$-ball has a weaker property: every individual psh function
is weakly continuous there.  The distinction is essential in the
counterexamples below.

\begin{proposition}\label{prop:bj-psh-continuous-bounded-valuations}
Every $\phi\in\PSH(\omega)$ is weakly continuous on
\[
                         \{v\in X^{\lin}\mid T_\omega(v)\leq C\}.
\tag{49.16}\label{eq:bj-T-ball}
\]
Every $\phi\in\mathcal E^1$ further satisfies
\[
 |\phi(v)|\leq A\max\{1,T_\omega(v)\}^{\,1-2^{-n}}+B
\tag{49.17}\label{eq:bj-e1-sublinear-growth}
\]
for constants controlled by $J(\phi)$ and $\sup\phi$.
\end{proposition}

\begin{proof}
These are \cite[Theorems~11.11--11.12]{BJ22}; we detail how they
follow from the earlier quantitative interfaces.
For bounded-set continuity, normalize $\sup\phi=0$.
The convex-composition rule \cite[Corollary~4.11]{BJ22} makes
$\exp(\phi)$ bounded psh and hence finite-energy.
On $\{T_\omega\le C\}$ it is bounded below by $e^{-C}$,
so continuity of its logarithm recovers continuity of $\phi$.
It therefore suffices to treat a finite-energy potential $\eta$.

Choose decreasing continuous fixed-class approximants
$\eta_j\downarrow\eta$ on a tail with bounded J-values.
Decreasing strong convergence gives $I(\eta_j,\eta)\to0$.
The uniform pairing estimate \cite[Theorem~10.3]{BJ22}, applied to
$\delta_v-\delta_{v_{\triv}}$, is uniform for $T_\omega(v)\le C$
because \eqref{eq:bj-dirac-energy-T} bounds the dual energies of
these Dirac masses.  Together with
$\eta_j(v_{\triv})\to\eta(v_{\triv})$, it yields
$\eta_j\to\eta$ uniformly on that ball.
A uniform limit of its continuous restrictions is continuous.

For the growth statement, apply the same pairing estimate to
$\phi,0$ and $\delta_v,\delta_{v_{\triv}}$.
The I/J comparison bounds its potential terms by $J(\phi)$,
and \eqref{eq:bj-dirac-energy-T} bounds its measure terms by
$\max\{1,T_\omega(v)\}$.  In Theorem~10.3 the potential exponent
is $\alpha=2^{-n}$.  The measure factors contribute exponents
$1/2$ and $1/2-\alpha$, whose sum is $1-2^{-n}$, giving
\[
 0\le\sup\phi-\phi(v)\le
 A_{J(\phi)}\max\{1,T_\omega(v)\}^{1-2^{-n}}.
\]
Adding $|\sup\phi|$ proves \eqref{eq:bj-e1-sublinear-growth}.
\end{proof}

\begin{remark}[The role of the two statements]
\label{rem:ch49-two-quantitative}
The proof of \Cref{prop:bj-psh-continuous-bounded-valuations}
identifies the exact quantitative input behind the published
interfaces of \cref{itm:ch49-import:equicontinuity}.
Their consequences are different.

\eqref{eq:bj-T-ball} is a statement about a \emph{fixed} $\phi$: it
upgrades upper semicontinuity to continuity on each $T$-ball.  It is used
once, in the completeness proof
\Cref{thm:bj-linear-topology-complete}\,\cref{itm:ch49-top:complete}, to
pass to the limit in the inequality $|\phi(v_i)-\phi(v_j)|\le\varepsilon$
in the second variable.  It is emphatically not equicontinuity of the
family $\PSH(\omega)$ on the ball --- that is false, by
\Cref{ex:ch49-weak-not-strong} --- and it is not weak-to-strong continuity
of the inclusion of the ball.

\eqref{eq:bj-e1-sublinear-growth} is a statement about a \emph{fixed}
$\phi$: it bounds its values at every $v\in X^{\lin}$ by a
sublinear function of $T_\omega(v)$.  It is sharper than the elementary finiteness bound
$|\phi(v)|\le|\sup\phi|+T_\omega(v)$, which is linear in $T_\omega(v)$
and follows at once from \eqref{eq:bj-T-valuation}.  Every use in this book
of \eqref{eq:bj-e1-sublinear-growth} that only needs finiteness can be
served by that trivial bound; the exponent $1-2^{-n}$ is needed only when
a growth rate is asserted.
\end{remark}

\begin{example}[Weak convergence without strong convergence, on a
$T$-ball]
\label{ex:ch49-weak-not-strong}
Take $X=\mathbf P^1_k$, whose set of closed points is infinite, with
$\omega=c_1(\cO(d))$, $d\ge1$.  Index a net by the finite subsets
$S\subset X(k)$, ordered by inclusion, and choose a closed point
$p_S\notin S$.  Put $v_S:=\ord_{p_S}\in X^{\divis}$.
\begin{enumerate}
\item\label{itm:ch49-wns:ball}
      By \Cref{ex:ch49-curve-case}, $T_\omega(v_S)=d$ for every $S$, so
      the whole net lies in the compact ball $\{T_\omega\le d\}$.
\item\label{itm:ch49-wns:weak}
      $v_S\to v_{\triv}$ weakly.  Indeed for $f\in k(X)^\times$ the set
      $Z_f$ of zeros and poles of $f$ is finite, and $v_S(f)=0$ as
      soon as $S\supset Z_f$; since the functions $v\mapsto v(f)$
      generate the weak topology on the valuation locus, this is weak convergence.
\item\label{itm:ch49-wns:strong}
      $d_{\infty,\omega}(v_S,v_{\triv})=T_\omega(v_S)=d$ for every $S$,
      by \Cref{prop:ch49-dinfty-basic}\,\cref{itm:ch49-dinf:vertex}.  So the
      net does not converge strongly to $v_{\triv}$, and it cannot converge strongly to any other point, since strong
      convergence implies weak convergence and the latter limit is unique.
\item\label{itm:ch49-wns:energy}
      The obstruction is visible in the energies.
      In dimension one \eqref{eq:bj-dirac-energy-T} gives
      \[
       E^\vee_\omega(\delta_{v_S})=\frac d2
       \qquad\text{while}\qquad
       E^\vee_\omega(\delta_{v_{\triv}})=0 ,
      \]
      so the second condition of
      \Cref{thm:bj-linear-topology-complete}\,\cref{itm:ch49-top:criterion}
      fails by exactly $d/2$.  For $n=1$ both constants in
      \eqref{eq:bj-dirac-energy-T} equal $\tfrac12$, and indeed
      $\tfrac12T_\omega(v_S)=d/2$.
\end{enumerate}
This is the valuative shadow of \Cref{ex:ch46-energy-not-weakly-continuous}
and of \Cref{ex:ch48-not-a-metric}: the same failure of weak continuity of
energy, seen on the base rather than on the potentials.
\end{example}

\begin{theorem}
\label{thm:ch49-global-radon-riesz-counterexample}
Let $k$ be an algebraically closed field with the trivial absolute value,
let $n\ge2$, and put
\[
             X=\mathbf P^n_k,
             \qquad L=\cO_{\mathbf P^n}(1),
             \qquad \omega=c_1(L).
\]
Fix a linear line $C\simeq\mathbf P^1_k$ in $X$.  For $n\ge3$,
let $v_C=\ord_{E_C}$, where $E_C$ is the exceptional divisor of the
blowup of $X$ along $C$; for $n=2$, set $v_C=\ord_C$.  Choose
pairwise distinct points $p_j\in C(k)$, and let $v_j=\ord_{E_j}$, where
$E_j$ is the exceptional divisor of the blowup of $X$ at $p_j$.  Then
\[
              v_j\longrightarrow v_C
              \quad\text{in the weak topology of }X^{\lin},
\]
and
\[
              T_\omega(v_j)=T_\omega(v_C)=1
              \qquad\text{for every }j,
\]
whereas
\[
              d_{\infty,\omega}(v_j,v_C)\ge1
              \qquad\text{for every }j.
\]
Consequently, weak convergence together with convergence of the
Boucksom--Jonsson exhaustion does not imply strong convergence on
$X^{\lin}$.  This failure occurs on $\mathbf P^n_k$ for every
$n\ge2$, in arbitrary characteristic.
\end{theorem}

\begin{proof}
All valuations in the statement are divisorial and hence have linear growth
by \cite[Proposition~1.28]{BJ22}.

We first prove weak convergence.  Denote by $\eta$ the generic point of
$C$, and by $\mathcal I$ the ideal sheaf of $C$.  We use the following
elementary specialization fact.  If $a$ is a nonzero regular function on
an open neighbourhood of $\eta$, put
\[
 r_a=\ord_{\mathcal I_\eta}(a).
\]
For $n\ge3$, the blowup construction identifies this number with
$\ord_{E_C}(a)$; for $n=2$, it is $\ord_C(a)$.  After shrinking
around $\eta$, the function $a$ lies in $\mathcal I^{r_a}$, and its
initial class
\[
 \operatorname{in}_C(a)
 \in \mathcal I^{r_a}/\mathcal I^{r_a+1}
 \simeq \operatorname{Sym}^{r_a}(\mathcal I/\mathcal I^2)
\tag{49.17a}\label{eq:ch49-pn-initial-class}
\]
is a regular section that is nonzero at $\eta$.  Its zero locus on the
curve $C$ is therefore finite, after also adding the finite complement of
the chosen neighbourhood.

At every remaining point $p\in C(k)$, the class in
\eqref{eq:ch49-pn-initial-class} has nonzero fibre.  Since $C$ is a smooth
linear subspace of the smooth variety $X$, the conormal injection
\[
 (\mathcal I/\mathcal I^2)\otimes k(p)
 \lhook\joinrel\longrightarrow
 \mathfrak m_p/\mathfrak m_p^2
\]
identifies the source with the transverse linear forms.  Its $r_a$-th
symmetric power is injective.  Thus the image of
$\operatorname{in}_C(a)(p)$ in
$\mathfrak m_p^{r_a}/\mathfrak m_p^{r_a+1}$ is nonzero.  Since
$a\in\mathfrak m_p^{r_a}$, this proves
\[
             \ord_{\mathfrak m_p}(a)=r_a
\tag{49.17b}\label{eq:ch49-pn-regular-specialization}
\]
outside a finite subset of $C(k)$.

Now fix $f\in k(X)^\times$.  On an affine open neighbourhood of $\eta$,
write $f=a/b$ with nonzero regular functions $a,b$.  Apply
\eqref{eq:ch49-pn-regular-specialization} to both functions and discard the
union of their finite bad sets.  The exceptional valuation of the blowup at
a smooth point $p$ is the $\mathfrak m_p$-adic order: the associated
graded ring is a polynomial ring over $k(p)$, so this order is additive
and extends to the fraction field.  Hence, at every remaining point,
\[
 \begin{aligned}
  \ord_{E_p}(f)
    &=\ord_{\mathfrak m_p}(a)-\ord_{\mathfrak m_p}(b)\\
    &=r_a-r_b=v_C(f).
 \end{aligned}
\tag{49.17c}\label{eq:ch49-pn-rational-specialization}
\]
Pairwise distinctness of the $p_j$ makes the sequence eventually avoid
every finite subset of $C(k)$.  Equation
\eqref{eq:ch49-pn-rational-specialization} therefore gives
$v_j(f)=v_C(f)$ for all sufficiently large $j$.  The weak topology on
the valuation locus is pointwise convergence on $k(X)^\times$, so
$v_j\to v_C$ weakly.

We next compute the exhaustion.  A nonzero homogeneous polynomial of degree
$m$ has multiplicity at a point at most $m$, and equality is attained by
the $m$-th power of a hyperplane through that point.  The definition of
$T_\omega$ and \cite[Lemma~4.46]{BJ22} thus give
\[
                              T_\omega(v_j)=1.
\tag{49.17d}\label{eq:ch49-pn-T-point}
\]
The homogeneous ideal of the linear line $C$ is generated by linear
forms.  A nonzero degree-$m$ homogeneous polynomial cannot lie in
$\mathcal I_C^{m+1}$, while the $m$-th power of any hyperplane
containing $C$ lies in $\mathcal I_C^m$.  It follows, both for
$\ord_{E_C}$ when $n\ge3$ and for $\ord_C$ when $n=2$, that
\[
                              T_\omega(v_C)=1.
\tag{49.17e}\label{eq:ch49-pn-T-center}
\]

Finally, choose a hyperplane $H_j$ through $p_j$ that does not contain
$C$, and let $\ell_j\in H^0(X,L)$ be a defining section.  The hyperplane
has multiplicity one at $p_j$, whereas $\ell_j$ is nonzero at the
generic point of $C$.  Therefore
\[
                          v_j(\ell_j)=1,
                          \qquad v_C(\ell_j)=0.
\tag{49.17f}\label{eq:ch49-pn-moving-hyperplane}
\]
The section formula \eqref{eq:bj-dinfty-sections} gives
\[
 d_{\infty,\omega}(v_j,v_C)
 \ge \lvert v_j(\ell_j)-v_C(\ell_j)\rvert=1.
\tag{49.17g}\label{eq:ch49-pn-strong-gap}
\]
Thus the sequence is not strongly convergent although it is weakly
convergent and its exhaustion values agree identically with the exhaustion
of its limit.  This is a negative answer to the implication asked for in
\cite[Question~11.16]{BJ22}, in every dimension $n\ge2$.
\end{proof}

\begin{remark}[Why the earlier curve example did not answer the question]
\label{rem:ch49-curve-versus-radon-riesz}
In \cref{ex:ch49-weak-not-strong}, the weak limit is $v_{\triv}$, so
$T_\omega(v_S)=d$ but $T_\omega(v_{\triv})=0$.  That example proves
that weak and strong convergence differ on a $T$-ball, but it does not
satisfy the exhaustion-convergence hypothesis of Question~11.16.  In
\cref{thm:ch49-global-radon-riesz-counterexample}, both the moving point
valuations and their line-centred limit have exhaustion exactly one; this is
the additional equality that gives the genuine counterexample.
\end{remark}

\section{The toric dictionary and its signs}
\label{sec:ch49-toric}

\subsection{Invariant valuations and the toric retraction}
\label{sec:ch49-toric-setup}

Let $T\simeq\mathbf G_m^n$, with character lattice $M$, cocharacter
lattice $N$, and pairing $\langle\,\cdot\,,\,\cdot\,\rangle$.
A $T(k)$-invariant valuation is determined by
\[
                       v(\chi^m)=\langle v,m\rangle,
\tag{49.18}\label{eq:bj-toric-valuation-pairing}
\]
so $N_{\mathbf R}$ embeds in $X^{\val}$.  There is a retraction
$\rho:X^{\val}\to N_{\mathbf R}$.

Let $X=X_\Sigma$ be proper and normal toric, and let the BJ polytope
$P_{\mathrm{BJ}}\subset M_{\mathbf R}$ represent $\omega$.  Put
\[
 f_P(v):=\sup_{m\in P_{\mathrm{BJ}}}\langle v,m\rangle.
\tag{49.19}\label{eq:bj-toric-support}
\]

We now make all of this precise, since the rest of the chapter computes
with it.  Fix a full-dimensional lattice polytope $P\subset M_{\mathbf R}$
whose inner normal fan is $\Sigma$, let $L=L_P$ be the corresponding
ample toric line bundle with a chosen rational toric section $s_P$, and set
$P_{\mathrm{BJ}}:=-P$.  For $u\in N_{\mathbf R}$ let $v_u$ be the
monomial valuation of $k(X)=k(M)$ determined by
\eqref{eq:bj-toric-valuation-pairing}, that is
\[
 v_u\Bigl(\sum_{m\in M}c_m\chi^m\Bigr)
 =\min\{\langle u,m\rangle\mid c_m\ne0\}.
\]

\begin{lemma}
\label{lem:ch49-toric-slice}
\begin{enumerate}
\item\label{itm:ch49-ts:val}
      For every $u\in N_{\mathbf R}$, $v_u$ is a valuation of
      $k(X)^\times$, trivial on $k^\times$, with a center on $X$; and
      $u\mapsto v_u$ is a homeomorphism of $N_{\mathbf R}$ onto the set
      of $T(k)$-invariant points of $X^{\val}$.  We write
      $\iota:N_{\mathbf R}\hookrightarrow X^{\val}$ for this embedding.
\item\label{itm:ch49-ts:retraction}
      The map $\rho:X^{\val}\to N_{\mathbf R}$ sending $v$ to the
      linear form $m\mapsto v(\chi^m)$ is continuous, satisfies
      $\rho\circ\iota=\mathrm{id}$, and
      $\iota(\rho(v))\le v$ in the order \eqref{eq:semivaluation-order}.
\item\label{itm:ch49-ts:section}
      For every $u\in N_{\mathbf R}$,
      \[
       v_u(s_P)=f_P(u),
       \qquad\text{equivalently}\qquad
       \log|s_P|(v_u)=-f_P(u),
      \tag{49.20}\label{eq:ch49-toric-section}
      \]
      and, for $m\ge1$ and
      $s=\sum_{m'\in mP\cap M}c_{m'}\chi^{m'}s_P^{\,m}\in H^0(X,mL)$
      nonzero,
      \[
       \frac{v_u(s)}{m}
       =\min_{p\in\Lambda(s)}\langle u,p\rangle+f_P(u),
       \qquad
       \Lambda(s):=\{m'/m\mid c_{m'}\ne0\}\subset P .
      \tag{49.21}\label{eq:ch49-toric-section-value}
      \]
\end{enumerate}
\end{lemma}

\begin{proof}
\cref{itm:ch49-ts:val}  Multiplicativity of $v_u$ is the fact that the
group algebra $k[M]$ is a domain graded by $M$: if
$\Lambda,\Lambda'$ are the sets of exponents where the minimum of
$\langle u,\cdot\rangle$ is attained for two Laurent polynomials, the
product of the corresponding initial forms is a nonzero element of
$k[M]$, so the minimum for the product is the sum of the two minima.  The
ultrametric inequality is clear, as is triviality on $k^\times$.  A
center exists because $X$ is proper.  Invariance under $T(k)$ holds
because $t\cdot\chi^m=\chi^m(t)\chi^m$ and $v_u$ is trivial on
$k^\times$; conversely an invariant valuation is determined by its values
on characters, since the invariant part of the filtration is spanned by
monomials.  Continuity in both directions follows because the weak topology
on the valuation locus is generated by the functions $v\mapsto v(f)$,
$f\in k(X)^\times$.

\cref{itm:ch49-ts:retraction}  The form $m\mapsto v(\chi^m)$ is additive
because $v$ is multiplicative, so $\rho(v)\in N_{\mathbf R}$;
continuity and $\rho\circ\iota=\mathrm{id}$ are immediate.  For the
inequality, $v(\sum c_m\chi^m)\ge\min_mv(\chi^m)
=v_{\rho(v)}(\sum c_m\chi^m)$ by the ultrametric inequality.  To obtain the valuation-order statement, choose a toric chart
$U_\sigma$ containing the center of $v$; the values on its monomial
generators show $\rho(v)\in\sigma$, so $v_{\rho(v)}$ also has a
center in this chart.  The inequality holds for every element of
$k[\sigma^\vee\cap M]$.  For a function regular at $c(v)$, write it
as $a/b$ with $b$ a unit there.  Then $v(b)=0$, and the inequality
and nonnegativity force $v_{\rho(v)}(b)=0$ as well.  Subtraction of
these zero denominator values preserves the inequality.
An arbitrary quotient of two inequalities would not suffice.

\cref{itm:ch49-ts:section}  Let $\sigma\in\Sigma$ be a maximal cone
containing $u$ and let $m_\sigma\in P$ be the corresponding vertex,
characterized by
$\langle u',m_\sigma\rangle=\min_{p\in P}\langle u',p\rangle$ for
$u'\in\sigma$.  On the affine chart $U_\sigma$, the line bundle $L$
is trivialized by the frame $\tau_\sigma:=\chi^{m_\sigma}s_P$, and the
center of $v_u$ lies in $U_\sigma$
\cite[Theorem~4.2.8 and \S6.1]{CLS11}.  Hence
\[
 v_u(s_P)=v_u\bigl(\chi^{-m_\sigma}\tau_\sigma\bigr)
 =v_u(\chi^{-m_\sigma})=-\langle u,m_\sigma\rangle,
\]
and $-\langle u,m_\sigma\rangle=-\min_{p\in P}\langle u,p\rangle
=\sup_{m\in-P}\langle u,m\rangle=f_P(u)$, that is
\[
 v_u(s_P)=f_P(u),
\]
which is \eqref{eq:ch49-toric-section}.

For the second formula, use the frame $\tau_\sigma^{\,m}$ of $mL$ and
write
\[
 s=\left(\sum_{m'}c_{m'}\chi^{m'-mm_\sigma}\right)\tau_\sigma^{\,m}.
\]
All
exponents $m'-mm_\sigma$ lie in the semigroup
$\sigma^\vee\cap M$.  Therefore
\[
 v_u(s)=\min_{c_{m'}\ne0}\langle u,m'-mm_\sigma\rangle
 =\min_{c_{m'}\ne0}\langle u,m'\rangle+m\,f_P(u).
\]
\end{proof}

\begin{warning}{Reflection relative to \Cref{chap:model-calculations}}
\label{warn:ch49-toric-signs}
\Cref{chap:model-calculations} fixes metric weights as $-\log\|s\|$ and denotes their recession
function by $\check h_P(v)=h_P(-v)$.  With the same toric section, the
polytope in BJ Appendix~B is $P_{\mathrm{BJ}}=-P$.  Hence
\[
                         f_{P_{\mathrm{BJ}}}=\check h_P.
\tag{49.22}\label{eq:bj-toric-sign-conversion}
\]
This explicit reflection prevents the two correct conventions from being
mixed.
\end{warning}

Indeed, by \eqref{eq:appJ-toric-recession},
$\check h_P(u)=h_P(-u)=\sup_{m\in P}\langle-u,m\rangle
=\sup_{m\in-P}\langle u,m\rangle=f_{P_{\mathrm{BJ}}}(u)$, which is
\eqref{eq:bj-toric-sign-conversion}; and
\eqref{eq:ch49-toric-section} says that $f_P$ is the additive weight of
the trivial metric on $L$ read on the toric slice, exactly as
$\check h_P$ is in \Cref{sec:toric-convex}.  The unit test is
$P=[0,d]$ on $\mathbf P^1$: both conventions give
$f_P(u)=\check h_P(u)=\max\{0,-du\}$, as in
\eqref{eq:p1-standard-weight}.

\subsection{Linear growth and \texorpdfstring{$d_\infty$}{d infinity} on
the toric slice}
\label{sec:ch49-toric-T}

Before the dictionary itself we record the two invariants of
\Cref{sec:ch49-growth,sec:ch49-topologies} on $N_{\mathbf R}$.  They are
computed exactly, with no limit and no approximation, from
\Cref{lem:ch49-toric-slice}\,\cref{itm:ch49-ts:section}.

\begin{proposition}
\label{prop:ch49-toric-T}
For every $u\in N_{\mathbf R}$ and every $m\ge1$,
\[
 \frac{a_m(v_u)}{m}=f_P(u)+f_P(-u),
 \qquad\text{hence}\qquad
 T_\omega(v_u)=\max_{p,q\in P}\langle u,p-q\rangle .
\tag{49.23}\label{eq:ch49-toric-T}
\]
In particular $N_{\mathbf R}\subset X^{\lin}$, the function
$T_\omega\circ\iota$ is the width function of the polytope $P$, it is
convex, positively homogeneous of degree one, and it vanishes exactly at
$u=0$, where $v_0=v_{\triv}$.
\end{proposition}

\begin{proof}
By \eqref{eq:ch49-toric-section-value}, $m^{-1}v_u(s)$ equals
$\min_{p\in\Lambda(s)}\langle u,p\rangle+f_P(u)$; the maximum over
nonzero $s\in H^0(X,mL)$ is attained when $\Lambda(s)$ is a single
lattice point, that is when $s$ is a monomial, and equals
\[
 \max\Bigl\{\langle u,p\rangle\ \Bigm|\ p\in\tfrac1m(mP\cap M)\Bigr\}
 +f_P(u).
\]
Since $P$ is a lattice polytope, its vertices lie in $\tfrac1m(mP\cap M)$
for every $m\ge1$, and a linear form attains its maximum on $P$ at a
vertex; so the displayed maximum is $\max_{p\in P}\langle u,p\rangle$.
Finally $\max_{p\in P}\langle u,p\rangle
=-\min_{p\in P}\langle-u,p\rangle=f_P(-u)$.  Thus $a_m(v_u)/m$ is
independent of $m$ and equal to $f_P(u)+f_P(-u)$, which is therefore
also $T_L(v_u)=T_\omega(v_u)$ by
\eqref{eq:bj-maximal-vanishing} and \Cref{prop:ch49-T-equals-T}.  The
identity $f_P(u)+f_P(-u)=\max_{p,q\in P}\langle u,p-q\rangle$ is the
definition of the width, and the listed properties are those of a support
function of the difference body $P-P$, which is a convex body symmetric
about the origin; it vanishes at $u$ only if $\langle u,\cdot\rangle$ is
constant on $P$, that is only if $u=0$, $P$ being
full-dimensional.
\end{proof}

\begin{proposition}
\label{prop:ch49-toric-dinfty}
For $u,u'\in N_{\mathbf R}$,
\[
 d_{\infty,\omega}(v_u,v_{u'})
 =\max\bigl\{\,f_P(u'-u)+f_P(u)-f_P(u'),
              \ f_P(u-u')+f_P(u')-f_P(u)\,\bigr\}.
\tag{49.24}\label{eq:ch49-toric-dinfty}
\]
Both entries of the maximum are nonnegative, the formula is symmetric in
$u,u'$, it reduces to \eqref{eq:ch49-toric-T} when $u'=0$, and it is
Lipschitz in $(u,u')$ with constant $2\sqrt2 R_P$ for the product
Euclidean norm, where $R_P:=\max_{p\in P}|p|$.
In particular $\iota$ is a $2R_P$-Lipschitz embedding of
$N_{\mathbf R}$ into $(X^{\lin},d_{\infty,\omega})$.
\end{proposition}

\begin{proof}
We use the section formula \eqref{eq:bj-dinfty-sections}, that is
\Cref{imp:ch49-imported}\,\cref{itm:ch49-import:sections}.  Fix $m\ge1$
and a nonzero $s\in H^0(X,mL)$ with exponent set
$S:=\Lambda(s)\subset P$.  By \eqref{eq:ch49-toric-section-value},
\[
 \frac{v_u(s)-v_{u'}(s)}{m}
 =\min_{p\in S}\langle u,p\rangle-\min_{p\in S}\langle u',p\rangle
  +f_P(u)-f_P(u') .
\]
Let $p^*\in S$ attain the second minimum.  Then
$\min_S\langle u,\cdot\rangle-\min_S\langle u',\cdot\rangle
\le\langle u-u',p^*\rangle\le\max_{p\in P}\langle u-u',p\rangle
=f_P(u'-u)$, so
\[
 \frac{v_u(s)-v_{u'}(s)}{m}\le f_P(u'-u)+f_P(u)-f_P(u') .
\]
Conversely, choose a vertex $p_0$ of $P$ attaining
$\max_{p\in P}\langle u-u',p\rangle$; it is a lattice point, so
$s:=\chi^{p_0}s_P\in H^0(X,L)$ is a monomial section with
$S=\{p_0\}$, and the displayed inequality is an equality for it.  Taking
the supremum over $m$ and $s$, and then the maximum with the same
expression with $u,u'$ interchanged, gives
\eqref{eq:ch49-toric-dinfty}.

Nonnegativity of each entry is subadditivity of the support function,
$f_P(u'-u)+f_P(u)\ge f_P(u')$.  Setting $u'=0$ and using
$f_P(0)=0$ gives $\max\{f_P(-u)+f_P(u),0\}=T_\omega(v_u)$, which is
\eqref{eq:ch49-toric-T} and
\Cref{prop:ch49-dinfty-basic}\,\cref{itm:ch49-dinf:vertex}.  For each exponent set $S\subset P$, both
$u\mapsto\min_{p\in S}\langle u,p\rangle$ and $f_P$ are
$R_P$-Lipschitz.  The section formula therefore gives
$d_{\infty,\omega}(v_u,v_{u'})\le2R_P|u-u'|$.
The triangle inequality for $d_{\infty,\omega}$ then yields
\[
 |d(v_u,v_{u'})-d(v_a,v_{a'})|
 \le2R_P(|u-a|+|u'-a'|)
 \le2\sqrt2 R_P|(u,u')-(a,a')|,
\]
which proves the stated product-norm bound.
\end{proof}

\begin{remark}[Two consistency checks on \eqref{eq:ch49-toric-dinfty}]
\label{rem:ch49-toric-dinfty-checks}
First, homogeneity: replacing $(u,u')$ by $(tu,tu')$ multiplies the
right-hand side by $t$, because $f_P$ is positively homogeneous; this
is \Cref{prop:ch49-dinfty-basic}\,\cref{itm:ch49-dinf:homog}.  Second, the
ray: for $u'=\lambda u$ with $\lambda\ge1$ one computes
$f_P((\lambda-1)u)+f_P(u)-f_P(\lambda u)=(\lambda-1)f_P(u)+f_P(u)
-\lambda f_P(u)=0$ and
$f_P((1-\lambda)u)+f_P(\lambda u)-f_P(u)
=(\lambda-1)f_P(-u)+(\lambda-1)f_P(u)=(\lambda-1)T_\omega(v_u)$, so
\[
 d_{\infty,\omega}(v_u,v_{\lambda u})=(\lambda-1)\,T_\omega(v_u),
\]
which is the equality case of
\Cref{lem:ch49-multiplicative-comparison} recorded in
\Cref{rem:ch49-comparison-sharp}.
\end{remark}

\subsection{Invariant psh functions are convex}
\label{sec:ch49-toric-dictionary}

\begin{theorem}\label{thm:bj-toric-psh-convex}
For a $T(k)$-invariant $\omega$-psh function $\phi$, set
\[
                         f_\phi:=(\phi+f_P)|_{N_{\mathbf R}}.
\tag{49.25}\label{eq:bj-toric-convex-function}
\]
Then $f_\phi$ is convex and $f_\phi\leq f_P+O(1)$.  The assignment is
a bijection from torus-invariant $\omega$-psh functions to such finite convex
functions, with inverse
\[
 f\longmapsto\text{the unique psh extension of }\rho^*(f-f_P)
                    \text{ from }X^{\val}.
\tag{49.26}\label{eq:bj-toric-psh-inverse}
\]
Legendre transform gives a second bijection with lower-semicontinuous
proper convex functions
\[
 g=f^\vee:P_{\mathrm{BJ}}\longrightarrow\mathbf R\cup\{+\infty\},
\qquad
 g(m)=\sup_{v\in N_{\mathbf R}}(\langle v,m\rangle-f(v)).
\tag{49.27}\label{eq:bj-toric-legendre}
\]
\end{theorem}

\begin{proof}
For any $\phi\in\PSH(\omega)$, a fixed-class Fubini--Study
regularization and \Cref{lem:ch49-toric-finite-level} give
$f_{\phi_j}\downarrow f_\phi$ on $N_{\mathbf R}$.
The limit is finite because $N_{\mathbf R}\subset X^{\lin}$,
hence convex.  Also $f_\phi\le f_P+\sup\phi$.

For invariant $\phi$, the retraction identity is the precise
published toric interface \cite[Appendix~B, p.~830]{BJ22}:
\[
 \phi|_{X^{\val}}=\rho^*(f_\phi-f_P).
 \tag{49.28}\label{eq:bj-toric-retraction-identity}
\]
The same interface supplies, for every finite convex
$f\le f_P+O(1)$, an invariant psh extension with this restriction.
Its uniqueness follows from divisorial determination of psh functions.
We quote this existence and invariant-extension statement explicitly;
the finite-level section calculation alone does not prove it.
The notation in \eqref{eq:bj-toric-psh-inverse} therefore does not
evaluate $\rho$ at a proper-support point.

It remains to verify the Legendre part without importing any boundary
convention.  Given such an $f$, extend $g=f^\vee$ to
$M_{\mathbf R}$.  If $m\notin P_{\mathrm{BJ}}$, strict separation gives
$v\in N_{\mathbf R}$ with
$\langle v,m\rangle>f_P(v)$.  Since
$f(tv)\leq t f_P(v)+C$, letting $t\to+\infty$ shows
$g(m)=+\infty$.  Thus $g$ is supported on
$P_{\mathrm{BJ}}$; it is lower semicontinuous and convex because it is a
supremum of continuous affine functions.  It is proper: a finite convex
function on all of $N_{\mathbf R}$ has a subgradient, and the preceding
separation argument forces every subgradient to lie in
$P_{\mathrm{BJ}}$.  Since $f$ is finite convex, it is continuous and
therefore closed, so Fenchel--Moreau gives $f=g^\vee$.

Conversely, let $g$ be a proper lower-semicontinuous convex function on
$P_{\mathrm{BJ}}$, extended by $+\infty$ outside the polytope.  Compactness
and lower semicontinuity give
$a:=\min_{P_{\mathrm{BJ}}}g\in\mathbf R$.  Its conjugate
\[
                  f(v)=g^\vee(v)
                     =\sup_{m\in P_{\mathrm{BJ}}}
                        \{\langle v,m\rangle-g(m)\}
\]
is finite convex and satisfies $f\leq f_P-a$.  Fenchel--Moreau gives
$f^\vee=g$, completing the second bijection.
\end{proof}

The extended value $+\infty$ on the polytope is essential for singular
psh functions.  No value $-\infty$ occurs for the convex function $f$
on $N_{\mathbf R}$, because $N_{\mathbf R}\subset X^{\lin}$.

The next lemmas give the finite section calculation and the retraction
comparison used in that proof; general invariant existence retains
the stated published input.

\begin{lemma}
\label{lem:ch49-toric-finite-level}
Let $\phi=\frac1m\max_{1\le i\le N}\{\log|s_i|+\lambda_i\}$ be a
Fubini--Study function for $L$, with $0\ne s_i\in H^0(X,mL)$ and
$\lambda_i\in\mathbf R$.  Then, on the toric slice,
\[
 f_\phi(u)=(\phi+f_P)(v_u)
 =\max_{1\le i\le N}\ \max_{p\in\Lambda(s_i)}
   \Bigl\{\langle u,-p\rangle+\frac{\lambda_i}{m}\Bigr\},
\tag{49.29}\label{eq:ch49-toric-fs-convex}
\]
with $\Lambda(s_i)\subset P$ as in
\eqref{eq:ch49-toric-section-value}.  This is a polyhedral convex
function with rational slopes and real affine constants; it is rational
polyhedral when the $\lambda_i$ are rational.  Its slopes $-p$ lie in
$P_{\mathrm{BJ}}=-P$; in particular
\[
 f_\phi\ \le\ f_P+\tfrac1m\max_i\lambda_i
 =f_P+\sup_{X^{\An}}\phi .
\]
Conversely, every rational polyhedral convex function on
$N_{\mathbf R}$ whose slopes lie in $P_{\mathrm{BJ}}$ and whose
constants are rational is of the form $f_\phi$ for an invariant
Fubini--Study function $\phi$ as above with monomial $s_i$.
\end{lemma}

\begin{proof}
By \eqref{eq:section-log-value} and
\eqref{eq:ch49-toric-section-value},
\[
 \frac{\log|s_i|(v_u)}{1}=-v_u(s_i)
 =-m\Bigl(\min_{p\in\Lambda(s_i)}\langle u,p\rangle+f_P(u)\Bigr),
\]
so that
\[
 \phi(v_u)+f_P(u)
 =\max_i\Bigl\{\frac{\lambda_i}{m}-\min_{p\in\Lambda(s_i)}\langle u,p\rangle
 \Bigr\}
 =\max_i\max_{p\in\Lambda(s_i)}
   \Bigl\{\frac{\lambda_i}{m}+\langle u,-p\rangle\Bigr\},
\]
because $-\min_p\langle u,p\rangle=\max_p\langle u,-p\rangle$.  This is
\eqref{eq:ch49-toric-fs-convex}: a finite maximum of affine functions,
hence polyhedral and convex, with rational slopes $-p\in-P$.  The upper
bound follows from $\langle u,-p\rangle\le f_P(u)$ for $-p\in-P$, and
$\sup_{X^{\An}}\phi=m^{-1}\max_i\lambda_i$ is
\Cref{cor:ch34-generic-constancy}.  For the converse, write the given function as
$\max_i\{\langle u,q_i\rangle+a_i\}$, where $q_i\in-P$ are rational.
Choose a common denominator $m$ with $-mq_i\in mP\cap M$.
The regular monomial sections $s_i=\chi^{-mq_i}s_P^m$ and
weights $\lambda_i=ma_i$ give the required presentation.
These are generically finite Fubini--Study functions; the selected
monomials need not be basepoint free.  Rational $a_i$ give rational
weights.  If a bounded difference from $f_P$ is required, all vertex
slopes must occur, and the full vertex system is basepoint free.
\end{proof}

\begin{lemma}
\label{lem:ch49-toric-retraction-dominates}
Let $\phi\in\PSH(\omega)$.  Then, on $X^{\val}$,
\[
 \phi\ \le\ \rho^*\bigl(f_\phi-f_P\bigr),
\]
with equality on $N_{\mathbf R}$.  Consequently, if
$\phi,\psi\in\PSH(\omega)$ are $T(k)$-invariant and
$f_\phi=f_\psi$, then $\phi$ and $\psi$ have the same restriction to
$N_{\mathbf R}$ and the same $\rho$-pullback; the assertion that they
are then equal is
\eqref{eq:bj-toric-retraction-identity} and is part of the imported toric
core.
\end{lemma}

\begin{proof}
By \Cref{lem:ch49-toric-slice}\,\cref{itm:ch49-ts:retraction} one has
$\iota\rho(v)\le v$ for $v\in X^{\val}$, and every $\omega$-psh
function is decreasing for \eqref{eq:semivaluation-order}; hence
$\phi(v)\le\phi(\iota\rho(v))=(f_\phi-f_P)(\rho(v))$.  Equality on
$N_{\mathbf R}$ is $\rho\circ\iota=\mathrm{id}$.
\end{proof}

\begin{remark}[The exact role of the toric input]
\label{rem:ch49-toric-core-use}
The convexity of the restriction of any psh function follows from
fixed-class regularization and the finite section calculation, without
requiring invariant approximants.  General invariant extension and
existence for arbitrary finite convex $f\le f_P+O(1)$ are the quoted
dictionary of Appendix~B.  A maximum of finitely many affine minorants
approximates a convex function from below; it does not by itself provide
a decreasing psh regularization.  For an already given invariant psh
function, the bounded-measure proof below constructs decreasing invariant
regularizations by replacing sections with their monomial terms.
\end{remark}

\section{Explicit toric measures and energies}
\label{sec:ch49-toric-ma}

Normalize lattice Lebesgue measure $\rd m$ so a fundamental parallelepiped
of $M$ has volume $1$, and let
\[
                         V_P:=\operatorname{vol}(P_{\mathrm{BJ}})
                            =\frac{(\omega^n)}{n!}.
\tag{49.30}\label{eq:bj-toric-volume}
\]

\begin{theorem}\label{thm:bj-toric-ma-energy}
For a bounded toric psh function,
\[
 \MA_\omega(\phi)=V_P^{-1}\iota_*
             \operatorname{MA}_{\mathbf R}(f_\phi),
\tag{49.31}\label{eq:bj-toric-ma}
\]
where $\iota:N_{\mathbf R}\hookrightarrow X^{\val}$, and
\[
                         E_\omega(\phi)
              =-\frac1{V_P}\int_{P_{\mathrm{BJ}}}f_\phi^\vee\,\rd m.
\tag{49.32}\label{eq:bj-toric-energy}
\]
Formula \eqref{eq:bj-toric-energy} extends to arbitrary toric psh
functions.  In particular,
\[
 \phi\in\mathcal E^1_{\mathrm{tor}}(\omega)
 \quad\Longleftrightarrow\quad
 f_\phi^\vee\in L^1(P_{\mathrm{BJ}}),
\tag{49.33}\label{eq:bj-toric-e1}
\]
and strong convergence of toric finite-energy potentials is equivalent to
$L^1(P_{\mathrm{BJ}})$-convergence of their Legendre transforms.
\end{theorem}

\begin{proof}[Proof architecture]
The proofs below proceed from the toric intersection formula for
rational PL potentials to bounded measures, then integrate the first
variation to obtain the bounded energy formula.  Decreasing truncations
give the singular energy formula and its $L^1$-finiteness criterion.
The full equivalence between strong convergence and $L^1$ convergence
remains the stated Appendix~B input of \cite{BJ22}.
After the reflection
\eqref{eq:bj-toric-sign-conversion}, it agrees with the unnormalized formula
of \Cref{chap:model-calculations}:
\[
 \MA_\omega^{\mathrm u}(\phi)
       =n!\,\iota_*\operatorname{MA}_{\mathbf R}(f_\phi).
\]
\end{proof}

The rest of this section supplies the details, in the order: the mass
normalization \eqref{eq:bj-toric-volume}; the measure
\eqref{eq:bj-toric-ma} at the polyhedral level and then for bounded
$\phi$; the energy \eqref{eq:bj-toric-energy} by an exact convex-analytic
integration of the first variation, which needs no further toric input; and
the singular case \eqref{eq:bj-toric-e1}.

\subsection{The mass normalization}
\label{sec:ch49-toric-mass}

\begin{samepage}
\begin{lemma}
\label{lem:ch49-toric-volume}
Let $f:N_{\mathbf R}\to\mathbf R$ be convex with $f-f_P$ bounded.
Then
\[
 \overline{\partial f(N_{\mathbf R})}=P_{\mathrm{BJ}},
 \qquad
 \MA_{\mathbf R}(f)(N_{\mathbf R})=\vol_M(P_{\mathrm{BJ}})=V_P,
 \qquad
 V=(\omega^n)=n!\,V_P .
\]
Moreover $\MA_{\mathbf R}(f)=(\nabla f^\vee)_*(\rd m|_{P_{\mathrm{BJ}}})$.
\end{lemma}
\end{samepage}

\begin{proof}
A bounded difference gives the same recession function, and the closure of
the total subgradient image of a finite convex function is the closure of
the domain of its conjugate, which for $f_P$ is $P_{\mathrm{BJ}}$; this
is \Cref{ex:appD-indicator-support} and
\Cref{prop:appD-legendre-pushforward}, whose \eqref{eq:appD-real-ma-mass}
and \eqref{eq:appD-legendre-pushforward} are the two displayed identities.
The relation $V=n!\,V_P$ is the toric degree formula
\eqref{eq:appJ-toric-degree}, applied to $P$, whose lattice volume equals
that of $-P=P_{\mathrm{BJ}}$.  The same statement over a nontrivially
valued base is \Cref{lem:toric-real-total-mass}.
\end{proof}

Thus \eqref{eq:bj-toric-ma} is consistent on total masses: both sides are
probability measures.  Equivalently, in unnormalized form,
\[
 \MA^{\mathrm u}_\omega(\phi)=n!\ \iota_*\MA_{\mathbf R}(f_\phi),
\]
which is exactly \eqref{eq:toric-ma-comparison} of
\Cref{chap:model-calculations} after the reflection
\eqref{eq:bj-toric-sign-conversion}, and which passes the $n!$-test of
\Cref{app:convexity}.

\subsection{The measure}
\label{sec:ch49-toric-measure}

\begin{proposition}
\label{prop:ch49-toric-ma-pl}
Let $\phi\in\PL(X)\cap\PSH(\omega)$ be $T(k)$-invariant.
Its restriction $f_\phi$ is rational polyhedral and convex:
rational PL presentations restrict to rational polyhedral functions
on the toric slice, and convexity follows from
\Cref{thm:bj-toric-psh-convex}.  Its slopes lie in $P_{\mathrm{BJ}}$.  Then
\eqref{eq:bj-toric-ma} holds; explicitly, $\MA_\omega(\phi)$ is the
finite atomic measure
\[
 \MA_\omega(\phi)
 =\frac1{V_P}\sum_{u}\vol_M\bigl(\partial f_\phi(u)\bigr)\,\delta_{v_u},
\]
the sum being over points $u$ with full-dimensional subdifferential
(the vertices of the induced polyhedral decomposition of
$f_\phi$).
\end{proposition}

\begin{proof}
By \Cref{prop:ch46-pl-measure} the measure $\MA^{\mathrm u}_\omega(\phi)$
is the finite atomic measure
$\sum_Eb_E(\cL^n\cdot E)\delta_{v_E}$, the sum over the components of the
central fibre of a test configuration $\cX$ determining $\phi$, with
$\cL=\omega_\cX+D_\phi$.  Since $\phi$ is invariant we may take $\cX$
toric: by \Cref{thm:tc-pl-correspondence} and
\Cref{lem:ch49-toric-finite-level} a positive multiple of $f_\phi$ is
the support-type function of a $T\times\mathbf G_m$-invariant polarized
toric degeneration of $(X,L)$ over $\mathbf A^1$, whose central-fibre
components $E_u$ correspond to the vertices $u$ of the corner locus of
$f_\phi$, with $v_{E_u}=v_u$ after the normalization of
\Cref{def:ch33-component-valuation}, and whose residual intersection
numbers are
\[
 b_{E_u}\,(\cL^n\cdot E_u)=n!\,\vol_M\bigl(\partial f_\phi(u)\bigr).
\]
The last identity is the toric intersection formula in the form
\cite[Theorem~13.4.3]{CLS11}; it is the same computation as in the proof of
\Cref{thm:toric-ma-comparison}, where it is carried out over a
nontrivially valued base.  Dividing by $V=n!V_P$ and using
\eqref{eq:appD-polyhedral-atom} gives the displayed formula, and the
displayed formula is \eqref{eq:bj-toric-ma} because a polyhedral convex
function has purely atomic Alexandrov measure supported on the vertices of
its corner locus (\Cref{lem:ch08-alexandrov-polyhedral}).
\end{proof}

\begin{proposition}
\label{prop:ch49-toric-ma-bounded}
Formula \eqref{eq:bj-toric-ma} holds for every bounded $T(k)$-invariant
$\phi\in\PSH(\omega)$, that is, whenever $f_\phi-f_P$ is bounded.
\end{proposition}

\begin{proof}
Take a decreasing fixed-class regularization $\eta_j\in\cH(L)$ of
$\phi$.  Expand each section in its toric monomial basis, and replace
it in the weighted maximum by all its nonzero monomial terms.
The resulting rational Fubini--Study function $\psi_j$ is invariant,
has the same restriction to $N_{\mathbf R}$ as $\eta_j$ by
\Cref{lem:ch49-toric-finite-level}, and is basepoint free:
a common zero of all the selected monomials would be a common zero
of the original basepoint-free sections.  On $X^{\val}$ it is
the retraction pullback of that restriction.  Therefore the
$\psi_j$ decrease on $X^{\val}$, and divisorial comparison extends
their order to $X^{\An}$.  Their limit is psh, since its value at
the trivial valuation is the finite number $\sup\phi$, and agrees
with $\phi$ on $X^{\val}$ by
\eqref{eq:bj-toric-retraction-identity}.  Divisorial determination
gives $\psi_j\downarrow\phi$ everywhere.

Pass to a cofinal tail below one $\psi_{j_0}$.
Since $\psi_j\ge\phi$ and $\phi$ is bounded, there are common finite
upper and lower bounds on this tail.  Thus
$f_{\psi_j}-f_P$ is uniformly bounded and
$f_{\psi_j}\downarrow f_\phi$ locally uniformly on $N_{\mathbf R}$
by the compact-set Dini argument for decreasing continuous convex
functions.  Alexandrov continuity gives vague convergence
$\MA_{\mathbf R}(f_{\psi_j})\to\MA_{\mathbf R}(f_\phi)$.
The limit and approximants all have total mass $V_P$ by
\Cref{lem:ch49-toric-volume}, so this convergence is weak against
bounded continuous functions as well.  Indeed, choose a compactly
supported cutoff capturing all but $\varepsilon$ of the limiting
mass; vague convergence and the common total mass bound the
remaining tails by $2\varepsilon$ eventually.

The rational intersection formula
\Cref{prop:ch49-toric-ma-pl} applies to every $\psi_j$.
On the analytic side, decreasing mixed-MA continuity gives
$\MA_\omega(\psi_j)\to\MA_\omega(\phi)$.
Testing the real measures by the bounded continuous restrictions of
continuous functions on $X^{\An}$ and passing to the limit proves
\eqref{eq:bj-toric-ma}.  This includes continuous real-PL potentials
and agrees with \cite[Appendix~B, p.~831]{BJ22}.
\end{proof}

\begin{corollary}
\label{cor:ch49-toric-ma-e1}
The measure formula \eqref{eq:bj-toric-ma} also holds for every
invariant $\phi\in\cE^1(\omega)$.
\end{corollary}

\begin{proof}
The energy formula of \Cref{prop:ch49-toric-energy-singular}, proved
from the bounded measure case without using this corollary, gives $g=f_\phi^\vee\in L^1(P_{\mathrm{BJ}})$.
Its convex effective domain contains the interior of the polytope:
a convex subset of full Lebesgue measure has every interior point
in its interior.  At every such point $m$, $g$ has a subgradient
$u$, and Fenchel reciprocity gives $m\in\partial f_\phi(u)$.
Conversely $\partial f_\phi(N_{\mathbf R})\subset P_{\mathrm{BJ}}$.
Thus the Alexandrov mass is exactly $V_P$.

The bounded invariant truncations
$\phi_c=\max\{\phi,\sup\phi-c\}$ decrease to $\phi$.
Their convex functions converge locally uniformly to $f_\phi$,
so Alexandrov continuity gives vague convergence of their real measures.
Since the limit and all approximants have mass $V_P$, the cutoff
argument in \Cref{prop:ch49-toric-ma-bounded} gives weak convergence
on $N_{\mathbf R}$.  Decreasing finite-energy continuity gives
$\MA_\omega(\phi_c)\to\MA_\omega(\phi)$ on $X^{\An}$.
Testing the bounded real-measure identity against restrictions of
continuous analytic functions proves the claim.
No use of this corollary is needed to establish the energy formula.
\end{proof}

\subsection{The energy}
\label{sec:ch49-toric-energy}

The energy formula now follows from the measure formula by an exact
computation in convex analysis, with no further toric geometry.  This is
the toric analogue of the derivation of \Cref{thm:ch26-toric-energy} over a
discretely valued base.  Both arguments integrate the first variation after the stated
normalization.

\begin{proposition}
\label{prop:ch49-toric-energy-bounded}
Let $\phi$ be a bounded $T(k)$-invariant $\omega$-psh function.
Then \eqref{eq:bj-toric-energy} holds:
\[
 E_\omega(\phi)=-\frac1{V_P}\int_{P_{\mathrm{BJ}}}f_\phi^\vee\,\rd m .
\]
\end{proposition}

\begin{proof}
Put $\phi_t:=t\phi$ for $t\in[0,1]$; each $\phi_t$ is
$\omega$-psh by \Cref{thm:bj-basic-closure}\,(3), invariant, and
bounded, and
\[
 f_{\phi_t}=(t\phi+f_P)|_{N_{\mathbf R}}=(1-t)f_P+t\,f_\phi
\]
because $\psi\mapsto f_\psi=\psi+f_P$ is affine.  By the segment
expansion \Cref{thm:ch47-segment-expansion}, or by
\Cref{prop:ch46-first-variation-pl} and approximation, $t\mapsto
E_\omega(\phi_t)$ is a polynomial of degree at most $n+1$ with
\[
 \frac{\rd}{\rd t}E_\omega(\phi_t)=\int_{X^{\An}}\phi\,\MA_\omega(\phi_t).
\]
By \Cref{prop:ch49-toric-ma-bounded} and
\Cref{lem:ch49-toric-volume}, and writing $u_t(m):=\nabla f_{\phi_t}^\vee(m)$
for the (almost everywhere defined) gradient of the conjugate,
\[
 \int_{X^{\An}}\phi\,\MA_\omega(\phi_t)
 =\frac1{V_P}\int_{N_{\mathbf R}}(f_\phi-f_P)\,\MA_{\mathbf R}(f_{\phi_t})
 =\frac1{V_P}\int_{P_{\mathrm{BJ}}}(f_\phi-f_P)\bigl(u_t(m)\bigr)\,\rd m .
\]
Put $g_t=f^\vee_{\phi_t}$ and $C=\|f_\phi-f_P\|_\infty$.
The variational definition of the conjugate gives, on $P_{\mathrm{BJ}}$,
\[
 |g_t(m)-g_s(m)|\le C|t-s|.
\]
Thus $t\mapsto g_t(m)$ is absolutely continuous, uniformly in $m$.
For each fixed $t$, the finite convex function $g_t$ is differentiable
at almost every point of the interior of $P_{\mathrm{BJ}}$.
At such a point, Fenchel reciprocity identifies the unique maximizer
in its conjugate formula as $u_t(m)=\nabla g_t(m)$.
For almost every pair $(t,m)$, differentiation in $t$ also exists.
Evaluating the variational formula at this maximizer for nearby
parameters gives the envelope derivative
\[
 \partial_t g_t(m)=-(f_\phi-f_P)(u_t(m)).
\]
Uniqueness here comes from differentiability in $m$, not merely from
differentiability of a convex function in $t$.
The uniform Lipschitz bound and Fubini's theorem now give
\[
 -\frac1{V_P}\int_{P_{\mathrm{BJ}}}(g_1-g_0)\,\rd m
 =\int_0^1\frac1{V_P}\int_{P_{\mathrm{BJ}}}
       (f_\phi-f_P)(u_t(m))\,\rd m\,\rd t
 =E_\omega(\phi)-E_\omega(0).
\]
The polytope boundary has Lebesgue measure zero.  Since $g_0=0$
on $P_{\mathrm{BJ}}$ and $E_\omega(0)=0$, this is
\eqref{eq:bj-toric-energy}.
\end{proof}

\begin{proposition}
\label{prop:ch49-toric-energy-singular}
Formula \eqref{eq:bj-toric-energy} holds for every $T(k)$-invariant
$\phi\in\PSH(\omega)$, both sides being elements of
$[-\infty,\sup_{X^{\An}}\phi]$.  Consequently
\eqref{eq:bj-toric-e1} holds: an invariant $\omega$-psh $\phi$ has
finite energy if and only if $f^\vee_\phi\in L^1(P_{\mathrm{BJ}},\rd m)$.
\end{proposition}

\begin{proof}
Let $\phi_c:=\max\{\phi,\sup\phi-c\}$ for $c>0$; these are bounded,
invariant, $\omega$-psh, and decrease to $\phi$ as $c\uparrow\infty$.
On the convex side, $f_{\phi_c}=\max\{f_\phi,f_P+\sup\phi-c\}$ decreases
to $f_\phi$, and the conjugates increase to $f^\vee_\phi$ pointwise:
the two suprema over $c$ and $u$ commute in
$\sup_c\sup_u(\langle u,m\rangle-f_{\phi_c}(u))$.
By \Cref{prop:ch46-energy-usc}\,\cref{itm:ch46-Eusc:decreasing},
$E_\omega(\phi_c)\downarrow E_\omega(\phi)$; by monotone convergence,
$\int f^\vee_{\phi_c}\rd m\uparrow\int f^\vee_\phi\rd m$ in
$(-\infty,+\infty]$, on the tail $c\ge1$ the integrands are bounded below by
$f^\vee_{\phi_1}\ge-\|f_{\phi_1}-f_P\|_\infty$.  Applying
\Cref{prop:ch49-toric-energy-bounded} to each $\phi_c$ and passing to the
limit gives \eqref{eq:bj-toric-energy}, with both sides possibly
$-\infty$.

For \eqref{eq:bj-toric-e1}, recall that $\phi\in\cE^1(\omega)$ means
$E_\omega(\phi)>-\infty$ (\Cref{def:bj-e1}); by what has just been proved
this is $\int_{P_{\mathrm{BJ}}}f^\vee_\phi\,\rd m<+\infty$, and
$f^\vee_\phi$ is bounded below on $P_{\mathrm{BJ}}$ by
$-\sup_{X^{\An}}\phi$, so finiteness of the integral is exactly
$f^\vee_\phi\in L^1$.
\end{proof}

\begin{proposition}
\label{prop:ch49-toric-strong-L1}
For any net of invariant $\phi_j,\phi\in\cE^1(\omega)$,
\[
 \phi_j\longrightarrow\phi\ \text{strongly}
 \quad\Longleftrightarrow\quad
 f^\vee_{\phi_j}\longrightarrow f^\vee_\phi
       \ \text{in }L^1(P_{\mathrm{BJ}},\rd m).
\]
No monotonicity or separate convergence-of-suprema hypothesis is required.
\end{proposition}

\begin{proof}
This is the full strong/L1 interface of
\cite[Appendix~B, p.~831]{BJ22}, quoted at its finite-energy domain.
For a decreasing net, one part can also be checked directly:
the conjugates increase pointwise, their integrals converge by the
energy formula and decreasing energy continuity, and therefore
their L1 distance to the limiting conjugate tends to zero.
The general converse requires the convex convergence statement supplied
by the published interface.  Almost-everywhere convergence and Fatou's
lemma alone do not justify passing a supremum through a limit, nor
do they prove the assertion for arbitrary nets.
\end{proof}

\subsection{Green functions and Dirac energies, in closed form}
\label{sec:ch49-toric-green}

\begin{example}[$\mathbf P^1$]
\label{ex:ch49-toric-trivial-p1}
Take the polytope $P=[0,d]$ of \Cref{chap:model-calculations}.  Then
$P_{\mathrm{BJ}}=[-d,0]$ and
\[
 f_P(u)=\max\{0,-du\}.
\]
For the trivial metric, $f_\phi=f_P$, whose Alexandrov measure is
$d\delta_0$.  Since $V_P=d$, the probability Monge--Amp\`ere measure is
$\delta_0=\delta_{v_{\triv}}$.  This verifies both the sign and the mass.
\end{example}

\begin{example}[Green function of a toric valuation]
\label{ex:ch49-toric-green-formula}
For $v\in N_{\mathbf R}$, the normalized Green function satisfies
\[
 f_{\phi_v}(w)=f_P(w-v)+f_P(v),
\qquad
 T_\omega(v)=f_P(v)+f_P(-v).
\tag{49.34}\label{eq:bj-toric-green}
\]
Its Monge--Amp\`ere measure is $\delta_v$.  The formula is a translated
support function, so its Alexandrov measure is concentrated at $v$, as
required.
\end{example}

\begin{proposition}
\label{prop:ch49-toric-green}
Let $u\in N_{\mathbf R}$ and let $\phi_u$ be the invariant
$\omega$-psh function with $f_{\phi_u}(w)=f_P(w-u)+f_P(u)$, as in
\eqref{eq:bj-toric-green}.  Then:
\begin{enumerate}
\item\label{itm:ch49-tg:ma}
      $\MA_\omega(\phi_u)=\delta_{v_u}$ and $\phi_u(v_u)=0$; thus
      $\phi_u$ is the Green function of $v_u$ in the sense of
      \Cref{ex:ch47-dirac-green}, and it is a rational polyhedral, hence
      $\mathbf Q$-PL, function whenever $u$ is rational.
\item\label{itm:ch49-tg:sup}
      $\sup_{X^{\An}}\phi_u=\phi_u(v_{\triv})=T_\omega(v_u)$; in
      particular this agrees with the general Green-height identity of
      \Cref{cor:ch49-green-invariants}.
\item\label{itm:ch49-tg:energy}
      With $b_P:=V_P^{-1}\int_Pp\,\rd p$ the barycenter of $P$,
      \[
       E^\vee_\omega(\delta_{v_u})=E_\omega(\phi_u)
       =f_P(u)+\langle u,b_P\rangle
       =\langle u,b_P\rangle-\min_{p\in P}\langle u,p\rangle .
      \tag{49.35}\label{eq:ch49-toric-dirac-energy}
      \]
\end{enumerate}
\end{proposition}

\begin{proof}
\cref{itm:ch49-tg:ma}  The function $w\mapsto f_P(w-u)$ is the support
function of $P_{\mathrm{BJ}}$ translated by $u$, so its subdifferential
is $P_{\mathrm{BJ}}$ at $w=u$ and a proper face elsewhere; by
\eqref{eq:appD-polyhedral-atom} its Alexandrov measure is
$\vol_M(P_{\mathrm{BJ}})\delta_u=V_P\delta_u$, and adding the constant
$f_P(u)$ changes nothing.  Hence $\MA_\omega(\phi_u)=\delta_{v_u}$ by
\eqref{eq:bj-toric-ma}.  Also
$\phi_u(v_u)=f_{\phi_u}(u)-f_P(u)=f_P(0)+f_P(u)-f_P(u)=0$.  Rationality
is clear from \Cref{lem:ch49-toric-finite-level}, the slopes being the
vertices of $P_{\mathrm{BJ}}$ and the constants $\langle u,-p\rangle$.

\cref{itm:ch49-tg:sup}
$\phi_u(v_{\triv})=f_{\phi_u}(0)-f_P(0)=f_P(-u)+f_P(u)=T_\omega(v_u)$ by
\eqref{eq:ch49-toric-T}, and $\sup\phi=\phi(v_{\triv})$.

\cref{itm:ch49-tg:energy}  For $m\in P_{\mathrm{BJ}}$, substituting
$w=u+w'$,
\[
 \begin{aligned}
  f^\vee_{\phi_u}(m)
   &=\sup_{w}\bigl(\langle w,m\rangle-f_P(w-u)-f_P(u)\bigr)\\
   &=\langle u,m\rangle-f_P(u)+f_P^\vee(m)\\
   &=\langle u,m\rangle-f_P(u),
 \end{aligned}
\]
since $f_P^\vee$ is the indicator of $P_{\mathrm{BJ}}$
(\Cref{ex:appD-indicator-support}).  Therefore, by
\eqref{eq:bj-toric-energy},
\[
 E_\omega(\phi_u)
 =-\frac1{V_P}\int_{P_{\mathrm{BJ}}}\bigl(\langle u,m\rangle-f_P(u)\bigr)
   \rd m
 =f_P(u)-\bigl\langle u,\bar m\bigr\rangle,
\]
where $\bar m=V_P^{-1}\int_{P_{\mathrm{BJ}}}m\,\rd m=-b_P$ is the
barycenter of $P_{\mathrm{BJ}}=-P$.  This gives
\eqref{eq:ch49-toric-dirac-energy}, and
$E^\vee_\omega(\delta_{v_u})=E_\omega(\phi_u)$ is
\Cref{cor:ch49-green-invariants}.
\end{proof}

\begin{corollary}
\label{cor:ch49-minkowski}
\index{barycenter!Minkowski inequality}
Let $P\subset M_{\mathbf R}\simeq\mathbf R^n$ be a full-dimensional
lattice polytope with barycenter $b_P$, and for $u\in N_{\mathbf R}$
put $h(u):=\min_{p\in P}\langle u,p\rangle$ and
$H(u):=\max_{p\in P}\langle u,p\rangle$.  Then
\[
 \frac1{n+1}\ \le\
 \frac{\langle u,b_P\rangle-h(u)}{H(u)-h(u)}\ \le\ \frac n{n+1}
 \qquad (u\ne0).
\]
\end{corollary}

\begin{proof}
Take $X=X_\Sigma$ with $\Sigma$ the inner normal fan of $P$ and
$L=L_P$, so that $\omega=c_1(L)$ is ample.  By
\Cref{prop:ch49-toric-T}, $T_\omega(v_u)=H(u)-h(u)$, and by
\Cref{prop:ch49-toric-green}\,\cref{itm:ch49-tg:energy},
$E^\vee_\omega(\delta_{v_u})=\langle u,b_P\rangle-h(u)$.  Now apply
\eqref{eq:bj-dirac-energy-T}.
\end{proof}

\begin{remark}[A classical statement, and a check]
\label{rem:ch49-minkowski-check}
\Cref{cor:ch49-minkowski} is the classical inequality of Minkowski for the
centroid of a convex body: the centroid of $P$ divides every slab
$\{h\le\langle u,\cdot\rangle\le H\}$ containing $P$ in a ratio between
$1:n$ and $n:1$; see \cite{Sch14}.  It extends from lattice polytopes to
arbitrary convex bodies by scaling and approximation, both sides being
continuous for the Hausdorff metric.  We have derived it from
\eqref{eq:bj-dirac-energy-T}; the logical direction could of course be
reversed, and the reader who prefers to regard the convex statement as
known may read \Cref{cor:ch49-minkowski} as a nontrivial confirmation of
the two constants in \Cref{thm:bj-dirac-finite-energy}.  The extremal cases
match: equality on the right holds when $P$ is a cone over a facet with
apex on $\{ \langle u,\cdot\rangle=h\}$, and on the left when the apex is
on $\{\langle u,\cdot\rangle=H\}$.  The reason is that the
cross-sectional measure of a cone at relative height $s$ above its apex
is proportional to $s^{n-1}$, so the centroid sits at relative height
$n/(n+1)$ \emph{above the apex}.  For the standard simplex
$P=\operatorname{conv}\{0,e_1^\vee,\dots,e_n^\vee\}$ and
$u=(1,\dots,1)$ one gets $h=0$, $H=1$,
$\langle u,b_P\rangle=n/(n+1)$, the right-hand extremal case; and for
$u=e_1$ one gets $h=0$, $H=1$,
$\langle u,b_P\rangle=1/(n+1)$, the left-hand one.  In pluripotential
terms, on $\mathbf P^n$ the two extremes are realized by
$v_{(1,\dots,1)}=\ord_0$ and by $v_{e_1}=\ord_{D_1}$, the order of
vanishing along a coordinate hyperplane; see
\Cref{ex:ch49-p2-dictionary}.
\end{remark}

\begin{warning}{Divisorial poles need not give model Green functions}
\label{warn:ch49-green-not-pl}
The toric formula above is a model-friendly special case.
For smooth projective varieties over an algebraically closed trivially
valued field of characteristic zero and a rational ample class,
Boucksom--Jonsson exhibit divisorial Green functions that are not $\mathbf Q$-PL
\cite[Theorem~60 and Examples~63, 65]{BJGreen24}.
Theorem~A of that paper has a discretely valued hypothesis and is
a separate result.  Divisorial support of the pole therefore
does not justify replacing a finite-energy Green function by a rational PL
one.
\end{warning}

\section{Worked examples}
\label{sec:ch49-examples}

The following computations illustrate the chapter.  The first is the toric dictionary on
$\mathbf P^2$, carried to explicit numbers and to a picture of the
polytope; the second is a quasimonomial valuation on a blown-up surface,
whose linear-growth constant is computed twice; the third is a valuation
without linear growth, with the finite-energy dichotomy exhibited on it;
the fourth is a non-toric divisorial point at which the lower constant of
\eqref{eq:bj-dirac-energy-T} is attained; and the fifth is a set of unit
tests against \Cref{chap:model-calculations},
\Cref{chap:bfj-energy}, \Cref{chap:bj-energy} and
\Cref{chap:bj-finite-energy}.

\subsection{The toric dictionary on
\texorpdfstring{$\mathbf P^2$}{P2}, with the polytope}
\label{sec:ch49-p2}

\begin{example}[The dictionary on $\mathbf P^2$]
\label{ex:ch49-p2-dictionary}
\index{toric geometry!projective plane}
Let $X=\mathbf P^2_k$ with homogeneous coordinates $z_0,z_1,z_2$, let
$T=\mathbf G_m^2$ act by $t\cdot[z_0:z_1:z_2]=[z_0:t_1z_1:t_2z_2]$, and
put $x:=z_1/z_0=\chi^{e^\vee_1}$, $y:=z_2/z_0=\chi^{e^\vee_2}$, so that
$M=\mathbf Z^2$ with basis $e^\vee_1,e^\vee_2$ and
$N=\mathbf Z^2$ with dual basis $e_1,e_2$.  Take
$L=\cO(1)$, $\omega=c_1(L)$, $n=2$, $V=(\omega^2)=1$.

\emph{The polytope.}  With the toric section $s_P=z_0$, the polytope of
$L$ is the standard simplex
\[
 P=\operatorname{conv}\{0,\ e^\vee_1,\ e^\vee_2\}
  =\{p\in M_{\mathbf R}\ :\ p_1\ge0,\ p_2\ge0,\ p_1+p_2\le1\},
\]
of lattice volume $\vol_M(P)=\tfrac12$; and
\[
 P_{\mathrm{BJ}}=-P=\operatorname{conv}\{0,-e^\vee_1,-e^\vee_2\},
 \qquad
 V_P=\vol_M(P_{\mathrm{BJ}})=\tfrac12=\frac{(\omega^2)}{2!},
\]
in accordance with \eqref{eq:bj-toric-volume}.  Its support function is
\[
 f_P(u)=\sup_{m\in P_{\mathrm{BJ}}}\langle u,m\rangle
 =\max\{0,-u_1,-u_2\},
\]
which by \eqref{eq:bj-toric-sign-conversion} is the reflected support
function $\check h_P$ of \Cref{ex:ch08-toric-pn}.

\begin{center}
\begin{tikzpicture}[scale=1.7,>=Stealth,font=\scriptsize]
 \draw[->,gray!70] (-1.35,0)--(1.35,0) node[right]{$m_1$};
 \draw[->,gray!70] (0,-1.35)--(0,1.3) node[above]{$m_2$};
 \fill[blue!12] (0,0)--(1,0)--(0,1)--cycle;
 \draw[blue!70!black,thick] (0,0)--(1,0)--(0,1)--cycle;
 \node[blue!55!black] at (0.33,0.31) {$P$};
 \fill[red!8] (0,0)--(-1,0)--(0,-1)--cycle;
 \draw[red!70!black,thick] (0,0)--(-1,0)--(0,-1)--cycle;
 \fill[red!32] (0,0)--(-0.55,0)--(0,-0.55)--cycle;
 \draw[densely dashed] (-0.55,0)--(0,-0.55);
 \node[red!55!black] at (-0.72,-0.15) {$P_{\mathrm{BJ}}$};
 \fill (0,0) circle (0.018);
 \fill[blue!70!black] (1,0) circle (0.018);
 \fill[blue!70!black] (0,1) circle (0.018);
 \fill[red!70!black] (-1,0) circle (0.018);
 \fill[red!70!black] (0,-1) circle (0.018);
 \node[below right] at (1,0) {$e^\vee_1$};
 \node[above left] at (0,1) {$e^\vee_2$};
 \draw[->,gray!60] (-1.55,0.72)--(-0.20,-0.14);
 \node[left,align=right] at (-1.55,0.72)
   {$\nabla f^\vee=(1,1)$,\\ area $c^2/2$};
 \draw[->,gray!60] (-1.55,-1.00)--(-0.32,-0.56);
 \node[left,align=right] at (-1.55,-1.00)
   {$\nabla f^\vee=(0,0)$,\\ area $(1-c^2)/2$};
 \draw[->,gray!60] (0.95,-0.85)--(-0.20,-0.37);
 \node[right] at (0.95,-0.85) {$m_1+m_2=-c$};
\end{tikzpicture}
\end{center}

\emph{Five toric points, identified.}  Under
\eqref{eq:bj-toric-valuation-pairing}:
$v_{(0,0)}=v_{\triv}$;
$v_{(1,1)}=\ord_0$, the multiplicity valuation at the torus-fixed point
$0=[1:0:0]$, which is $\ord_{E_1}$ for the blowup of $\mathbf P^2$
at $0$;
$v_{(1,0)}=\ord_{D_1}$ with $D_1=\{z_1=0\}$, since
$v(x)=1$ and $v(y)=0$;
$v_{(-1,-1)}=\ord_{D_0}$ with $D_0=\{z_0=0\}$, since
$\ord_{D_0}(z_1/z_0)=-1$;
and $v_{(1,2)}=\ord_{E_2}$, the second exceptional divisor of
\Cref{ex:ch49-blowup-surface} below.

\emph{The two invariants.}  By \eqref{eq:ch49-toric-T} and
\eqref{eq:ch49-toric-dirac-energy}, with barycenter
$b_P=(\tfrac13,\tfrac13)$:
\begin{center}
\begin{tabular}{ccccc}
\toprule
$u$ & valuation & $T_\omega$ & $E^\vee_\omega(\delta_{v_u})$
& $E^\vee_\omega/T_\omega$\\
\midrule
$(0,0)$ & $v_{\triv}$ & $0$ & $0$ & undefined\\
$(1,1)$ & $\ord_0$ & $1$ & $\tfrac23$ & $\tfrac23$\\
$(1,0)$ & $\ord_{D_1}$ & $1$ & $\tfrac13$ & $\tfrac13$\\
$(-1,-1)$ & $\ord_{D_0}$ & $1$ & $\tfrac13$ & $\tfrac13$\\
$(1,2)$ & $\ord_{E_2}$ & $2$ & $1$ & $\tfrac12$\\
$(1,-1)$ & $v_{(1,-1)}$ & $2$ & $1$ & $\tfrac12$\\
\bottomrule
\end{tabular}
\end{center}
Every defined entry of the last column lies in
$[\tfrac13,\tfrac23]=[\tfrac1{n+1},\tfrac n{n+1}]$, as
\eqref{eq:bj-dirac-energy-T} requires, and both endpoints occur: at
$\ord_0$ on the right, at $\ord_{D_1}$ and $\ord_{D_0}$ on the left.
By \Cref{cor:ch49-minkowski} this is Minkowski's centroid inequality for
the triangle $P$, whose extremal directions are exactly those normal to
a facet or pointing at a vertex.

\emph{The metric.}  From \eqref{eq:ch49-toric-dinfty}, for instance
\[
 d_{\infty,\omega}\bigl(v_{(1,1)},v_{(1,0)}\bigr)
 =\max\{f_P(0,-1)+0-0,\ f_P(0,1)+0-0\}=\max\{1,0\}=1,
\]
\[
 d_{\infty,\omega}\bigl(v_{(1,1)},v_{(-1,-1)}\bigr)
 =\max\{f_P(-2,-2)+0-1,\ f_P(2,2)+1-0\}=\max\{1,1\}=1 .
\]
Both are consistent with
$|T_\omega(v)-T_\omega(w)|\le d_{\infty,\omega}(v,w)\le
T_\omega(v)+T_\omega(w)$, which here reads $0\le d\le2$.

\emph{The Monge--Amp\`ere measure of the model family.}  Let
$c\in[0,1]$ and, as in \Cref{ex:ch46-toric-p2},
\[
 \begin{aligned}
  \phi_c&:=\max\{c\log|\mathfrak m_0|,-c\},\\
  f_{\phi_c}&=f_P+\phi_c,\\
  f^\vee_{\phi_c}(m)&=(c+m_1+m_2)_+
    \quad\text{on }P_{\mathrm{BJ}}.
 \end{aligned}
\]
The conjugate has gradient $(1,1)$ on
$\{m\in P_{\mathrm{BJ}}:m_1+m_2>-c\}$, the shaded triangle of the figure,
of area $c^2/2$, and gradient $(0,0)$ on the complement, of area
$(1-c^2)/2$.  By \eqref{eq:appD-legendre-pushforward} and
\eqref{eq:bj-toric-ma},
\[
 \MA_{\mathbf R}(f_{\phi_c})
 =\frac{c^2}2\,\delta_{(1,1)}+\frac{1-c^2}2\,\delta_{(0,0)},
 \qquad
 \MA_\omega(\phi_c)
 =c^2\,\delta_{\ord_0}+(1-c^2)\,\delta_{v_{\triv}} ,
\]
of total mass $1$; this is the measure computed by intersection theory in
\Cref{ex:ch46-toric-p2}.  By \eqref{eq:bj-toric-energy},
\[
 E_\omega(\phi_c)
 =-2\int_{P_{\mathrm{BJ}}}(c+m_1+m_2)_+\,\rd m
 =-2\int_0^c(c-t)\,t\,\rd t=-\frac{c^3}3 ,
\]
again the number of \Cref{ex:ch46-toric-p2}.

\emph{The Green function of $\ord_0$.}  Taking $c=1$ gives
$\MA_\omega(\phi_1)=\delta_{\ord_0}$ and
$\phi_1(\ord_0)=-1$, so the normalized Green function of
\Cref{ex:ch47-dirac-green} is $\phi_{\ord_0}=\phi_1+1$.  On the convex
side, \eqref{eq:bj-toric-green} with $u=(1,1)$ gives
\[
 f_{\phi_{\ord_0}}(w)=f_P(w-(1,1))+f_P(1,1)
 =\max\{0,\,1-w_1,\,1-w_2\},
\]
and one checks the four values
$f_{\phi_{\ord_0}}(0,0)=1$,
$f_{\phi_{\ord_0}}(1,1)=0$,
$f_{\phi_{\ord_0}}(2,2)=0$,
$f_{\phi_{\ord_0}}(-1,-1)=2$ against
$\phi_1(v_w)=-\min\{1,\max\{0,\min(w_1,w_2)\}\}$ of
\Cref{ex:ch46-toric-p2}, shifted by $1$ and added to
$f_P(w)=\max\{0,-w_1,-w_2\}$: the two lists agree.  Finally
$E_\omega(\phi_{\ord_0})=E_\omega(\phi_1)+1=-\tfrac13+1=\tfrac23
=E^\vee_\omega(\delta_{\ord_0})$, the table entry, and
$\sup\phi_{\ord_0}=1=T_\omega(\ord_0)$, confirming
\Cref{prop:ch49-toric-green}\,\cref{itm:ch49-tg:sup}.
\end{example}

\subsection{A quasimonomial valuation on a blown-up surface}
\label{sec:ch49-blowup}

\begin{example}[Two blowups, and the linear-growth constant]
\label{ex:ch49-blowup-surface}
\index{quasimonomial valuation!linear growth constant}
We continue \Cref{ex:ch32-blowup-surface}, compactifying $\mathbf A^2$ to
$\mathbf P^2$ and polarizing by $L=\cO(1)$.  Thus
$p=[1:0:0]$ is the origin of the chart $\Spec k[x,y]$;
$\pi_1:Y_1\to\mathbf P^2$ is the blowup of $p$, with exceptional prime
divisor $E_1$; $q\in E_1$ is the point corresponding to the tangent
direction $y=0$; $\pi_2:Y_2\to Y_1$ is the blowup of $q$, with
exceptional prime divisor $E_2$ and $E_1'$ the strict transform of
$E_1$; and in the chart $(z_1,z_2)$ with $E_1'=\{z_1=0\}$,
$E_2=\{z_2=0\}$ one has, by \eqref{eq:ch32-blowup-coordinates},
\[
 x=z_1z_2,\qquad y=z_1z_2^{\,2}.
\]
For $\alpha=(\alpha_1,\alpha_2)\in\mathbf R^2_{\ge0}\setminus\{0\}$ let
$v_\alpha$ be the quasimonomial valuation with
$v_\alpha(z_i)=\alpha_i$; it lies on the dual \emph{cone} complex of the SNC pair
$(Y_2,E_1'+E_2)$, whose two-dimensional cone has these nonnegative
weights.  The slice $\alpha_1+\alpha_2=1$ is the segment with endpoints
$\ord_{E_1}$ and $\ord_{E_2}$.  The calculations below first apply to
the whole cone; statements about the segment use that normalization.
An SNC pair on a birational surface is distinguished here from an
SNC test configuration over $\mathbf A^1$.

\emph{Step 1: $v_\alpha$ is a toric point.}  Put
\[
 u:=(u_1,u_2):=(\alpha_1+\alpha_2,\ \alpha_1+2\alpha_2)
 \in\mathbf R^2_{\ge0},
\]
so that $v_\alpha(x)=u_1$ and $v_\alpha(y)=u_2$ by
\Cref{ex:ch32-blowup-surface}\,\cref{itm:ch32-bs:values}.  We claim
$v_\alpha=v_u$, the toric valuation of
\Cref{lem:ch49-toric-slice}.  Indeed a monomial satisfies
\[
 x^ay^b=z_1^{\,a+b}z_2^{\,a+2b},
\]
and the linear map $(a,b)\mapsto(a+b,a+2b)$ of $\mathbf Z^2$ has
determinant $1$, hence is injective; so distinct monomials of
$k[x,y]$ map to distinct monomials of $k[z_1,z_2]$ with the same
coefficients, and no cancellation occurs.  Since $v_\alpha$ is monomial
in $(z_1,z_2)$,
\[
 v_\alpha\Bigl(\sum c_{ab}x^ay^b\Bigr)
 =\min_{c_{ab}\ne0}\bigl\{(a+b)\alpha_1+(a+2b)\alpha_2\bigr\}
 =\min_{c_{ab}\ne0}\{au_1+bu_2\},
\]
which is $v_u$.  This is
\Cref{ex:ch32-blowup-surface}\,\cref{itm:ch32-bs:consistency} for general
$\alpha$; at $\alpha=(1,1)$ it gives $u=(2,3)$, the case recorded
there.

\emph{Step 2: the linear-growth constant.}  Since $u\ge0$ one has
\[
 f_P(u)=\max\{0,-u_1,-u_2\}=0,
 \qquad
 f_P(-u)=\max\{0,u_1,u_2\}=u_2,
\]
because $u_2-u_1=\alpha_2\ge0$.  By
\eqref{eq:ch49-toric-T},
\[
 \boxed{\ T_\omega(v_\alpha)=u_2=\alpha_1+2\alpha_2\ }
\]
\emph{Independently, and without toric geometry}: a section of
$\cO(m)$ is, in the chart, a polynomial $f$ of degree $\le m$, and
\[
 v_\alpha(f)=\min_{c_{ab}\ne0}\{au_1+bu_2\}
 \le m\max\{u_1,u_2\}=mu_2 ,
\]
with equality for $f=y^m$; so $a_m(v_\alpha)=mu_2$ for every
$m\ge1$ and $T_L(v_\alpha)=u_2$, by \eqref{eq:bj-maximal-vanishing}.
The two computations agree.

\emph{Step 3: the vertices, checked against pseudoeffectivity.}  At
$\alpha=(1,0)$, $v_\alpha=\ord_{E_1}=\ord_p$ and $u=(1,1)$, so
$T_\omega=1$; this is
\Cref{ex:ch49-p2-divisorial}\,\cref{itm:ch49-p2d:point}.  At
$\alpha=(0,1)$, $v_\alpha=\ord_{E_2}$ and $u=(1,2)$, so
$T_\omega=2$.  We verify the latter by
\Cref{prop:ch49-psef-threshold}.  On $Y_2$ write $H$ for the pullback
of the hyperplane class, $F_1:=\pi_2^*E_1=E_1'+E_2$, and $E_2$; then
\[
 H^2=1,\quad F_1^2=E_2^2=-1,\quad
 H\cdot F_1=H\cdot E_2=F_1\cdot E_2=0 ,
\]
This surface is toric, so every effective divisor class is represented
by an effective torus-invariant divisor (\cite[Ch.~6]{CLS11}).
The invariant prime-divisor classes reduce to the three generators
$E_1'=F_1-E_2$, $E_2$, and the strict transform
$\widetilde\ell_0=H-F_1-E_2$ of the line $\{y=0\}$, which passes
through $p$ with tangent direction $y=0$ and hence through $q$.
Solving
\[
 H-cE_2=a\,(H-F_1-E_2)+b\,(F_1-E_2)+e\,E_2,
 \qquad a,b,e\ge0,
\]
gives $a=1$, $b=1$, $e=2-c$; so $H-cE_2$ is pseudoeffective exactly
for $c\le2$, that is $\tau_L(E_2)=2=T_\omega(\ord_{E_2})$.

\emph{Step 4: the other invariants on the cone and segment.}  With
$b_P=(\tfrac13,\tfrac13)$, \eqref{eq:ch49-toric-dirac-energy} gives
\[
 E^\vee_\omega(\delta_{v_\alpha})
 =0+\frac{u_1+u_2}{3}
 =\frac{2\alpha_1+3\alpha_2}{3}
 =\frac{A_{\mathbf P^2}(v_\alpha)}{3},
\]
by \eqref{eq:ch32-blowup-discrepancy}: on this cone the dual energy of the
Dirac mass is one third of the log discrepancy.  The ratio
\[
 \frac{E^\vee_\omega(\delta_{v_\alpha})}{T_\omega(v_\alpha)}
 =\frac{2\alpha_1+3\alpha_2}{3(\alpha_1+2\alpha_2)}
\]
decreases from $\tfrac23$ at $\alpha=(1,0)$ to $\tfrac12$ at
$\alpha=(0,1)$, staying inside $[\tfrac13,\tfrac23]$ as
\eqref{eq:bj-dirac-energy-T} requires.  Finally
\eqref{eq:ch49-toric-dinfty} gives, for all $\alpha,\alpha'$ in the
nonnegative cone and the corresponding $u,u'$ in the nonnegative quadrant,
\[
 d_{\infty,\omega}(v_\alpha,v_{\alpha'})
 =\max\{|u_1-u_1'|,\ |u_2-u_2'|\}
 =\max\bigl\{|\Delta\alpha_1+\Delta\alpha_2|,\
 |\Delta\alpha_1+2\Delta\alpha_2|\bigr\},
\]
$\Delta\alpha:=\alpha-\alpha'$; a norm in $\Delta\alpha$.  So the dual
complex embeds bi-Lipschitzly here, up to and including the endpoints,
which is the situation described in
\Cref{rem:ch49-dual-complex-lipschitz}.

\emph{Step 5: the irrational point.}  Take
$\alpha=(1,\sqrt2)$, the point of
\Cref{ex:ch32-blowup-surface}\,\cref{itm:ch32-bs:irrational}: it is
Abhyankar of rational rank two, not divisorial, and
\[
 u=(1+\sqrt2,\ 1+2\sqrt2),
 \qquad
 T_\omega(v_\alpha)=1+2\sqrt2\approx3.828,
\]
\[
 E^\vee_\omega(\delta_{v_\alpha})=\frac{2+3\sqrt2}{3}\approx2.081 .
\]
The ratio is $\approx0.5436\in[\tfrac13,\tfrac23]$.  For instance
$\alpha'=(1,\tfrac{140}{99})$ has rational coordinates, hence defines a
divisorial valuation by
\Cref{thm:ch32-quasimonomial-invariants}\,\cref{itm:ch32-qi:divisorial}.
The explicit formula of Step~4 gives
\[
 d_{\infty,\omega}(v_\alpha,v_{\alpha'})
 =2\left|\sqrt2-\frac{140}{99}\right|<2\cdot10^{-4}.
\]
This is a concrete local instance of divisorial approximation with an
explicit constant.
\end{example}

\subsection{A valuation without linear growth}
\label{sec:ch49-no-growth}

\begin{example}[A curve valuation along a transcendental branch]
\label{ex:ch49-no-linear-growth}
\index{valuation!without linear growth}
Let $k$ be any algebraically closed field, $X=\mathbf P^2_k$,
$L=\cO(1)$, and keep the chart $\Spec k[x,y]$ with origin $p$.  Put
\[
 \hat y:=\sum_{j\ge1}x^{j!}\in k[[x]],
 \qquad
 P_N:=\sum_{j=1}^{N}x^{j!}\in k[x]\quad(N\ge1),
\]
so that $\deg P_N=N!$ and
$\ord_x(\hat y-P_N)=(N+1)!$.

\emph{Step 1: $\hat y$ is transcendental over $k(x)$.}  Suppose
$F(x,Y)\in k[x,Y]$ is nonzero with $F(x,\hat y)=0$; let
$d:=\deg_YF$ and $e:=\deg_xF$.  If $d=0$ then
$F=F(x)\ne0$ and $F(x,\hat y)=F(x)\ne0$, a contradiction; so
$d\ge1$.  Since $F(x,Y)-F(x,Z)$ is divisible by $Y-Z$ in
$k[x,Y,Z]$, we may write
\[
 F(x,P_N)=F(x,P_N)-F(x,\hat y)=(P_N-\hat y)\,G_N,
 \qquad G_N\in k[[x]],
\]
whence $\ord_xF(x,P_N)\ge(N+1)!$ whenever $F(x,P_N)\ne0$.  But
$F(x,P_N)$ is a polynomial in $x$ of degree at most
$e+d\cdot N!$; so $F(x,P_N)\ne0$ forces
\[
 (N+1)!\le e+d\cdot N!,
 \qquad\text{that is}\qquad
 N+1\le\frac{e}{N!}+d ,
\]
which fails, for example, when $N\ge\max\{d+1,e,2\}$.  Hence $F(x,P_N)=0$ for all
large $N$.  But $F(x,\cdot)$ is a nonzero polynomial of degree $d$
over the field $k(x)$, so it has at most $d$ roots there, while the
$P_N$ are pairwise distinct.  Contradiction.

\emph{Step 2: the valuation.}  Define, for $0\ne f\in k[x,y]$,
\[
 v(f):=\ord_x\bigl(f(x,\hat y)\bigr)\in\mathbf N ,
\]
finite by Step~1, and extend to $k(x,y)^\times=k(X)^\times$ by
$v(f/g)=v(f)-v(g)$.  Then $v$ is a valuation: multiplicativity is
multiplicativity of $\ord_x$ on the domain $k[[x]]$, and the
ultrametric inequality is clear.  It is trivial on $k^\times$, has value
group $\mathbf Z$ and residue field $k$, and
$v(x)=1$, $v(y)=\ord_x\hat y=1$, so $v\ge0$ on $k[x,y]$ with
$v>0$ on $\mathfrak m_p$: the center of $v$ on $\mathbf P^2$ is
$p$, and $v\in X^{\val}$.  Note
$\operatorname{rat.rk}(v)+\operatorname{trdeg}(v)=1+0=1<2=\dim X$, so
$v$ is \emph{not} Abhyankar
(\Cref{thm:abhyankar-inequality}); it is a curve valuation along a formal
branch that is not algebraic.

\emph{Step 3: linear growth fails, with numbers.}  For $N\ge2$ put
\[
 f_N:=y-P_N\in k[x,y],
 \qquad \deg f_N=N!,\qquad v(f_N)=(N+1)! .
\]
Homogenizing, $f_N$ is a nonzero section $F_N\in H^0(X,\cO(N!))$ with
$v(F_N)=(N+1)!$, so
\[
 \frac{a_{N!}(v)}{N!}\ \ge\ \frac{(N+1)!}{N!}=N+1
 \qquad\text{for every }N\ge2 ,
\]
and therefore $T_L(v)=+\infty$ by \eqref{eq:bj-maximal-vanishing}:
$v\notin X^{\lin}$.  The growth of $a_m/m$ is at least
$\sim\log m/\log\log m$ along the subsequence $m=N!$, so it is not
merely unbounded but unbounded with a computable rate.

\emph{Step 4: the dichotomy, exhibited.}  Put
\[
 \begin{aligned}
  \phi_N&:=\frac1{N!}\log|F_N|\in\PSH(\omega),\\
  \sup_{X^{\An}}\phi_N&=0,\\
  \phi_N(v)&=-\frac{v(F_N)}{N!}=-(N+1),
 \end{aligned}
\]
by \eqref{eq:section-log-value}.  Choose
$c_N:=\bigl(N(N+1)\bigr)^{-1}$ for $N\ge2$, so that
$\sum_{N\ge2}c_N=\tfrac12\le1$ while
$\sum_{N\ge2}c_N(N+1)=\sum_{N\ge2}N^{-1}=+\infty$.  The partial sums
\[
 \psi_K:=\sum_{N=2}^{K}c_N\phi_N
\]
are $\omega$-psh by \Cref{thm:bj-basic-closure}\,(3), being convex
combinations of $\phi_2,\dots,\phi_K$ and the constant $0$ with total
weight one; they decrease, each $\phi_N$ being $\le0$; and their limit
$\psi$ satisfies $\psi(v_{\triv})=0$, hence is $\omega$-psh by
\Cref{thm:bj-basic-closure}\,(4).  Finally
\[
 \psi(v)=\sum_{N\ge2}c_N\phi_N(v)=-\sum_{N\ge2}\frac{N+1}{N(N+1)}=-\infty .
\]
So $\{v\}$ is pluripolar, $\Capa^*_\omega(\{v\})=0$,
$\delta_v\notin\cM^1(\omega)$ and
$E^\vee_\omega(\delta_v)=+\infty$; this is
\Cref{cor:ch49-nonpluripolar} in the four equivalent forms, and
\Cref{thm:bj-dirac-finite-energy} in the negative direction.  By contrast,
every $v_\alpha$ of \Cref{ex:ch49-blowup-surface} has
$E^\vee_\omega(\delta_{v_\alpha})<\infty$, and its Green function is even
$\mathbf Q$-PL when $\alpha$ is rational
(\Cref{prop:ch49-toric-green}\,\cref{itm:ch49-tg:ma}).  The two examples
sit at the two sides of the dichotomy, on the same variety and with the
same polarization.

\emph{Step 5: what does and does not distinguish them.}  Both $v$ and
$v_\alpha$ are centered at $p$, both have rank one, and both are limits
of divisorial valuations in the weak topology.  The quasimonomial valuation $v_\alpha$ attains equality in
Abhyankar's inequality, whereas $v$ does not.  In characteristic zero,
an Abhyankar valuation is monomial on a smooth SNC pair
\cite[\S3]{JM12}.  Larger rational weights on its active
cone give a divisorial valuation dominating it up to scale, so
monotonicity of $T_\omega$ proves linear growth.
This uses an SNC pair, not an SNC test configuration of singular $X$.

The example proves that full support alone does not imply linear
growth.  Restriction to an algebraic branch would instead give
proper support and $T_\omega=+\infty$ by
\Cref{prop:bj-linear-growth-potential}; it is not a non-Abhyankar
example in $X^{\lin}$.  No converse classification is asserted.
\end{example}

\subsection{A non-toric divisorial point: the conic}
\label{sec:ch49-conic}

\begin{example}[The order of vanishing along a smooth conic]
\label{ex:ch49-conic}
Let $X=\mathbf P^2$, $L=\cO(1)$, $V=1$, and let $C\subset X$ be a
smooth conic, with equation $q\in H^0(X,\cO(2))$.  Put $v:=\ord_C$, a
divisorial point that is not torus-invariant for the standard torus:
its invariant prime divisors are the coordinate lines.

\emph{The linear-growth constant.}  A nonzero $s\in H^0(X,\cO(m))$ with
$\ord_C(s)=j$ is divisible by $q^j$, so $2j\le m$; and $q^{j}$
realizes the bound when $m=2j$.  Hence $a_m(v)=\lfloor m/2\rfloor$ and
\[
 T_\omega(\ord_C)=T_L(\ord_C)=\tfrac12 ,
\]
in agreement with \Cref{prop:ch49-psef-threshold}, the class
$\cO(1)-cC=\cO(1-2c)$ being pseudoeffective exactly for $c\le\tfrac12$.

\emph{The Green function by intersection theory.}  Let
$\overline\cX:=\operatorname{Bl}_{C\times\{0\}}(\mathbf P^2\times\mathbf P^1)$
with exceptional divisor $E$ and strict transform $\widetilde P$ of
$\mathbf P^2\times\{0\}$, so that $\cX_0=\widetilde P+E$ with both
multiplicities $1$, and $v_E=\ord_C$: this is the deformation to the
normal cone of $C$, and the exceptional divisor of the blowup of a smooth
center $W$ computes $\ord_W$.  Write $H$ for the pullback of the
hyperplane class of $\mathbf P^2$ and $F$ for the fibre class of
$\mathbf P^2\times\mathbf P^1\to\mathbf P^1$, so that
$\pi^*F=\widetilde P+E$.  Since $\widetilde P\cong\mathbf P^2$, with
$H|_{\widetilde P}$ the hyperplane class and
$E|_{\widetilde P}=C$ of degree $2$, we get
\[
 H^3=0,\quad
 H^2\cdot\pi^*F=1,\quad
 H^2\cdot\widetilde P=1,\quad
 H\cdot E\cdot\widetilde P=2,\quad
 E^2\cdot\widetilde P=C^2=4 ,
\]
and $E\cdot\pi^*F=0$, whence
\[
 H^2\cdot E=0,\qquad
 H\cdot E^2=-2,\qquad
 E^3=-4 .
\]
For $c\ge0$ put $\cL:=H-cE$, that is
$\phi_c:=\phi_{-cE}\in\PL_{\mathbf R}(X)$
(and $\phi_c\in\PL(X)$ when $c$ is rational).  Necessity of $0\le c\le\tfrac12$ follows by testing a line
in $\widetilde P$ and a fibre of $E\to C$.
For sufficiency, use the section $q\in H^0(X,2L)$:
\[
 \phi_{1/2}=\frac12\max\{\log|q|,-1\}
\]
is $\omega$-psh by the section construction and maximum stability.
For $0\le c\le\tfrac12$, $\phi_c=2c\phi_{1/2}$ is a convex
combination of $\phi_{1/2}$ and $0$, hence $\omega$-psh.
Thus the exact relative-nef interval is proved without inferring
nefness from two tests alone.  By
\Cref{prop:ch46-pl-measure},
\[
 \MA^{\mathrm u}_\omega(\phi_c)
 =(1-2c)^2\,\delta_{v_{\triv}}+4c(1-c)\,\delta_{\ord_C},
\]
of total mass $(1-2c)^2+4c-4c^2=1=V$, and
\[
 (\omega,\phi_c)^3=(H-cE)^3=-6c^2+4c^3,
 \qquad
 E_\omega(\phi_c)=\frac{(H-cE)^3}{3V}=-2c^2+\frac{4c^3}{3}.
\]

\emph{The numbers at the endpoint.}  At $c=\tfrac12$ the measure
degenerates to $\delta_{\ord_C}$; $\phi_{1/2}(\ord_C)=-\tfrac12$ by
\Cref{prop:ch34-affine-on-faces}, and $\sup\phi_{1/2}=0$.  Hence the
normalized Green function is $\phi_{\ord_C}=\phi_{1/2}+\tfrac12$, with
\[
 S=\sup\phi_{\ord_C}=\tfrac12=T_\omega(\ord_C),
 \qquad
 E_\omega(\phi_{1/2})=-\tfrac12+\tfrac16=-\tfrac13,
\]
\[
 E^\vee_\omega(\delta_{\ord_C})
 =E_\omega(\phi_{\ord_C})=-\tfrac13+\tfrac12=\tfrac16
 =\frac{1}{n+1}\,T_\omega(\ord_C).
\]
The lower constant of \eqref{eq:bj-dirac-energy-T} is therefore attained at
a non-toric divisorial point, with the same ratio $\tfrac13$ as at the
line $\ord_{D_1}$ of \Cref{ex:ch49-p2-dictionary}; and
$S=T_\omega$ here, as in \Cref{rem:ch49-S-vs-T}.  Note also that the
truncation estimate \Cref{lem:ch49-truncation-energy} is sharp here: with
$a=\tfrac12$ and $\psi=\phi_{1/2}$, one has
$3E_\omega(\psi)=-1=-na$.
\end{example}

\subsection{Unit tests}
\label{sec:ch49-unit-tests}

\begin{example}[Five checks of the toric formulae]
\label{ex:ch49-unit-tests}
Each of the following is a computation carried out elsewhere in this book
by a different method; we reproduce it from
\eqref{eq:bj-toric-ma} and \eqref{eq:bj-toric-energy}.

\emph{(i) The trivial metric, on $\mathbf P^n$.}  With $P$ the
standard simplex, $f_0=f_P$ is the support function of
$P_{\mathrm{BJ}}$, so $\partial f_P(0)=P_{\mathrm{BJ}}$ and
$\MA_{\mathbf R}(f_P)=V_P\delta_0$ by
\eqref{eq:appD-polyhedral-atom}; hence
$\MA_\omega(0)=\delta_{v_{\triv}}$, and
$f^\vee_P\equiv0$ on $P_{\mathrm{BJ}}$ gives $E_\omega(0)=0$.  These
are \Cref{lem:ch47-three-normalizations} and
\eqref{eq:bj-trivial-energy-zero}.  In unnormalized form the mass is
$n!\,V_P=(\omega^n)$, which is \Cref{ex:ch08-toric-pn}.

\emph{(ii) $\mathbf P^2$ and \Cref{ex:ch46-toric-p2}.}  Done in
\Cref{ex:ch49-p2-dictionary}: the measure
$c^2\delta_{\ord_0}+(1-c^2)\delta_{v_{\triv}}$ and the energy
$-c^3/3$ agree with the intersection-theoretic computation there, and
with the polarization computation of the mixed measure recorded there.

\emph{(iii) $\mathbf P^1\times\mathbf P^1$ and
\Cref{ex:ch26-toric-surface}, over the two bases.}  Take
$X=\mathbf P^1\times\mathbf P^1$, $\omega=c_1(\cO(1,1))$,
$P=[0,1]^2$, $P_{\mathrm{BJ}}=[-1,0]^2$, $V_P=1$,
$V=2!\cdot1=2$, and
$f_P(u)=\max\{0,-u_1\}+\max\{0,-u_2\}$.  For $s>0$ let
$\phi^{(s)}$ be the invariant psh function with
\[
 f_{\phi^{(s)}}(u)=\max\{-u_1,-s\}+\max\{-u_2,-s\},
\]
which is the function $F_s$ of \Cref{ex:ch26-toric-surface} read in the
present convention.  A one-variable computation gives, for
$g(t)=\max\{-t,-s\}$ and $m\in[-1,0]$,
\[
 g^\vee(m)=\sup_{t\in\mathbf R}\bigl(tm+\min\{t,s\}\bigr)=s(m+1),
\]
so that $f^\vee_{\phi^{(s)}}(m)=s(m_1+1)+s(m_2+1)$ and
\eqref{eq:bj-toric-energy} gives
\[
 E_\omega(\phi^{(s)})
 =-\int_{[-1,0]^2}\bigl(s(m_1+1)+s(m_2+1)\bigr)\rd m
 =-\Bigl(\frac s2+\frac s2\Bigr)=-s .
\]
\Cref{thm:ch26-toric-energy} computed the \emph{unnormalized} energy over a
discretely valued base and obtained
$E(F_s)-E(F_0)=-2s$ in \eqref{eq:ch26-toric-surface-energy}.  The two
agree: $-2s=V\cdot(-s)$, the factor $V=2$ being exactly
\Cref{warn:ch46-normalization}.  A reader who obtains $+s$ has dropped
the minus sign of \eqref{eq:bj-toric-energy}, which comes from the
reflection \eqref{eq:bj-toric-sign-conversion}; a reader who obtains
$-2s$ here has used the unnormalized energy.

\emph{(iv) The Calabi--Yau solutions of \Cref{ex:ch47-toric-cy}.}  There,
$\MA_{\mathbf R}(f)=V_P\mu$ is solved on
$\mathbf P^1\times\mathbf P^1$ for the diffuse datum
$\mu=\nu\otimes\nu$, $\nu=(1+u)^{-2}\rd u$, with
$f(u)=g(u_1)+g(u_2)$ and $g$ as recorded there, and the energy
$E_\omega(\phi)=-1$ is obtained from \eqref{eq:bj-toric-energy}.  The
same number is obtained there a second time from
\eqref{eq:ch46-pure-pairing-as-integrals}, giving
$6E_\omega(\phi)=0+(-2)+(-4)=-6$.  Both computations use the formulae
proved in \Cref{sec:ch49-toric-ma}; neither is circular, the second being
independent of the toric dictionary.

\emph{(v) The Dirac energies of \Cref{ex:ch46-p1-energy}.}  On
$\mathbf P^1$ with $\omega=c_1(\cO(d))$, $P=[0,d]$,
$b_P=d/2$, $V_P=d$.  For $u>0$, \eqref{eq:ch49-toric-T} gives
$T_\omega(v_u)=f_P(u)+f_P(-u)=0+du=du$, and
\eqref{eq:ch49-toric-dirac-energy} gives
$E^\vee_\omega(\delta_{v_u})=0+u\cdot\tfrac d2=\tfrac{du}{2}
=\tfrac12T_\omega(v_u)$, the common value of the two constants
$\tfrac1{n+1}=\tfrac n{n+1}$ for $n=1$.  Setting $u=1$, so that
$v_1=\ord_p$, and comparing with \Cref{ex:ch46-p1-energy} at $c=d$,
where $\MA^{\mathrm u}_\omega(\phi_d)=d\,\delta_{\ord_p}$,
$E_\omega(\phi_d)=-d/2$ and $\phi_d(\ord_p)=-d$, gives
$E^\vee_\omega(\delta_{\ord_p})=-\tfrac d2+d=\tfrac d2$: the same
number.
\end{example}

\begin{remark}[The exceptional-cube sign]
\label{rem:ch49-draft-correction}
The reader comparing \Cref{ex:ch49-conic} with a naive application of the
blowup formulae is warned of one sign that is easy to get wrong.  For a
\emph{curve} $W$ in a smooth threefold $Y$ one has
$E^3=-\deg N_{W/Y}$ and $H\cdot E^2=-\deg(H|_W)$, not $+$; the fastest
check is $Y=\mathbf P^3$, $W$ a line, where
$\widetilde H=H-E$ is the class of the strict transform of a plane
through the line, hence a fibre of the induced map to $\mathbf P^1$, so
that $(H-E)^2=0$ and therefore $E^3=-2=-\deg N_{W/Y}$.  With the wrong
sign one finds $\MA^{\mathrm u}_\omega(\phi_c)=(4c+4c^2)\delta_{\ord_C}$
in \Cref{ex:ch49-conic}, whose mass exceeds $V=1$ already at
$c=\tfrac14$; the total-mass test of \Cref{app:normalizations} detects
this immediately.  We record the trap because the same computation is the
model for every deformation to the normal cone of a positive-dimensional
center.
\end{remark}

\section{Scope and hypotheses}
\label{sec:ch49-scope}

\subsection{Where each hypothesis is used}
\label{sec:ch49-hypothesis-use}

\begin{center}
\begin{tabular}{p{.30\textwidth}p{.62\textwidth}}
\toprule
hypothesis & where it is used\\
\midrule
$k$ algebraically closed &
Throughout, through
\Cref{chap:trivial-analytification}: the description of $X^{\An}$ by
support and center, and the identification of closed points with
$k$-points, which is used in \Cref{ex:ch49-weak-not-strong} and in the
toric charts.\\
$k$ trivially valued &
Everywhere.  It is what makes $X^{\An}$ a cone, $v_{\triv}$ a
distinguished point, $\sup\phi=\phi(v_{\triv})$ a maximum principle, and
$T_\omega$ an exhaustion; see \Cref{sec:ch49-part3}.\\
$X$ integral &
\Cref{lem:ch49-am-finite} (a nonzero section is a regular section) and
the description \eqref{eq:triv-full-support-locus} of $X^{\val}$.  Valuation and section constructions can be made componentwise.
The global probability-normalized energy assertions are stated for
integral $X$: if there is more than one top-dimensional component,
a Dirac mass has infinite dual energy \cite[Corollary~9.13]{BJ22}.\\
$X$ unibranch &
The envelope property supplies the sharp upper Dirac-energy bound,
the Green functions and their identities, and the capacity equivalence;
see \Cref{imp:ch49-envelope,cor:ch49-green-invariants}.
The Green route is also used in \Cref{prop:ch49-embedding-easy}.
The lower bound $E^\vee_\omega(\delta_v)\ge T_\omega(v)/(n+1)$,
the finiteness equivalence \eqref{eq:bj-dirac-e1-iff}, and the basic
height and metric constructions do not need that envelope input.\\
$\omega$ ample &
\Cref{prop:ch49-T-equals-T} (Fubini--Study regularization in a fixed
class), \Cref{prop:ch49-dinfty-basic}\,\cref{itm:ch49-dinf:metric} and
\cref{itm:ch49-dinf:polarization}, and everywhere in
\Cref{sec:ch49-toric,sec:ch49-toric-ma}, where $P$ is the polytope of an
ample toric bundle.  The proofs use fixed ample-class regularization and psh
point separation; they do not establish the same metric conclusions
for an arbitrary big and nef class.\\
$\operatorname{char}k=0$ &
In strong divisorial density, Corollary~A.3; and in the smooth
SNC-test-configuration package, Theorems~A.1 and~A.4.
The pointwise psh topology proof has no characteristic restriction.
Explicit toric formulas are characteristic-free; see
\Cref{warn:ch49-characteristic}.\\
$X$ toric, $L$ toric &
\Cref{sec:ch49-toric,sec:ch49-toric-ma} and
\Cref{ex:ch49-p2-dictionary,ex:ch49-blowup-surface,ex:ch49-unit-tests}.
\Cref{ex:ch49-conic,ex:ch49-no-linear-growth} are deliberately non-toric.\\
\bottomrule
\end{tabular}
\end{center}

\subsection{What is not claimed}
\label{sec:ch49-not-claimed}

\begin{enumerate}
\item\label{itm:ch49-nc:question}
      We do not claim that comparability between $T_\omega$ and
      $E^\vee_\omega(\delta_\cdot)$ promotes the former to a strong
      convergence criterion.  The proposed replacement is false by
      \cref{thm:ch49-global-radon-riesz-counterexample};
      \cref{rem:ch49-open-question-detail} explains why comparability alone
      could never prove it.
\item\label{itm:ch49-nc:S}
      The Green identity $\sup\phi_v=T_\omega(v)$ is proved at the
      ample envelope signature of \Cref{cor:ch49-green-invariants}.
      Its proof uses the solution of the Dirac equation and domination;
      the mere comparison of dual energy and $T_\omega$ would give
      only a factor-$n$ bound.
\item\label{itm:ch49-nc:lipschitz}
      The face-interior estimate $C_\cX/c$ does not itself give a
      boundary constant.  The separate exact source
      \cite[Theorem~A.4]{BJ22} does give equi-Lipschitz continuity
      on closed SNC faces at its smooth characteristic-zero signature.
\item\label{itm:ch49-nc:equicontinuity}
      We do not claim that the family $\PSH^{\sup}(\omega)$ is
      equicontinuous on a $T_\omega$-ball, nor that the weak and strong
      topologies agree there.  \Cref{ex:ch49-weak-not-strong} refutes both.
      What is true is that each \emph{individual} $\omega$-psh function
      is weakly continuous there
      (\Cref{prop:bj-psh-continuous-bounded-valuations}), and that
      statement is imported.
\item\label{itm:ch49-nc:pl}
      We do not claim that the Green function of a divisorial point is
      $\mathbf Q$-PL.  It is for rational toric parameters
      (\Cref{prop:ch49-toric-green}\,\cref{itm:ch49-tg:ma}), and in
      \Cref{ex:ch49-conic}; but \cite[Theorem~60 and Examples~63, 65]{BJGreen24} exhibits a
      divisorial point whose Green function is not, and
      \Cref{warn:ch49-green-not-pl} records the consequence.
\item\label{itm:ch49-nc:extreme}
      Stability under domination is justified by the full
      $\cE^1$ integrability characterization, Theorem~9.20.
      Integrability of only continuous psh test functions would not
      prove it.  \Cref{cor:bj-extremal-measures-dirac} uses the former
      published interface explicitly.
\item\label{itm:ch49-nc:transcendence}
      The example \Cref{ex:ch49-no-linear-growth} produces \emph{one}
      valuation without linear growth, on $\mathbf P^2$ over any
      algebraically closed field.  We do not classify such valuations, and
      the relation between linear growth and the finiteness of other
      asymptotic invariants is not discussed here.  For valuations and
      asymptotic invariants of graded sequences of ideals, see
      \cite{JM12}.
\end{enumerate}

\section{Chapter dictionary}
\label{sec:ch49-dictionary}

\begingroup
\begin{longtable}{p{.24\textwidth}p{.68\textwidth}}
\toprule
object & content and reference\\
\midrule
\endfirsthead
\toprule
object & content and reference\\
\midrule
\endhead
maximal vanishing  & $a_m(v)=\max\{v(s)\}$, superadditive; $T_L=\lim a_m/m$ may be infinite; each $a_m$ is finite and
attained for $v\in X^{\val}$\par\smallskip \eqref{eq:bj-max-vanishing}--\eqref{eq:bj-maximal-vanishing},
\Cref{lem:ch49-am-finite}\\

linear-growth locus  & $X^{\lin}=\{T_L<\infty\}$; a cone with vertex $v_{\triv}$,
$\sigma$-compact, independent of the polarization\par\smallskip \Cref{def:bj-linear-growth},
\Cref{prop:ch49-change-polarization},
\Cref{prop:ch49-T-elementary}\\

divisorial case  & Divisorial linear growth; if $E$ is $\mathbf Q$-Cartier,
$T_L(t\ord_E)=t\,\tau_L(E)$\par\smallskip \Cref{prop:ch49-psef-threshold},
\eqref{eq:ch49-psef-threshold}\\

the two exhaustions  & $T_\omega=T_L$ exactly, with no constant; $T_\omega=+\infty$ off
$X^{\val}$\par\smallskip \eqref{eq:bj-T-valuation}, \Cref{prop:ch49-T-equals-T}\\

truncation estimate  & $c-a\le\psi\le c=\sup\psi\ \Rightarrow\
E_\omega(\psi)\ge c-\frac n{n+1}a$, sharp\par\smallskip \Cref{lem:ch49-truncation-energy}\\

the dichotomy  & $\delta_v\in\cM^1\iff v\in X^{\lin}$, with
$\frac{T_\omega}{n+1}\le E^\vee_\omega(\delta_v)\le
\frac{nT_\omega}{n+1}$\par\smallskip \Cref{thm:bj-dirac-finite-energy},
\eqref{eq:bj-dirac-e1-iff}, \eqref{eq:bj-dirac-energy-T}\\

nonpluripolar locus  & $X^{\lin}=\{v:\{v\}\text{ not pluripolar}\}
=\{v:\Capa^*_\omega(\{v\})>0\}$\par\smallskip \Cref{cor:ch49-nonpluripolar}\\

the metric  & $d_{\infty,\omega}(v,w)=\sup_\phi|\phi(v)-\phi(w)|$, the smallest metric
making psh functions $1$-Lipschitz; $d(v,v_{\triv})=T_\omega(v)$\par\smallskip \Cref{def:bj-dinfty},
\Cref{prop:ch49-dinfty-basic}\\

comparison lemma  & $v\le w\le\lambda v\Rightarrow d_{\infty,\omega}(v,w)\le
(\lambda-1)T_\omega(v)$; Lipschitz for Thompson's part metric\par\smallskip \Cref{lem:ch49-multiplicative-comparison},
\eqref{eq:ch49-multiplicative-comparison}\\

topological theorem  & weak topology of $\PSH$ is pointwise on $X^{\lin}$; completeness;
bi-Lipschitz embedding into $(\cM^1,I^\vee)$; strong criterion\par\smallskip \Cref{thm:bj-linear-topology-complete}\\

easy half of the embedding  & $|E^\vee(\delta_v)-E^\vee(\delta_w)|\le d_{\infty,\omega}$ and
$I^\vee(\delta_v,\delta_w)\le2d_{\infty,\omega}$\par\smallskip \Cref{prop:ch49-embedding-easy},
\eqref{eq:ch49-embedding-easy}\\

dual complexes  & For smooth $X$ in characteristic zero:
$\Delta_\cX$ is strongly compact with $T_\omega$ bounded;
rational points are $d_\infty$-dense; retractions decrease\par\smallskip \Cref{prop:ch49-dual-complex},
\Cref{lem:ch49-retraction-decreases}\\

divisorial density  & $X^{\divis}$ is $d_\infty$-dense in $X^{\lin}$ when
$\operatorname{char}k=0$\par\smallskip \Cref{thm:bj-strong-divisorial-density},
\Cref{warn:ch49-characteristic}\\

toric slice  & $N_{\mathbf R}\hookrightarrow X^{\val}$, retraction $\rho$ with
$\iota\rho(v)\le v$; $v_u(s_P)=f_P(u)$\par\smallskip \Cref{lem:ch49-toric-slice},
\eqref{eq:ch49-toric-section}\\

sign conversion  & $f_{P_{\mathrm{BJ}}}=\check h_P$, $P_{\mathrm{BJ}}=-P$\par\smallskip \Cref{warn:ch49-toric-signs},
\eqref{eq:bj-toric-sign-conversion}\\

toric $T_\omega$  & $T_\omega(v_u)=f_P(u)+f_P(-u)$, the width of $P$ in the direction
$u$; exact for every $m$\par\smallskip \Cref{prop:ch49-toric-T},
\eqref{eq:ch49-toric-T}\\

toric $d_\infty$  & $\max\{f_P(u'-u)+f_P(u)-f_P(u'),\ \cdots\}$\par\smallskip \Cref{prop:ch49-toric-dinfty},
\eqref{eq:ch49-toric-dinfty}\\

toric dictionary  & invariant $\omega$-psh $\leftrightarrow$ convex $f\le f_P+O(1)$
$\leftrightarrow$ lsc convex $g$ on $P_{\mathrm{BJ}}$\par\smallskip \Cref{thm:bj-toric-psh-convex},
\eqref{eq:bj-toric-convex-function}--\eqref{eq:bj-toric-legendre}\\

toric Monge--Amp\`ere  & For invariant bounded or finite-energy $\phi$,
$\MA_\omega(\phi)=V_P^{-1}\iota_*\MA_{\mathbf R}(f_\phi)$;
unnormalized $=n!\,\iota_*\MA_{\mathbf R}(f_\phi)$\par\smallskip \eqref{eq:bj-toric-ma},
\Cref{prop:ch49-toric-ma-pl},
\Cref{prop:ch49-toric-ma-bounded}\\

toric energy  & $E_\omega(\phi)=-V_P^{-1}\int_{P_{\mathrm{BJ}}}f^\vee_\phi\,\rd m$;
finite energy $\iff f^\vee_\phi\in L^1$\par\smallskip \eqref{eq:bj-toric-energy}, \eqref{eq:bj-toric-e1},
\Cref{prop:ch49-toric-energy-bounded},
\Cref{prop:ch49-toric-energy-singular}\\

toric Green functions  & $f_{\phi_u}(w)=f_P(w-u)+f_P(u)$;
$E^\vee_\omega(\delta_{v_u})=f_P(u)+\langle u,b_P\rangle$;
$\sup\phi_u=T_\omega(v_u)$\par\smallskip \eqref{eq:bj-toric-green},
\Cref{prop:ch49-toric-green},
\eqref{eq:ch49-toric-dirac-energy}\\

convex-geometric shadow  & \eqref{eq:bj-dirac-energy-T} on the toric slice is Minkowski's centroid
inequality\par\smallskip \Cref{cor:ch49-minkowski}\\
\bottomrule
\end{longtable}
\endgroup

\section{Source and provenance}
\label{sec:ch49-source}

\subsection{Source note}
\label{sec:ch49-source-note}

Linear growth and maximal vanishing are \cite[\S\S1.5 and~4.5]{BJ22}.  The
Dirac-energy criterion, $d_\infty$, completeness, sublinear growth, and
strong divisorial density are \cite[\S\S11.1--11.5]{BJ22}.  The
characteristic-zero dual-complex input is \cite[Appendix~A]{BJ22}.  The
toric dictionary and formulae are \cite[Appendix~B]{BJ22}.  Its normalization
in the printed text suppresses
$V_P^{-1}$ in the displayed energy formula by integrating against
normalized Lebesgue probability measure on $P$; we wrote the factor
explicitly because $\rd m$ above is lattice Lebesgue measure.

Statement by statement, the attributions are as follows.
\Cref{def:bj-linear-growth} and the elementary theory of $a_m$ and
$T_L$ are \cite[\S1.5]{BJ22}, where $T_L$ is written $T$;
\Cref{prop:bj-linear-growth-potential} is
\cite[Definition~1.26 and Proposition~4.38]{BJ22}.
\Cref{thm:bj-dirac-finite-energy} is \cite[Proposition~11.1]{BJ22} and
\Cref{cor:bj-extremal-measures-dirac} is
\cite[Proposition~11.3]{BJ22}.  On the topological side,
\Cref{thm:bj-linear-topology-complete} collects
\cite[Theorems~11.4, 11.9 and~11.13 and Corollary~11.15]{BJ22}, the
section formula \eqref{eq:bj-dinfty-sections} is
\cite[Lemma~11.7]{BJ22}, the bounded-set continuity and growth statements of
\Cref{prop:bj-psh-continuous-bounded-valuations} are
\cite[Theorems~11.11 and~11.12]{BJ22}, and
\Cref{thm:bj-strong-divisorial-density} is
\cite[Theorem~11.18]{BJ22}, resting on \cite[Appendix~A]{BJ22}.  The
question recorded in \cref{rem:ch49-open-question} is
\cite[Question~11.16]{BJ22}; the counterexample in
\cref{thm:ch49-global-radon-riesz-counterexample} is the verified
book-level result proved above.  The toric statements
\Cref{thm:bj-toric-psh-convex} and \Cref{thm:bj-toric-ma-energy}, together
with \eqref{eq:bj-toric-green}, are \cite[Appendix~B]{BJ22}.

The non-BJ inputs are: the convex analysis of \Cref{app:convexity}, in
particular \Cref{thm:appD-alexandrov-measure} and
\Cref{prop:appD-legendre-pushforward}, which follow \cite{Roc70} and
\cite{Sch14}; the toric intersection formula
\cite[Theorem~13.4.3]{CLS11} and the toric section--polytope dictionary
\cite[Ch.~6]{CLS11}, used in \Cref{lem:ch49-toric-slice} and
\Cref{prop:ch49-toric-ma-pl}; the quasimonomial description of Abhyankar
valuations \cite[\S3]{JM12}, quoted once in
\Cref{ex:ch49-no-linear-growth} and not used later; and the
non-$\mathbf Q$-PL Green function of \cite[Theorem~60 and Examples~63, 65]{BJGreen24}, quoted
in \Cref{warn:ch49-green-not-pl}.  Over a nontrivially valued base the
toric comparison is \Cref{thm:toric-ma-comparison}, following \cite{CL06}
and \cite{BGPS14}, with the tropical refinements of \cite{BGJKK25} and
\cite{Vil20}.

\subsection{Proof and example provenance}
\label{sec:ch49-additions}

The published Section~11 supplies the valuation-theoretic conclusions.
The chapter expands their arguments and records explicit constants,
without claiming novelty or priority.  The finite flag proof
\Cref{lem:ch49-am-finite} establishes the finite-degree maximum.
The polarization estimates give the explicit constants $p,q$.
The divisorial threshold proof supplies the $\mathbf Q$-Cartier
hypothesis and conductor descent needed on a nonnormal base.
The exact identity $T_\omega=T_L$ is already
\cite[Lemma~4.46]{BJ22}, not merely a comparison up to normalization.

The truncation estimate and Green-function proof of the Dirac bounds
make the dependence on the envelope property visible.
The pointwise-topology proof follows the fixed-measure truncation
argument of Theorem~11.4 in every characteristic.
Completeness is derived here from compact $T$-sublevels and bounded-set
continuity.  The upper embedding constant $2$ is proved from the
Fenchel defect; the lower comparison is Theorem~11.13.
SNC boundary equicontinuity is Theorem~A.4, not a missing conclusion
of the elementary interior estimate.

The toric section calculation proves convexity and the closed
$T$ and $d_\infty$ formulas directly.  General invariant extension
is the published Appendix~B dictionary.
Rational intersection measures pass to bounded invariant potentials
through monomial regularization and Alexandrov continuity.
The energy formula is derived by integrating the first variation,
with almost-everywhere conjugate differentiability and Fubini's theorem.
The finite-energy measure consequence uses bounded truncations and
the absence of lost Alexandrov mass.  The full strong/L1 correspondence
is quoted at its exact Appendix~B scope.
The Green-function and Dirac-energy formulas also appear in that
appendix; the barycenter comparison explains their constants.

The worked examples provide explicit section, blowup, polytope,
capacity, and energy calculations.  In particular, the factorial-gap
formal branch is shown transcendental in every characteristic.
The Radon--Riesz counterexample
\Cref{thm:ch49-global-radon-riesz-counterexample} is the book-level
result recorded in the existing verification receipts; the published
source supplies Question~11.16 and its definitions.
The present direct editorial review does not extend a historical
independent-verifier verdict to newly edited chapter bytes.
The conic and toric numerical examples are checked against
\Cref{chap:bj-energy,chap:bj-finite-energy};
\Cref{rem:ch49-draft-correction} explains the exceptional-cube sign.

\subsection{Historical note}
\label{sec:ch49-history}

The idea that the analytically relevant valuations are those whose
vanishing on sections is linearly bounded goes back to the study of
asymptotic invariants of graded linear series.  Jonsson and Musta\c{t}\u{a}
\cite{JM12} organized the subject around the functions
$v\mapsto v(\mathfrak a_\bullet)$ attached to graded sequences of
ideals, and isolated the classes of valuations --- Abhyankar,
quasimonomial, divisorial --- on which those functions are computable; the
inequality $\ord_E(s)\le C\deg s$ of
\Cref{prop:ch32-divisorial-linear-growth} is the oldest ingredient and is
Zariski's.  In the nontrivially valued setting, Boucksom, Favre and
Jonsson \cite[Theorem~6.1]{BFJ16} establish the uniform facewise
Lipschitz control on SNC skeleta used by \cite[Theorem~A.4]{BJ22}.
What is specific to the trivially valued theory is
that $T_\omega$ becomes an \emph{exhaustion} of a cone, with the
proper-support points at infinity; this is
\cite[\S1.5]{BJ22}, and it is the structural reason the chapter exists.

The identification of the finite-energy locus with the linear-growth locus,
and the two constants of \eqref{eq:bj-dirac-energy-T}, are
\cite[\S11]{BJ22}.  For a positive-dimensional compact Kähler manifold, every point
is pluripolar and its Dirac mass has infinite dual energy, since
finite-energy measures do not charge pluripolar sets \cite{GZ07}.
The non-Archimedean result has a different pointwise finiteness locus:
function-field valuations can be nonpluripolar.
Alexander--Taylor capacity remains a useful comparison with the
complex theory; see \Cref{sec:ch48-alexander-taylor}.
Here the metric and measure-gauge comparison is the precise
Theorem~11.13 of \cite{BJ22}, rather than an identification with a
complex singularity metric.

The toric dictionary is much older in spirit.  That a torus-invariant
metric on a polarized toric variety is a convex function on
$N_{\mathbf R}$, and that its Monge--Amp\`ere measure is the real
Monge--Amp\`ere measure of that function, is the non-Archimedean form of a
statement going back to the analytic theory of toric metrics; the
arithmetic version, with heights and Legendre transforms, is the theory of
Burgos Gil, Philippon and Sombra \cite{BGPS14}, and the model-metric
comparison over a nontrivially valued field is
\Cref{thm:toric-ma-comparison}, following Chambert-Loir \cite{CL06}.  The
trivially valued version, with the reflected polytope $P_{\mathrm{BJ}}$
and the Legendre transform on it, is \cite[Appendix~B]{BJ22}.  The
tropical-toric pluripotential theory of Burgos Gil, Gubler, Jell and
K\"unnemann \cite{BGJKK25} develops the same dictionary intrinsically on
tropical toric varieties, and Vilsmeier's local comparison \cite{Vil20}
supplies the polystable analogue.

Two threads leave this chapter.  The first is \Cref{chap:bj-stability},
where $X^{\lin}$, $T_\omega$ and the toric formulae become the
computational engine of the valuative criteria: a filtration is tested on
divisorial valuations, and its energy is read off a polytope.  The second
is \Cref{part:comparisons}, where the toric formulae are the first place at
which the trivially valued, discretely valued and complex theories are
compared on the same object and produce the same number; the master
dictionary \Cref{chap:master-dictionary} records the comparison, and
\Cref{ex:ch49-unit-tests} is its toric row.
